% Marketing et publicité — Allemand Tle A — Cours signé OnBuch+
% Programme officiel MINESEC. Compiler avec : tectonic AL07-marketing-et-publicite.tex
\newcommand{\DOCMATIERE}{Allemand}
\newcommand{\DOCNIVEAU}{Tle A}
\newcommand{\DOCMODULE}{Chapitre 2 : Vie économique}
\newcommand{\DOCLECON}{AL07}
\newcommand{\DOCTITRE}{Marketing et publicité}
% OnBuch+ — préambule « cours signés » ENRICHI (compilé avec tectonic / XeLaTeX)
% Système visuel « Pop » d'OnBuch+ : fond crème (jamais blanc), bordures encre
% épaisses, ombres « dures » décalées, titres Archivo Black en capitales.
% Version MATHÉMATIQUES : graphes (pgfplots/tikz), boîtes pédagogiques
% (définition, propriété, méthode, attention, le savais-tu, exercice résolu,
% exercice, TP, bilan), page de garde et signature OnBuch+ ancrée en bas.
%
% Macros attendues dans main.tex AVANT \input{preamble} :
%   \DOCMATIERE  \DOCNIVEAU  \DOCMODULE  \DOCLECON (numéro)  \DOCTITRE
\documentclass[11pt]{article}

\usepackage{fontspec}
\usepackage[ngerman,french]{babel} % français = langue principale ; l'allemand via \all{..} / textebox
\usepackage[a4paper,margin=2cm,top=3.4cm,bottom=2.8cm,headheight=1.4cm,footskip=1.3cm]{geometry}
\usepackage{amsmath,amssymb,amsfonts}
\usepackage{mathtools} % \xmapsto, \coloneqq... souvent utilisés par les modèles
\usepackage{cancel} % simplifications barrées (bilans de forces)
% \abs{z} (module, valeur absolue) et \norm{u} (norme), utilisés par les modèles
\providecommand{\abs}{}\let\abs\relax\DeclarePairedDelimiter{\abs}{\lvert}{\rvert}
\providecommand{\norm}{}\let\norm\relax\DeclarePairedDelimiter{\norm}{\lVert}{\rVert}
\usepackage[table]{xcolor} % [table] : fournit aussi \rowcolors (alternance de lignes)
\usepackage{pagecolor}
\usepackage{tikz}
\usetikzlibrary{arrows.meta,positioning,calc,shapes.geometric,shapes.misc,shapes.arrows,decorations.pathmorphing,decorations.pathreplacing,decorations.markings,patterns,shadows,mindmap,trees,fit,backgrounds,matrix,intersections,angles,quotes}
\usepackage{pgfplots}
\usepackage[version=4]{mhchem} % équations nucléaires \ce{^{235}_{92}U}
\usepackage[european,siunitx]{circuitikz} % schémas électriques
\pgfplotsset{compat=1.18}
\usepgfplotslibrary{fillbetween,groupplots} % aires entre courbes (calcul intégral)
\usepackage{enumitem}
\usepackage{fancyhdr}
% [most] charge toutes les bibliothèques tcolorbox (listings, minted, raster,
% documentation...) et gonfle inutilement la mémoire TeX sur un document long
% (~300 boîtes) : on ne charge que ce qui est réellement utilisé.
\usepackage[skins,breakable]{tcolorbox}
\usepackage{graphicx}
\usepackage{colortbl}
\usepackage{booktabs}
\usepackage{tabularx}
\usepackage{diagbox} % case d'en-tête diagonale des tableaux à double entrée
% Colonnes extensibles centrée / gauche / droite (C, L, R), souvent utilisées par les modèles
\newcolumntype{C}{>{\centering\arraybackslash}X}
\newcolumntype{L}{>{\raggedright\arraybackslash}X}
\newcolumntype{R}{>{\raggedleft\arraybackslash}X}
\usepackage{array}
\usepackage{multicol}
\usepackage{multirow}
\usepackage{siunitx}
% parse-numbers=false : accepte aussi \SI{2,5 \times 10^{-3}}{...}
\sisetup{locale=FR,output-decimal-marker={,},group-minimum-digits=4,parse-numbers=false}
\DeclareSIUnit\ohm{\ensuremath{\symup{\Omega}}} % U+2126 absent de la police
\DeclareSIUnit\division{div} % réglages d'oscilloscope (ms/div, V/div)
\DeclareSIUnit\tr{tr} % tours (vitesse de rotation, ex. \SI{1500}{\tr\per\minute})
\DeclareSIUnit\molar{M} % concentration molaire (ex. \SI{3}{\molar}, \SI{1}{\micro\molar})
\DeclareSIUnit\an{an} % année (le modèle confond parfois avec \year, un compteur TeX) — ex. \SI{5}{\kilo\meter\per\an}
\DeclareSIUnit\FCFA{FCFA} % franc CFA (ex. \SI{500}{\FCFA\per\kilo\gram})
\DeclareSIUnit\degreeBrix{\textdegree Brix} % teneur en sucre d'un sirop/jus (réfractométrie)
\DeclareSIUnit\month{mois} % le modèle utilise parfois \month, un compteur TeX, comme unité de durée
\usepackage{qrcode}
% \degree : commande fréquemment utilisée par les modèles pour un angle ou une
% température en dehors de \SI{}{} (non fournie par les paquets chargés).
\providecommand{\degree}{^{\circ}}
% \qed (fin de démonstration), utilisé par les modèles sans amsthm
\providecommand{\qed}{\hfill$\square$}
% \barre : double barre des valeurs interdites dans un tableau de variations (tkz-tab)
\providecommand{\barre}{\ensuremath{\|}}
% environnement proof (amsthm non chargé) utilisé par les modèles
\ifdefined\proof\else\newenvironment{proof}[1][Démonstration]{\par\noindent\textit{#1.}\ }{\qed\par}\fi
\usepackage{lastpage}
\usepackage{hyperref}
\hypersetup{hidelinks}

% ============================================================
% Palette « Pop » OnBuch+ — mêmes hex que la classe P (plus_ui.dart)
% ============================================================
\definecolor{poporange}{HTML}{F59321}
\definecolor{popdark}{HTML}{E07A0C}
\definecolor{popcream}{HTML}{FFF6E8}
\definecolor{popink}{HTML}{1C1714}
\definecolor{poppurple}{HTML}{7A5AE0}
\definecolor{popgreen}{HTML}{2E9E6A}
\definecolor{poppink}{HTML}{E0457B}
\definecolor{popblue}{HTML}{2F7DE1}
\definecolor{popgold}{HTML}{C98A12}
\definecolor{popmuted}{HTML}{8A7F76}
\definecolor{popline}{HTML}{EADFD2}
\definecolor{popgray}{HTML}{8A7F76}
\colorlet{popgrey}{popgray}
% Alias tolérants (le modèle nomme parfois les couleurs différemment)
\colorlet{popwhite}{white}
\colorlet{popblack}{popink}
\colorlet{popred}{poppink}
\colorlet{popyellow}{popgold}
% Teintes claires (fonds de tableaux/encadrés) — le modèle nomme parfois
% les couleurs avec un suffixe L (ex. popblueL) sans les avoir définies.
\colorlet{poporangeL}{poporange!20}
\colorlet{popdarkL}{popdark!20}
\colorlet{poppurpleL}{poppurple!20}
\colorlet{popgreenL}{popgreen!20}
\colorlet{poppinkL}{poppink!20}
\colorlet{popblueL}{popblue!20}
\colorlet{popgoldL}{popgold!20}
\colorlet{popredL}{popred!20}
\colorlet{popgrayL}{popgray!20}
% Teintes claires pour fonds de tableaux / zones
\colorlet{popblueL}{popblue!10!white}
\colorlet{poporangeL}{poporange!14!white}
\colorlet{popgreenL}{popgreen!12!white}
\colorlet{poppurpleL}{poppurple!12!white}
\colorlet{poppinkL}{poppink!12!white}
\colorlet{popgoldL}{popgold!14!white}

\pagecolor{popcream}

% ============================================================
% Typographie
% ============================================================
\defaultfontfeatures{Ligatures=TeX}
\setmainfont{PlusJakartaSans-Medium.ttf}[Path=../../../fonts/,BoldFont=PlusJakartaSans-Bold.ttf]
% Maths en sans-serif (Fira Math) avec chiffres et lettres droites de Plus
% Jakarta, pour que les formules se fondent dans le texte.
\usepackage[math-style=ISO,bold-style=ISO]{unicode-math}
\setmathfont{FiraMath-Regular.otf}[Path=../../../fonts/]
\setmathfont{PlusJakartaSans-Medium.ttf}[Path=../../../fonts/,range={up/num,up/{Latin,latin},\mathup}]
\usepackage{icomma} % virgule décimale sans espace en mode math
% Caractères mathématiques Unicode absents des polices de texte (Plus Jakarta,
% Archivo Black) : rendus par leur équivalent mathématique, en texte comme en
% formule — sinon ils s'affichent en carré vide (ex. « Arithmétique dans ℤ »).
\usepackage{newunicodechar}
\newunicodechar{ℝ}{\ensuremath{\mathbb{R}}}
\newunicodechar{ℤ}{\ensuremath{\mathbb{Z}}}
\newunicodechar{ℂ}{\ensuremath{\mathbb{C}}}
\newunicodechar{ℕ}{\ensuremath{\mathbb{N}}}
\newunicodechar{ℚ}{\ensuremath{\mathbb{Q}}}
\newunicodechar{𝔻}{\ensuremath{\mathbb{D}}}
\newunicodechar{α}{\ensuremath{\alpha}}
\newunicodechar{β}{\ensuremath{\beta}}
\newunicodechar{γ}{\ensuremath{\gamma}}
\newunicodechar{δ}{\ensuremath{\delta}}
\newunicodechar{ε}{\ensuremath{\varepsilon}}
\newunicodechar{θ}{\ensuremath{\theta}}
\newunicodechar{λ}{\ensuremath{\lambda}}
\newunicodechar{μ}{\ensuremath{\mu}}
\newunicodechar{σ}{\ensuremath{\sigma}}
\newunicodechar{τ}{\ensuremath{\tau}}
\newunicodechar{φ}{\ensuremath{\varphi}}
\newunicodechar{ψ}{\ensuremath{\psi}}
\newunicodechar{ω}{\ensuremath{\omega}}
\newunicodechar{Σ}{\ensuremath{\Sigma}}
% majuscules produites par \MakeUppercase dans les titres (α-aminés -> Α)
\newunicodechar{Α}{\ensuremath{\alpha}}
\newunicodechar{Β}{\ensuremath{\beta}}
\newunicodechar{Γ}{\ensuremath{\Gamma}}
\newunicodechar{Δ}{\ensuremath{\Delta}}
\newunicodechar{Ε}{\ensuremath{\varepsilon}}
\newunicodechar{Θ}{\ensuremath{\Theta}}
\newunicodechar{Λ}{\ensuremath{\Lambda}}
\newunicodechar{Μ}{\ensuremath{\mu}}
\newunicodechar{Τ}{\ensuremath{\tau}}
\newunicodechar{Φ}{\ensuremath{\Phi}}
\newunicodechar{Ψ}{\ensuremath{\Psi}}
\newunicodechar{Ω}{\ensuremath{\Omega}}
\newunicodechar{↦}{\ensuremath{\mapsto}}
\newunicodechar{∈}{\ensuremath{\in}}
\newunicodechar{∉}{\ensuremath{\notin}}
\newunicodechar{∧}{\ensuremath{\wedge}}
\newunicodechar{⇔}{\ensuremath{\Leftrightarrow}}
\newunicodechar{⟺}{\ensuremath{\Longleftrightarrow}}
\newunicodechar{⇒}{\ensuremath{\Rightarrow}}
\newunicodechar{⟹}{\ensuremath{\Longrightarrow}}
\newunicodechar{∘}{\ensuremath{\circ}}
\newunicodechar{∪}{\ensuremath{\cup}}
\newunicodechar{∩}{\ensuremath{\cap}}
\newunicodechar{≡}{\ensuremath{\equiv}}
\newunicodechar{⊂}{\ensuremath{\subset}}
\newunicodechar{⊥}{\ensuremath{\perp}}
\newunicodechar{∃}{\ensuremath{\exists}}
\newunicodechar{∀}{\ensuremath{\forall}}
\newunicodechar{∛}{\ensuremath{\sqrt[3]{\,}}}
\newunicodechar{∜}{\ensuremath{\sqrt[4]{\,}}}
\newunicodechar{⃗}{\ensuremath{{}^{\to}}}
\newunicodechar{ⁿ}{\ifmmode^{n}\else\textsuperscript{\ensuremath{n}}\fi}
\newunicodechar{ˣ}{\ifmmode^{x}\else\textsuperscript{\ensuremath{x}}\fi}
\newunicodechar{ᵇ}{\ifmmode^{b}\else\textsuperscript{\ensuremath{b}}\fi}
\newunicodechar{ʳ}{\ifmmode^{r}\else\textsuperscript{\ensuremath{r}}\fi}
\newunicodechar{ᵘ}{\ifmmode^{u}\else\textsuperscript{\ensuremath{u}}\fi}
\newunicodechar{ʸ}{\ifmmode^{y}\else\textsuperscript{\ensuremath{y}}\fi}
\newunicodechar{ᵃ}{\ifmmode^{a}\else\textsuperscript{\ensuremath{a}}\fi}
\newunicodechar{ᵉ}{\ifmmode^{e}\else\textsuperscript{\ensuremath{e}}\fi}
\newunicodechar{ᵏ}{\ifmmode^{k}\else\textsuperscript{\ensuremath{k}}\fi}
\newunicodechar{ᵖ}{\ifmmode^{p}\else\textsuperscript{\ensuremath{p}}\fi}
\newunicodechar{ᴺ}{\ifmmode^{N}\else\textsuperscript{\ensuremath{N}}\fi}
\newunicodechar{ᶜ}{\ifmmode^{c}\else\textsuperscript{\ensuremath{c}}\fi}
\newunicodechar{ʰ}{\ifmmode^{h}\else\textsuperscript{\ensuremath{h}}\fi}
\newunicodechar{ᵠ}{\ifmmode^{q}\else\textsuperscript{\ensuremath{q}}\fi}
\newunicodechar{ₙ}{\ifmmode_{n}\else\textsubscript{\ensuremath{n}}\fi}
\newunicodechar{ₐ}{\ifmmode_{a}\else\textsubscript{\ensuremath{a}}\fi}
\newunicodechar{ᵢ}{\ifmmode_{i}\else\textsubscript{\ensuremath{i}}\fi}
\newunicodechar{ᵣ}{\ifmmode_{r}\else\textsubscript{\ensuremath{r}}\fi}
\newunicodechar{ₖ}{\ifmmode_{k}\else\textsubscript{\ensuremath{k}}\fi}
\newunicodechar{ᵦ}{\ifmmode_{\beta}\else\textsubscript{\ensuremath{\beta}}\fi}
\newunicodechar{ₓ}{\ifmmode_{x}\else\textsubscript{\ensuremath{x}}\fi}
\newfontfamily\popdisplay{ArchivoBlack-Regular.ttf}[Path=../../../fonts/]
\newfontfamily\poplogo{SpaceGrotesk-Bold.ttf}[Path=../../../fonts/]
\newfontfamily\popsemi{PlusJakartaSans-SemiBold.ttf}[Path=../../../fonts/]
\newfontfamily\popxbold{PlusJakartaSans-ExtraBold.ttf}[Path=../../../fonts/]
\newfontfamily\popxboldls{PlusJakartaSans-ExtraBold.ttf}[Path=../../../fonts/,LetterSpace=3.0]

\setlist[enumerate]{leftmargin=1.7em,itemsep=5pt,topsep=4pt}
\setlist[itemize]{leftmargin=1.7em,itemsep=4pt,topsep=4pt,label={\color{poporange}\textbullet}}
\setlength{\parindent}{0pt}
\setlength{\parskip}{5pt}
\renewcommand{\arraystretch}{1.25}

% pgfplots : style Pop par défaut
\pgfplotsset{
  every axis/.append style={axis line style={popink,line width=1pt},tick style={popink},
    grid style={popline,line width=0.5pt},label style={font=\small},tick label style={font=\footnotesize}},
  popcurve/.style={poporange,line width=1.8pt,no markers,smooth},
}
% Même style disponible dans un \draw TikZ simple (hors pgfplots)
\tikzset{popcurve/.style={poporange,line width=1.8pt}}
\tikzset{popgrid/.style={popline,very thin}}
\colorlet{popcurve}{poporange} % le nom du style est parfois utilisé comme couleur

% ============================================================
% Pill / squiggle
% ============================================================
\newcommand{\poppill}[3]{%
  \tikz[baseline=-0.55ex]\node[draw=popink,line width=1.2pt,rounded corners=2.6pt,%
    fill=#2,inner xsep=6.5pt,inner ysep=3pt]{%
    \popxboldls\fontsize{7}{7}\selectfont\color{#3}\MakeUppercase{#1}};%
}
\newcommand{\popsquiggle}[1]{%
  \tikz[baseline=-0.4ex]{
    \draw[#1,line width=2.6pt,line cap=round]
      plot [smooth] coordinates {(0,0.09) (0.34,0.01) (0.68,0.21) (1.02,0.01) (1.36,0.10)};
  }%
}
% Mot-clé surligné (marqueur orange sous le texte)
\newcommand{\cle}[1]{{\popxbold\color{popdark}#1}}

% ============================================================
% En-tête / pied
% ============================================================
\pagestyle{fancy}
\fancyhf{}
\renewcommand{\headrulewidth}{0pt}
\renewcommand{\footrulewidth}{0pt}
\lhead{\raisebox{-0.15ex}{\poplogo\fontsize{13}{13}\selectfont\color{popink}OnBuch\color{poporange}+}}
\rhead{\poppill{\DOCMATIERE\ \textbullet\ \DOCNIVEAU\ \textbullet\ Leçon \DOCLECON}{white}{popink}}
\lfoot{\popxbold\fontsize{7.6}{7.6}\selectfont\color{popmuted}\MakeUppercase{Cours signé OnBuch+}\ \textbullet\ {\popsemi\DOCTITRE}}
\rfoot{\tikz[baseline=-0.6ex]\node[rounded corners=8.5pt,fill=popink,minimum height=17pt,minimum width=17pt,inner xsep=5pt,inner sep=0]{\popxbold\fontsize{8}{8}\selectfont\color{popcream}\thepage\,/\,\pageref*{LastPage}};}

% ============================================================
% Titres
% ============================================================
\newcommand{\chaptitre}[2]{%
  \begin{tcolorbox}[enhanced,colback=poporange,colframe=popink,boxrule=2.6pt,arc=11pt,
      drop shadow=popink,left=15pt,right=15pt,top=12pt,bottom=11pt,before skip=3pt,after skip=18pt]
    \poppill{#2}{popink}{popcream}\par\vspace{9pt}
    {\raggedright\hyphenpenalty=10000\exhyphenpenalty=10000\popdisplay\fontsize{22}{24}\selectfont\color{popcream}\MakeUppercase{#1}\par}
  \end{tcolorbox}
}
% Section numérotée : pastille orange avec le numéro + titre Archivo
\newcounter{popsec}
\newcommand{\coursec}[1]{%
  \stepcounter{popsec}\setcounter{popsub}{0}%
  \par\vspace{16pt}\noindent
  \begin{minipage}{\linewidth}
  \tikz[baseline=-0.7ex]\node[draw=popink,line width=1.6pt,fill=poporange,rounded corners=4pt,minimum size=22pt,inner sep=0]{\popdisplay\fontsize{12}{12}\selectfont\color{popcream}\thepopsec};%
  \hspace{8pt}{\popdisplay\fontsize{15}{17}\selectfont\color{popink}\MakeUppercase{#1}}\par
  \vspace{4pt}\noindent{\color{poporange}\rule{36pt}{3.6pt}}
  \end{minipage}\par\nopagebreak\vspace{8pt}
}
\newcounter{popsub}
\newcommand{\courssub}[1]{%
  \stepcounter{popsub}%
  \par\vspace{9pt}\noindent{\popxbold\fontsize{11.5}{13}\selectfont\color{popink}{\color{poporange}\thepopsec.\thepopsub}\hspace{6pt}#1}\par\nopagebreak\vspace{4pt}
}

% ============================================================
% Boîtes pédagogiques — même recette : fond blanc, bordure épaisse colorée,
% ombre dure encre, badge plein. Titre optionnel [#1].
% ============================================================
\tcbset{popbase/.style={breakable,enhanced,colback=white,boxrule=2.3pt,arc=9pt,
  shadow={3pt}{-3pt}{0pt}{popink},left=14pt,right=14pt,top=11pt,bottom=11pt,before skip=11pt,after skip=15pt,
  fontupper=\popsemi\fontsize{10.6}{14.5}\selectfont\color{popink},
  fontlower=\popsemi\fontsize{10.6}{14.5}\selectfont\color{popink}}}
% Coche dessinée (le glyphe ✓ n'existe pas dans les polices du document)
\newcommand{\popcheck}[1]{\tikz[baseline=-0.5ex]\draw[#1,line width=1.6pt,line cap=round,line join=round](-0.13,0.0)--(-0.04,-0.09)--(0.13,0.1);}
\providecommand{\coche}{\popcheck{popgreen}}
\providecommand{\resultat}[1]{\par\smallskip\noindent\fcolorbox{popgreen}{popgreenL}{\parbox{\dimexpr\linewidth-2\fboxsep-2\fboxrule\relax}{\textbf{Résultat.} #1}}\par\smallskip}
\newcommand{\popboxhead}[3]{\poppill{#1}{#2}{white}\ifx\relax#3\relax\else\hspace{6pt}{\popxbold\fontsize{10.6}{12}\selectfont\color{popink}#3}\fi\par\vspace{7pt}}

\newtcolorbox{aretenir}[1][]{popbase,colframe=popgreen,before upper={\popboxhead{À retenir}{popgreen}{#1}}}
% Texte en allemand (document de lecture, dialogue, extrait) : cadre
% d'étude. Un titre optionnel (source, auteur) s'affiche dans l'en-tête.
\newtcolorbox{textebox}[1][]{popbase,colframe=popblue,before upper={\popboxhead{Text}{popblue}{#1}\begin{otherlanguage}{ngerman}},after upper={\end{otherlanguage}}}
% Marques ✓ / ✗ (forme correcte / fautive) souvent utilisées par les modèles
\providecommand{\cross}{\textcolor{poppink}{\ensuremath{\times}}}
\providecommand{\xmark}{\cross}\providecommand{\crossmark}{\cross}\providecommand{\cmark}{\ensuremath{\checkmark}}
% Ligne de réponse / justification à compléter (inventée par les modèles)
\providecommand{\justification}[1]{\par\noindent\textit{Justification~:} \dotfill\par}
% \textipa (package tipa non chargé) : la police ne couvre pas l'API -> simple passage
\providecommand{\textipa}[1]{#1}
% Soulignement (ulem non chargé) : \uline, \ul
\providecommand{\uline}[1]{\underline{#1}}\providecommand{\ul}[1]{\underline{#1}}
% \alert (beamer) inventée par les modèles : texte mis en évidence
\providecommand{\alert}[1]{\textbf{\textcolor{poppink}{#1}}}
% variantes ASCII d'ordinaux écrites en macro
\providecommand{\iere}{\textsuperscript{re}}\providecommand{\ieme}{\textsuperscript{e}}\providecommand{\ere}{\textsuperscript{re}}\providecommand{\eme}{\textsuperscript{e}}\providecommand{\er}{\textsuperscript{er}}\providecommand{\ie}{\textsuperscript{e}}
% Ordinaux français écrits en macro par les modèles : 1\ière, 3\ième
\newcommand{\ière}{\textsuperscript{re}}\newcommand{\ième}{\textsuperscript{e}}
% \exemple (inventée par les modèles) : étiquette « Exemple : »
\providecommand{\exemple}{\textbf{Exemple~:}\ }
% \square utilisable aussi hors mode math (cases à cocher)
\AtBeginDocument{\let\squareMath\square\renewcommand{\square}{\ensuremath{\squareMath}}}
% Mot ou phrase en allemand à l'intérieur d'un paragraphe français
\newcommand{\all}[1]{\ifmmode\text{\textit{#1}}\else\begingroup\selectlanguage{ngerman}\itshape #1\endgroup\fi}
\let\esp\all % alias toléré
\newtcolorbox{exemplebox}[1][]{popbase,colframe=poporange,before upper={\popboxhead{Exemple}{poporange}{#1}}}
\newtcolorbox{definition}[1][]{popbase,colframe=popblue,before upper={\popboxhead{Définition}{popblue}{#1}}}
\newtcolorbox{propriete}[1][]{popbase,colframe=poppurple,before upper={\popboxhead{Propriété}{poppurple}{#1}}}
\newtcolorbox{methode}[1][]{popbase,colframe=popgold,before upper={\popboxhead{Méthode}{popgold}{#1}}}
\newtcolorbox{attention}[1][]{popbase,colframe=poppink,before upper={\popboxhead{Attention}{poppink}{#1}}}
\newtcolorbox{experience}[1][]{popbase,colframe=popblue,colback=popblueL,before upper={\popboxhead{Expérience}{popblue}{#1}}}
% « Le savais-tu ? » : carte violette pleine (contexte, culture, Cameroun)
\newtcolorbox{savaistu}[1][]{popbase,colback=poppurple,colframe=popink,
  fontupper=\popsemi\fontsize{10.4}{14.2}\selectfont\color{white},
  before upper={\poppill{Le savais-tu\,?}{white}{poppurple}\ifx\relax#1\relax\else\hspace{6pt}{\popxbold\color{white}#1}\fi\par\vspace{7pt}}}
% Exercice résolu : énoncé en haut, \tcblower puis la solution sur fond crème
\newcounter{popexo}\newcounter{popexr}
\newtcolorbox{exoresolu}[1][]{popbase,colframe=poporange,colbacklower=popcream,
  segmentation style={popink,line width=1.2pt,dashed},
  before upper={\stepcounter{popexr}\popboxhead{Exercice résolu \thepopexr}{poporange}{#1}},
  before lower={\poppill{Solution}{popink}{popcream}\par\vspace{6pt}}}
% Exercice (non résolu en place) — la correction est dans la partie Corrigés
\newtcolorbox{exercice}[1][]{popbase,colframe=popink,
  before upper={\stepcounter{popexo}\popboxhead{Exercice \thepopexo}{popink}{#1}}}
% Corrigé
\newtcolorbox{corrige}[1][]{popbase,colframe=popgreen,colback=popgreenL,
  before upper={\popboxhead{Corrigé}{popgreen}{#1}}}
% Objectifs / prérequis / bilan
\newtcolorbox{objectifs}{popbase,colframe=popink,colback=poporangeL,
  before upper={\popboxhead{Objectifs de la leçon}{popink}{}}}
\newtcolorbox{prerequis}{popbase,colframe=popink,colback=popblueL,
  before upper={\popboxhead{Prérequis}{popblue}{}}}
\newtcolorbox{bilan}{popbase,colframe=popink,colback=white,boxrule=2.8pt,
  before upper={\popboxhead{Fiche bilan}{poporange}{L'essentiel à connaître pour l'examen}}}
% Figure encadrée avec légende
\newtcolorbox{popfigure}[1][]{enhanced,breakable=false,colback=white,colframe=popline,boxrule=1.4pt,arc=8pt,
  left=10pt,right=10pt,top=10pt,bottom=8pt,before skip=10pt,after skip=12pt,halign=center,
  lower separated=false,halign lower=center,
  fontlower=\popsemi\fontsize{9}{11.5}\selectfont\color{popmuted},
  #1}
\newcommand{\legende}[1]{\par\vspace{6pt}{\popsemi\fontsize{9}{11.5}\selectfont\color{popmuted}#1}}

% ============================================================
% Signature OnBuch+
% ============================================================
\newcommand{\poplogoinline}[1]{{\poplogo\fontsize{#1}{#1}\selectfont\color{popink}OnBuch\color{poporange}+}}
\newcommand{\signatureonbuch}{%
  \begin{tcolorbox}[enhanced,colback=popink,colframe=popink,boxrule=2.2pt,arc=10pt,
      drop shadow=poporange,left=14pt,right=14pt,top=10pt,bottom=10pt,before skip=0pt,after skip=0pt]
    \tikz[baseline=-0.7ex]\node[circle,fill=popgreen,draw=popcream,line width=1.3pt,minimum size=22pt,inner sep=0]{\popcheck{white}};%
    \hspace{9pt}\begin{minipage}[c]{0.62\linewidth}
      {\popxbold\fontsize{12}{13}\selectfont\color{popcream}Signé\ \textbullet\ {\poplogo OnBuch\color{poporange}+}}\par\vspace{2pt}
      {\raggedright\popsemi\fontsize{8}{10.5}\selectfont\color{popline}\MakeUppercase{Équipe pédagogique OnBuch+}\\\MakeUppercase{Conforme au programme officiel MINESEC}\par}
    \end{minipage}\hfill
    {\poplogo\fontsize{22}{22}\selectfont\color{popcream}OnBuch\color{poporange}+}
  \end{tcolorbox}%
}

% ============================================================
% Page de garde
% #1 titre · #2 matière · #3 niveau · #4 module · #5 n° leçon · #6 durée ·
% #7 description / objectifs (texte ou itemize)
% ============================================================
\newcommand{\pagegarde}[7]{%
  \thispagestyle{empty}
  \newgeometry{a4paper,margin=1.7cm,top=1.6cm,bottom=1.4cm}%
  \begin{tikzpicture}[remember picture,overlay]
    \fill[poporange!18] ([xshift=-3.2cm,yshift=-3.6cm]current page.north east) circle (4.6cm);
    \fill[poppurple!14] ([xshift=2.4cm,yshift=7.6cm]current page.south west) circle (3.8cm);
  \end{tikzpicture}%
  {\poplogo\fontsize{30}{30}\selectfont\color{popink}OnBuch\color{poporange}+}\hfill
  \raisebox{0.6ex}{\poppill{Cours signé \textbullet\ Premium}{popink}{popcream}}\par
  \vspace{1pt}\popsquiggle{poppurple}\par
  \vspace{8pt}
  \begin{tcolorbox}[enhanced,colback=poporange,colframe=popink,boxrule=2.8pt,arc=13pt,
      drop shadow=popink,left=18pt,right=18pt,top=12pt,bottom=13pt,after skip=10pt]
    \poppill{#2 \textbullet\ #3}{popink}{popcream}\hspace{5pt}\poppill{#4}{white}{popink}\par\vspace{8pt}
    {\popdisplay\fontsize{14}{14}\selectfont\color{popink}LEÇON #5}\par\vspace{4pt}
    {\raggedright\hyphenpenalty=10000\exhyphenpenalty=10000\popdisplay\fontsize{25}{27}\selectfont\color{popcream}\MakeUppercase{#1}\par}
    \vspace{8pt}
    {\popxbold\fontsize{9.5}{9.5}\selectfont\color{popink}\MakeUppercase{Durée conseillée : #6}}
  \end{tcolorbox}
  \begin{tcolorbox}[enhanced,colback=white,colframe=popink,boxrule=2.2pt,arc=10pt,
      drop shadow=popink,left=16pt,right=16pt,top=10pt,bottom=10pt,after skip=8pt]
    \poppill{Dans cette leçon}{poppurple}{white}\par\vspace{5pt}
    \popsemi\fontsize{9.3}{12.6}\selectfont\color{popink}#7
  \end{tcolorbox}
  \noindent\begin{minipage}[t]{0.487\linewidth}
    \begin{tcolorbox}[enhanced,colback=popgreen,colframe=popink,boxrule=2.2pt,arc=10pt,
        drop shadow=popink,left=11pt,right=11pt,top=10pt,bottom=10pt,height=3.1cm,valign=center]
      \tcbox[colback=white,colframe=popink,boxrule=1.3pt,arc=5pt,boxsep=0pt,nobeforeafter,
        left=3pt,right=3pt,top=3pt,bottom=3pt,valign=center]{\qrcode[height=1.9cm]{https://play.google.com/store/apps/details?id=com.luvvix.onbuch&hl=fr}}%
      \hspace{8pt}\begin{minipage}[c]{0.5\linewidth}
        {\popdisplay\fontsize{11}{12.5}\selectfont\color{white}\MakeUppercase{Télécharge OnBuch}}\par\vspace{4pt}
        {\raggedright\popsemi\fontsize{8.4}{11}\selectfont\color{white}Gratuit \textbullet\ Google Play\\Tuteur IA, annales, concours, résultats\par}
      \end{minipage}
    \end{tcolorbox}
  \end{minipage}\hfill
  \begin{minipage}[t]{0.487\linewidth}
    \begin{tcolorbox}[enhanced,colback=poppurple,colframe=popink,boxrule=2.2pt,arc=10pt,
        drop shadow=popink,left=11pt,right=11pt,top=10pt,bottom=10pt,height=3.1cm,valign=center]
      \tikz[baseline=-0.7ex]\node[circle,fill=white,draw=popink,line width=1.3pt,minimum size=1.6cm,inner sep=0]{\popdisplay\fontsize{24}{24}\selectfont\color{poppurple}+};%
      \hspace{8pt}\begin{minipage}[c]{0.6\linewidth}
        {\popdisplay\fontsize{11}{12.5}\selectfont\color{white}\MakeUppercase{Passe à OnBuch+}}\par\vspace{4pt}
        {\raggedright\popsemi\fontsize{8.4}{11}\selectfont\color{white}Tous les cours signés, Tuteur IA illimité, fiches PDF — onglet OnBuch+ de l'app\par}
      \end{minipage}
    \end{tcolorbox}
  \end{minipage}
  \vspace{8pt}
  \popcontenu
  \vfill
  \signatureonbuch
  \restoregeometry
  \newpage
}

% Rangée « Ce cours contient » (page de garde)
\newcommand{\poptile}[3]{% #1 couleur #2 gros texte #3 légende
  \begin{tcolorbox}[enhanced,colback=#1,colframe=popink,boxrule=2pt,arc=9pt,shadow={2.5pt}{-2.5pt}{0pt}{popink},
      left=6pt,right=6pt,top=9pt,bottom=9pt,halign=center,valign=center,height=1.75cm,width=\linewidth,nobeforeafter]
    {\popdisplay\fontsize{12.5}{13}\selectfont\color{white}\MakeUppercase{#2}}\par\vspace{3pt}
    {\centering\popsemi\fontsize{7.8}{9.5}\selectfont\color{white}#3\par}
  \end{tcolorbox}}
\newcommand{\popcontenu}{%
  \par\noindent{\popxboldls\fontsize{7.5}{7.5}\selectfont\color{popmuted}CE COURS SIGNÉ CONTIENT}\par\vspace{7pt}
  \noindent\begin{minipage}[t]{0.235\linewidth}\poptile{popblue}{Cours}{complet, illustré, pas à pas}\end{minipage}\hfill
  \begin{minipage}[t]{0.235\linewidth}\poptile{poporange}{Méthodes}{et exercices résolus}\end{minipage}\hfill
  \begin{minipage}[t]{0.235\linewidth}\poptile{poppink}{Exercices}{type examen + corrigés}\end{minipage}\hfill
  \begin{minipage}[t]{0.235\linewidth}\poptile{popgreen}{Fiche bilan}{l'essentiel à retenir}\end{minipage}\par}

% Sommaire
\newtcolorbox{sommaire}{enhanced,colback=white,colframe=popink,boxrule=2.2pt,arc=10pt,
  drop shadow=popink,left=18pt,right=18pt,top=13pt,bottom=13pt,before skip=6pt,after skip=20pt,
  before upper={\poppill{Sommaire}{poppurple}{white}\par\vspace{9pt}\popsemi\fontsize{10.6}{14.5}\selectfont\color{popink}}}

% Signature de fin de cours — poussée tout en bas de la dernière page
\newcommand{\findecours}{%
  \par\vfill\nopagebreak
  \begin{center}\popsquiggle{poporange}\end{center}\vspace{4pt}
  \signatureonbuch
}
\AtBeginDocument{\def\llbracket{\mathopen{[\mkern-3.5mu[}}\def\rrbracket{\mathclose{]\mkern-3.5mu]}}} % intervalles d'entiers (glyphe ⟦ absent de la police)
% Glyphes absents de Fira Math : redéfinis à partir de glyphes disponibles
\AtBeginDocument{%
  \def\setminus{\mathbin{\backslash}}\let\smallsetminus\setminus
  \def\vdots{\mathord{\vbox{\baselineskip3.2pt\lineskiplimit0pt\kern4pt\hbox{.}\hbox{.}\hbox{.}}}}%
  \def\ddots{\mathinner{\mkern1mu\raise6.5pt\hbox{.}\mkern2mu\raise3.6pt\hbox{.}\mkern2mu\raise0.7pt\hbox{.}\mkern1mu}}%
  \def\triangle{\mathord{\scalebox{0.9}{$\symup{\Delta}$}}}%
  \def\nabla{\mathord{\raisebox{\depth}{\scalebox{1}[-1]{$\symup{\Delta}$}}}}%
  \def\checkmark{\popcheck{popdark}}%
  \def\star{\mathbin{\ast}}\def\bigstar{\mathord{\ast}}%
  \def\diamond{\mathbin{\scalebox{0.6}{$\Diamond$}}}%
  \def\bigcup{\mathop{\mathchoice{\raisebox{-0.3em}{\scalebox{1.7}{$\cup$}}}{\scalebox{1.25}{$\cup$}}{\cup}{\cup}}\displaylimits}%
  \def\bigcap{\mathop{\mathchoice{\raisebox{-0.3em}{\scalebox{1.7}{$\cap$}}}{\scalebox{1.25}{$\cap$}}{\cap}{\cap}}\displaylimits}%
  \def\lesseqqgtr{\mathrel{\lessgtr}}\def\bigoplus{\mathop{\oplus}\displaylimits}%
}
\providecommand{\ssi}{si et seulement si} % abréviation fréquente
\AtBeginDocument{\providecommand{\wideparen}{\overparen}} % arc de cercle

\begin{document}

\pagegarde{Marketing et publicité}{Allemand}{Tle A}{Chapitre 2 : Vie économique}{AL07}{2 h}{Cette leçon de type texte propose la lecture et le commentaire d'un texte original sur le marketing et la publicité dans les pays germanophones. Elle développe le lexique économique thématique et consolide deux points de grammaire essentiels au Bac : l'impératif et l'infinitif avec zu exprimant la finalité (um ... zu, damit).
\vspace{4pt}
\begin{multicols}{2}\small
\begin{itemize}
\item À la fin de cette leçon, je sais comprendre globalement puis en détail un texte allemand sur le marketing et la publicité
\item À la fin de cette leçon, je sais employer correctement le vocabulaire thématique (die Werbung, das Produkt, der Kunde, das Angebot, der Verkauf, die Marke) avec articles et pluriels
\item À la fin de cette leçon, je sais former et utiliser l'impératif aux trois personnes (du, ihr, Sie) dans des slogans publicitaires
\item À la fin de cette leçon, je sais construire l'infinitif avec zu et exprimer la finalité avec um ... zu et damit
\item À la fin de cette leçon, je sais distinguer um ... zu et damit selon que les sujets sont identiques ou différents
\item À la fin de cette leçon, je sais répondre par écrit et à l'oral à des questions de compréhension sur un texte
\end{itemize}\end{multicols}}

\begin{sommaire}
\begin{multicols}{2}
\begin{enumerate}
\item Texte support : « Werbung überall »
\item Compréhension du texte
\item Lexique thématique : Marketing und Werbung
\item Grammaire 1 : L'impératif (der Imperativ)
\item Grammaire 2 : L'infinitif avec zu et la finalité (um ... zu / damit)
\item Méthode du commentaire et production guidée
\item Activité d'intégration
\item Méthodes et exercices résolus
\item Exercices
\item Corrigés des exercices
\item Fiche bilan
\end{enumerate}
\end{multicols}
\end{sommaire}

\begin{objectifs}
\begin{itemize}
\item À la fin de cette leçon, je sais comprendre globalement puis en détail un texte allemand sur le marketing et la publicité
\item À la fin de cette leçon, je sais employer correctement le vocabulaire thématique (die Werbung, das Produkt, der Kunde, das Angebot, der Verkauf, die Marke) avec articles et pluriels
\item À la fin de cette leçon, je sais former et utiliser l'impératif aux trois personnes (du, ihr, Sie) dans des slogans publicitaires
\item À la fin de cette leçon, je sais construire l'infinitif avec zu et exprimer la finalité avec um ... zu et damit
\item À la fin de cette leçon, je sais distinguer um ... zu et damit selon que les sujets sont identiques ou différents
\item À la fin de cette leçon, je sais répondre par écrit et à l'oral à des questions de compréhension sur un texte
\item À la fin de cette leçon, je sais rédiger un résumé et un court paragraphe de commentaire sur un texte publicitaire
\item À la fin de cette leçon, je sais créer un slogan publicitaire en allemand en utilisant l'impératif et la finalité
\end{itemize}
\end{objectifs}

\begin{prerequis}
\begin{itemize}
\item Conjugaison des verbes au présent (verbes réguliers, irréguliers, à particule) — classe de Première
\item Place du verbe et inversion du sujet dans la phrase déclarative et interrogative
\item Articles définis et indéfinis, déclinaison de base du nominatif et de l'accusatif
\item Vocabulaire général de la vie quotidienne et de la consommation (kaufen, Geld, Geschäft)
\item Subordonnées introduites par une conjonction (weil, dass) et place du verbe en fin de proposition
\end{itemize}
\end{prerequis}

\newpage

\coursec{Texte support : « Werbung überall »}

\courssub{Situation d'accroche et problématique}

À Douala comme à Berlin, impossible d'échapper à la publicité : sur les panneaux d'affichage du marché Mokolo, à la radio, sur les réseaux sociaux, les marques nous interpellent sans cesse. En Allemagne, la publicité est une industrie puissante : les entreprises dépensent des milliards d'euros pour convaincre les consommateurs. Mais comment une publicité réussit-elle à nous faire acheter ? Quelles stratégies se cachent derrière un simple slogan comme \all{„Komm gut an!“} ?

Dans cette leçon, nous allons lire un texte sur le marketing et la publicité, enrichir notre vocabulaire économique et apprendre à formuler des injonctions et des buts en allemand — exactement ce que font les publicitaires !

\textbf{Problématique :} \emph{Comment la publicité influence-t-elle nos choix de consommation, et quels procédés linguistiques utilise-t-elle pour nous convaincre ?}

\courssub{Avant la lecture : activer ses connaissances}

Avant de lire, observons le titre : \all{„Werbung überall – Wie Firmen uns zum Kaufen bringen“} (« La publicité partout – Comment les entreprises nous amènent à acheter »). Le titre annonce déjà deux idées : la publicité est \cle{omniprésente} (\all{überall}) et elle a un \cle{but} : nous pousser à l'achat (\all{zum Kaufen bringen}).

Bonne nouvelle pour les francophones : ce thème contient de nombreux \cle{mots transparents}, c'est-à-dire des mots allemands proches du français ou de l'anglais que vous connaissez déjà :

\begin{center}
\begin{tabularx}{\linewidth}{|X|X|X|}
\hline
\rowcolor{popblueL} Mot allemand & Français & Remarque \\
\hline
\all{das Produkt, -e} & le produit & mot transparent \\
\hline
\all{die Marke, -n} & la marque & mot transparent \\
\hline
\all{der Kunde, -n} & le client & déjà connu \\
\hline
\all{kaufen} & acheter & déjà connu \\
\hline
\all{das Radio, -s} & la radio & mot transparent \\
\hline
\all{der Preis, -e} & le prix & faux ami partiel : aussi « le prix » d'un concours \\
\hline
\end{tabularx}
\end{center}

\begin{attention}[Faux amis à surveiller]
\all{das Angebot} ne signifie pas « l'ange » mais \textbf{l'offre} (promotion, proposition commerciale). \all{die Firma, die Firmen} signifie \textbf{l'entreprise}, et non « la firme » au sens juridique seulement. \all{bekannt} veut dire \textbf{connu, célèbre} : \all{ein Produkt bekannt machen} = « faire connaître un produit ».
\end{attention}

\courssub{Lecture du texte}

\textbf{Première lecture silencieuse} (2 minutes) : lisez le texte sans dictionnaire, en vous aidant uniquement du glossaire. \textbf{Deuxième lecture à voix haute} : le professeur lit, puis des élèves lisent à tour de rôle, en soignant la prononciation : \all{Werbung} se prononce « vère-boung », \all{Zielgruppe} « tsil-groupé », \all{überzeugen} « u-ber-tsoi-guen ».

\begin{textebox}[Werbung überall – Wie Firmen uns zum Kaufen bringen]
Schon am Morgen sehen wir die erste Werbung: im Radio, auf dem Handy, an der Bushaltestelle. Firmen geben jedes Jahr Milliarden aus, um ihre Produkte bekannt zu machen. Sie wollen den Kunden überzeugen, damit er ihre Marke kauft. Ein guter Slogan ist kurz und direkt: „Probier es einfach!“ oder „Kauf jetzt und spare Geld!“ Die Werbung kennt ihre Zielgruppe genau: Jugendliche sieht man in den sozialen Medien, ältere Menschen eher im Fernsehen. Manche Kritiker sagen: Die Werbung manipuliert uns. Wir kaufen Produkte, die wir nicht brauchen. Trotzdem kann man ohne Werbung kaum ein neues Produkt verkaufen. Sie informiert auch über Angebote und Preise. Die Frage bleibt: Beeinflusst die Werbung unsere Wünsche – oder zeigt sie nur, was wir schon wollen?
\end{textebox}

\textbf{Glossaire (mots difficiles) :}

\begin{center}
\begin{tabularx}{\linewidth}{|X|X|}
\hline
\rowcolor{popblueL} Mot allemand & Traduction française \\
\hline
\all{die Werbung} (Sg.) & la publicité \\
\hline
\all{die Bushaltestelle, -n} & l'arrêt de bus \\
\hline
\all{ausgeben (gibt aus, gab aus, hat ausgegeben)} & dépenser (de l'argent) \\
\hline
\all{bekannt machen} & faire connaître, rendre célèbre \\
\hline
\all{überzeugen (+ Akk.)} & convaincre (quelqu'un) \\
\hline
\all{der Slogan, -s} & le slogan \\
\hline
\all{sparen} & économiser, épargner \\
\hline
\all{die Zielgruppe, -n} & le groupe cible \\
\hline
\all{die sozialen Medien} (Pl.) & les réseaux sociaux \\
\hline
\all{manipulieren} & manipuler \\
\hline
\all{brauchen} & avoir besoin de \\
\hline
\all{der Verbraucher, die Verbraucher} & le consommateur \\
\hline
\all{der Einfluss} (\all{die Einflüsse}) & l'influence \\
\hline
\all{der Wunsch} (\all{die Wünsche}) & le souhait, le désir \\
\hline
\end{tabularx}
\end{center}

\begin{definition}[Das Marketing]
\all{Das Marketing} (sans pluriel) désigne l'\textbf{ensemble des méthodes qu'une entreprise utilise pour vendre un produit} : l'étude du marché (\all{die Marktforschung}), la fixation du prix (\all{der Preis}), la publicité (\all{die Werbung}) et la distribution (\all{der Vertrieb}). La publicité n'est donc qu'\emph{une partie} du marketing — ne confondez pas les deux termes dans vos réponses au Bac.
\end{definition}

\courssub{Première compréhension globale}

Après la première lecture, répondez à ces trois questions fondamentales :

\begin{enumerate}
\item \textbf{De quoi parle le texte ?} — \all{Der Text handelt von der Werbung und ihren Strategien.} (Le texte parle de la publicité et de ses stratégies.)
\item \textbf{Qui parle ?} — Un journaliste ou un auteur anonyme qui décrit la publicité de manière objective, mais qui pose à la fin une question critique.
\item \textbf{Quel est le but de l'auteur ?} — \all{Der Autor will informieren und zum Nachdenken anregen.} (L'auteur veut informer et inciter à réfléchir.) Le texte est à la fois \cle{informatif} (il explique les stratégies) et \cle{critique} (il mentionne la manipulation).
\end{enumerate}

\begin{methode}[Lire un texte au Bac : les 5 étapes]
\begin{enumerate}
\item \textbf{Lire le titre et la source} : ils annoncent le thème et le type de texte.
\item \textbf{Première lecture rapide} : chercher l'idée générale, sans s'arrêter sur les mots inconnus.
\item \textbf{Deuxième lecture avec le glossaire} : comprendre les détails, phrase par phrase.
\item \textbf{Souligner les mots-clés} : les noms importants, les verbes d'action, les connecteurs (\all{damit}, \all{um ... zu}, \all{trotzdem}).
\item \textbf{Répondre aux questions en phrases complètes}, en allemand, en reformulant — jamais en recopiant mot à mot de longs passages du texte.
\end{enumerate}
\end{methode}

\courssub{Deuxième lecture détaillée : repérage linguistique}

Relisons le texte avec un œil de grammairien. Deux phénomènes dominent dans ce texte publicitaire — et ce sont précisément les deux points de grammaire de notre leçon :

\textbf{1. Les slogans à l'impératif.} La publicité donne des ordres amicaux au consommateur :

\begin{exemplebox}[Les impératifs du texte]
\all{„Probier es einfach!“} — « Essaie-le simplement ! » (impératif \all{du} de \all{probieren})

\all{„Kauf jetzt und spare Geld!“} — « Achète maintenant et économise de l'argent ! » (impératifs \all{kauf} et \all{spare}, forme abrégée typique des slogans)
\end{exemplebox}

Remarquez la place du verbe : à l'impératif, le verbe se place \cle{en première position} de la phrase, comme en français (« Achète ! » = \all{Kauf !}).

\textbf{2. Les phrases de finalité.} La publicité explique toujours \emph{pourquoi} les entreprises agissent :

\begin{exemplebox}[La finalité dans le texte]
\all{Firmen geben Milliarden aus, um ihre Produkte bekannt zu machen.} — « Les entreprises dépensent des milliards \textbf{pour} faire connaître leurs produits. » (structure \all{um ... zu} + infinitif)

\all{Sie wollen den Kunden überzeugen, damit er ihre Marke kauft.} — « Elles veulent convaincre le client \textbf{afin qu'}il achète leur marque. » (subordonnée avec \all{damit} + verbe conjugué à la fin)
\end{exemplebox}

\begin{aretenir}[Deux structures, une même idée de but]
\begin{itemize}
\item \all{um ... zu} + infinitif : le sujet des deux actions est \textbf{le même} (les entreprises dépensent, les entreprises font connaître).
\item \all{damit} + verbe conjugué en fin de phrase : les sujets sont \textbf{différents} (les entreprises convainquent, le client achète).
\end{itemize}
Ces deux structures seront étudiées en détail dans la section grammaire de cette leçon.
\end{aretenir}

\courssub{Carte mentale : les idées du texte}

Le texte s'organise autour de quatre grandes idées, que cette carte mentale résume :

\begin{popfigure}
\begin{tikzpicture}[
  mindmap,
  grow cyclic,
  every node/.style={concept, align=center, text=white},
  root concept/.append style={concept color=popblue, font=\large\bfseries},
  level 1/.append style={level distance=3.2cm, sibling angle=90},
  level 2/.append style={level distance=2.2cm, sibling angle=50, font=\small}
]
\node[root concept]{die Werbung}
  child[concept color=poporange]{ node {Strategien}
    child { node {Slogans} }
    child { node {Image der Marke} }
  }
  child[concept color=popgreen]{ node {Zielgruppe}
    child { node {Jugendliche} }
    child { node {ältere Menschen} }
  }
  child[concept color=poppurple]{ node {Medien}
    child { node {Radio, Handy} }
    child { node {Fernsehen, soziale Medien} }
  }
  child[concept color=poppink]{ node {Kritik}
    child { node {Manipulation} }
    child { node {Konsum} }
  };
\end{tikzpicture}
\legende{Carte mentale du texte : les quatre axes de la publicité}
\end{popfigure}

\courssub{Entraînement à la compréhension}

\begin{exoresolu}[Comprendre l'idée générale]
\textbf{Énoncé :} Répondez en allemand, par une phrase complète : \all{Was ist die Hauptidee des Textes?} (Quelle est l'idée principale du texte ?)

\tcblower

\textbf{Solution détaillée :} On cherche d'abord le thème (\all{die Werbung}), puis ce que l'auteur en dit (ses stratégies et sa critique). On formule une phrase complète avec un verbe conjugué en deuxième position :

\all{Die Hauptidee des Textes ist, dass die Werbung überall ist und dass Firmen mit klugen Strategien die Kunden zum Kaufen bringen wollen.}

« L'idée principale du texte est que la publicité est partout et que les entreprises veulent pousser les clients à acheter grâce à des stratégies intelligentes. »

\textbf{Méthode :} introduire la réponse par \all{Die Hauptidee ist, dass ...} avec le verbe de la subordonnée \textbf{à la fin} — une structure qui rapporte des points au Bac.
\end{exoresolu}

\begin{exercice}[À vous de jouer : compréhension détaillée]
Répondez en allemand, en phrases complètes :
\begin{enumerate}
\item \all{Wo sehen wir am Morgen die erste Werbung?} (Où voyons-nous la première publicité le matin ?)
\item \all{Warum geben Firmen so viel Geld aus?} (Pourquoi les entreprises dépensent-elles tant d'argent ?)
\item \all{Was sagen die Kritiker über die Werbung?} (Que disent les critiques de la publicité ?)
\item Trouvez dans le texte deux slogans à l'impératif et traduisez-les.
\item Relevez la phrase avec \all{damit} et expliquez pourquoi l'auteur n'utilise pas \all{um ... zu}.
\end{enumerate}
\end{exercice}

\begin{corrige}[Exercice « À vous de jouer »]
\begin{enumerate}
\item \all{Am Morgen sehen wir die erste Werbung im Radio, auf dem Handy und an der Bushaltestelle.} (« Le matin, nous voyons la première publicité à la radio, sur le portable et à l'arrêt de bus. »)
\item \all{Firmen geben so viel Geld aus, um ihre Produkte bekannt zu machen und den Kunden zu überzeugen.} (« Les entreprises dépensent tant d'argent pour faire connaître leurs produits et convaincre le client. »)
\item \all{Die Kritiker sagen, dass die Werbung uns manipuliert und dass wir Produkte kaufen, die wir nicht brauchen.} (« Les critiques disent que la publicité nous manipule et que nous achetons des produits dont nous n'avons pas besoin. »)
\item \all{„Probier es einfach!“} = « Essaie-le simplement ! » ; \all{„Kauf jetzt und spare Geld!“} = « Achète maintenant et économise de l'argent ! »
\item \all{Sie wollen den Kunden überzeugen, damit er ihre Marke kauft.} On utilise \all{damit} car les sujets sont différents : \all{sie} (les entreprises) et \all{er} (le client). Avec \all{um ... zu}, le sujet devrait être identique dans les deux parties.
\end{enumerate}
\end{corrige}

\begin{savaistu}[La publicité en Allemagne]
L'Allemagne est le premier marché publicitaire d'Europe : les entreprises y investissent chaque année des sommes considérables dans les spots télévisés, les affiches et la publicité en ligne. La ville de \textbf{Düsseldorf}, sur le Rhin, est surnommée la \emph{capitale allemande de la publicité} : c'est là que se sont installées de nombreuses grandes agences. Et au Cameroun ? La publicité envahit aussi nos écrans et nos rues — pensez aux panneaux des opérateurs de téléphonie mobile à Yaoundé et à Douala : comme en Allemagne, ils utilisent des slogans courts, directs, à l'impératif !
\end{savaistu}

\begin{attention}[Erreurs fréquentes des francophones]
\begin{itemize}
\item \all{die Werbung} est presque toujours au \textbf{singulier} en allemand, même quand le français dit « les publicités ». On dit : \all{Ich sehe viel Werbung.} (« Je vois beaucoup de publicités. »)
\item N'oubliez pas la \textbf{majuscule} à tous les noms : \all{das Produkt}, \all{der Kunde}, \all{die Marke} — jamais \emph{produkt}, \emph{kunde}.
\item Attention à la place du verbe dans \all{damit er ihre Marke kauft} : le verbe conjugué va \textbf{à la toute fin}, contrairement au français « afin qu'il achète ».
\item \all{ausgeben} = dépenser (de l'argent) ; ne le confondez pas avec \all{verbringen} (passer du temps) ni avec le faux ami « donner ».
\end{itemize}
\end{attention}

\coursec{Lexique thématique : Marketing und Werbung}

\courssub{Familles de mots et collocations}

Le vocabulaire économique s'organise en familles. Maîtriser les dérivés vous permet de varier vos formulations à l'écrit comme à l'oral.

\begin{center}
\begin{tabularx}{\linewidth}{|X|X|X|}
\hline
\rowcolor{popblueL} Verbe de base & Noms (avec article \& pluriel) & Adjectifs / Participles \\
\hline
\all{werben} (bewirbt, warb, hat beworben) & \all{die Werbung} (Sg.), \all{der Werbespot, -s}, \all{die Werbetafel, -n}, \all{der Werbeblock, -s} & \all{werbend}, \all{werbefrei} \\
\hline
\all{verkaufen} & \all{der Verkauf}, \all{der Verkäufer, -}, \all{die Verkäuferin, -nen}, \all{das Verkaufsgespräch, -e} & \all{verkäuflich}, \all{verkauft} \\
\hline
\all{produzieren} & \all{das Produkt, -e}, \all{die Produktion}, \all{der Produzent, -en}, \all{die Produktpalette, -n} & \all{produktiv}, \all{produziert} \\
\hline
\all{konsumieren} & \all{der Konsum}, \all{der Konsument, -en}, \all{das Konsumverhalten} & \all{konsumierend}, \all{konsumierbar} \\
\hline
\all{überzeugen} & \all{die Überzeugung}, \all{der Überzeugungsversuch, -e} & \all{überzeugend}, \all{überzeugt} \\
\hline
\all{manipulieren} & \all{die Manipulation}, \all{der Manipulator, -en} & \all{manipulativ}, \all{manipuliert} \\
\hline
\end{tabularx}
\end{center}

\begin{propriete}[Collocations utiles pour le Bac (expression écrite/orale)]
Apprenez ces blocs figés (\cle{chunks}) par cœur :
\begin{itemize}
\item \all{Werbung machen für + Akk.} (faire de la pub pour)
\item \all{eine Kampagne starten / durchführen} (lancer / mener une campagne)
\item \all{die Zielgruppe ansprechen / erreichen} (cibler / atteindre le groupe cible)
\item \all{ein Produkt auf den Markt bringen} (lancer un produit sur le marché)
\item \all{den Umsatz steigern / ankurbeln} (augmenter / stimuler le chiffre d'affaires)
\item \all{Marktforschung betreiben} (faire une étude de marché)
\item \all{Markenimage pflegen} (soigner l'image de marque)
\end{itemize}
\end{propriete}

\begin{exercice}[Entraînement lexical]
Complétez les phrases avec le mot juste de la famille indiquée (mettez la bonne forme) :
\begin{enumerate}
\item Die Firma \all{\_\_\_\_\_\_\_\_\_\_} (werben) viel Geld für das neue Handy.
\item Der \all{\_\_\_\_\_\_\_\_\_\_} (verkaufen) des Autos war sehr erfolgreich.
\item Wir müssen die \all{\_\_\_\_\_\_\_\_\_\_} (produzieren) steigern.
\item Das ist ein sehr \all{\_\_\_\_\_\_\_\_\_\_} (überzeugen) Argument.
\item Kritiker warnen vor der \all{\_\_\_\_\_\_\_\_\_\_} (manipulieren) durch Werbung.
\end{enumerate}
\end{exercice}

\begin{corrige}[Entraînement lexical]
\begin{enumerate}
\item \all{bewirbt} (verbe \all{bewerben} + Akk., ou \all{Werbung machen für} — ici verbe simple) / \all{macht ... Werbung}
\item \all{Verkauf} (nom masculin)
\item \all{Produktion} (nom féminin)
\item \all{überzeugend} (adjectif)
\item \all{Manipulation} (nom féminin)
\end{enumerate}
\end{corrige}

\coursec{Grammaire 1 : L'impératif (der Imperativ)}

\courssub{Formation}

L'impératif exprime un ordre, une consigne, un conseil ou une invitation. Il existe trois formes selon la personne à qui l'on s'adresse.

\begin{propriete}[Tableau de formation de l'impératif]
\begin{center}
\begin{tabular}{|l|c|c|c|}
\hline
\rowcolor{popblueL} \textbf{Verbe} & \textbf{du (familier sg)} & \textbf{ihr (familier pl)} & \textbf{Sie (formel sg/pl)} \\
\hline
\all{regulär: kaufen} & \all{kauf(e)!} & \all{kauft!} & \all{kaufen Sie!} \\
\hline
\all{regulär: arbeiten} & \all{arbeite!} & \all{arbeitet!} & \all{arbeiten Sie!} \\
\hline
\all{verbe à particule: ausgeben} & \all{gib aus!} & \all{gebt aus!} & \all{geben Sie aus!} \\
\hline
\all{verbe à voyelle changeante: lesen} & \all{lies!} & \all{lest!} & \all{lesen Sie!} \\
\hline
\all{verbe à voyelle changeante: fahren} & \all{fähr! / fahre!} & \all{fährt!} & \all{fahren Sie!} \\
\hline
\all{sein} & \all{sei!} & \all{seid!} & \all{seien Sie!} \\
\hline
\all{haben} & \all{hab(e)!} & \all{habt!} & \all{haben Sie!} \\
\hline
\all{werden} & \all{werde!} & \all{werdet!} & \all{werden Sie!} \\
\hline
\end{tabular}
\end{center}
\end{propriete}

\begin{aretenir}[Règles clés pour l'impératif \all{du}]
\begin{enumerate}
\item Radical du présent sans \all{-st} : \all{du kaufst} $\rightarrow$ \all{kauf!}
\item Verbes en \all{-eln}, \all{-ern} : gardent souvent le \all{-e} (\all{handeln} $\rightarrow$ \all{handle!}, \all{lächeln} $\rightarrow$ \all{lächele!}).
\item Verbes en \all{-igen}, \all{-nen}, \all{-men} : gardent le \all{-e} pour l'euphonie (\all{regnen} $\rightarrow$ \all{regne!}, \all{öffnen} $\rightarrow$ \all{öffne!}).
\item Verbes à mutation vocalique (\all{a$\rightarrow$ä}, \all{e$\rightarrow$i/ie}) : la mutation \textbf{subsiste} à l'impératif \all{du} (\all{fahren} $\rightarrow$ \all{fähr!}, \all{lesen} $\rightarrow$ \all{lies!}, \all{stoßen} $\rightarrow$ \all{stöß!}). Pas de \all{-e} final.
\item Verbes en \all{-x, -z, -ß, -tz, -ss} : pas de \all{-e} (\all{setzen} $\rightarrow$ \all{setz!}, \all{heißen} $\rightarrow$ \all{heiß!}).
\end{enumerate}
\end{aretenir}

\begin{aretenir}[Position de la particule séparable]
À l'impératif, la particule va \textbf{à la fin de la phrase} :
\all{Ruf mich an!} (Appelle-moi !) — \all{Gib das Geld aus!} (Dépense l'argent !)
\end{aretenir}

\begin{exemplebox}[Impératif dans la publicité (style télégraphique)]
Souvent, le pronom \all{du} / \all{ihr} / \all{Sie} est omis, le verbe reste en position 1 :
\all{„Probier es!“} (au lieu de \all{Probier du es!})
\all{„Kaufen Sie jetzt!“} (formel, pronom \all{Sie} conservé après le verbe)
\all{„Spare 20\%!“} (style nominal / infinitif utilisé comme impératif impersonnel — très fréquent sur les affiches)
\end{exemplebox}

\begin{exercice}[Entraînement : L'impératif]
Mettez les verbes entre parenthèses à l'impératif pour la personne indiquée :
\begin{enumerate}
\item \all{\_\_\_\_\_\_\_\_\_\_} (lesen) Sie die Zeitung! (formel)
\item \all{\_\_\_\_\_\_\_\_\_\_} (kaufen) du das Buch! (familier sg)
\item \all{\_\_\_\_\_\_\_\_\_\_} (ausgeben) ihr nicht so viel Geld! (familier pl)
\item \all{\_\_\_\_\_\_\_\_\_\_} (sein) du vorsichtig! (familier sg)
\item \all{\_\_\_\_\_\_\_\_\_\_} (anrufen) Sie mich morgen! (formel, verbe séparable)
\end{enumerate}
\end{exercice}

\begin{corrige}[Entraînement : L'impératif]
\begin{enumerate}
\item \all{Lesen} Sie die Zeitung!
\item \all{Kauf(e)} du das Buch! (les deux formes correctes, \all{Kauf} plus courant)
\item \all{Gebt} ihr nicht so viel Geld \all{aus}!
\item \all{Sei} du vorsichtig!
\item \all{Rufen} Sie mich morgen \all{an}!
\end{enumerate}
\end{corrige}

\coursec{Grammaire 2 : L'infinitif avec \all{zu} et la finalité (\all{um ... zu} / \all{damit})}

\courssub{L'infinitif avec \all{zu} (der Infinitiv mit \all{zu})}

En allemand, l'infinitif présent est souvent précédé de \all{zu} quand il dépend d'un verbe, d'un adjectif ou d'un nom.

\begin{propriete}[Construction de \all{zu} + infinitif]
\begin{itemize}
\item \textbf{Verbe simple :} \all{zu} + infinitif $\rightarrow$ \all{zu kaufen}, \all{zu lesen}.
\item \textbf{Verbe à particule séparable :} \all{zu} s'intercale \textbf{entre la particule et le radical} $\rightarrow$ \all{aufzumachen}, \all{auszugeben}, \all{anzurufen}.
\item \textbf{Verbe composé (préfixe inséparable) :} \all{zu} devant le verbe entier $\rightarrow$ \all{zu verkaufen}, \all{zu besuchen}, \all{zu manipulieren}.
\end{itemize}
\end{propriete}

\begin{exemplebox}[Verbes / adjectifs / noms appelant \all{zu} + infinitif]
\begin{itemize}
\item \textbf{Verbes :} \all{versuchen} (essayer), \all{planen} (planifier), \all{hoffen} (espérer), \all{versprechen} (promettre), \all{vergessen} (oublier), \all{unterlassen} (s'abstenir), \all{bitten} (prier), \all{auffordern} (inviter), \all{raten} (conseiller).
\item \textbf{Adjectifs :} \all{schön}, \all{leicht}, \all{schwer}, \all{möglich}, \all{notwendig}, \all{bereit} (+ \all{zu}).
\item \textbf{Noms :} \all{die Absicht}, \all{der Plan}, \all{die Möglichkeit}, \all{die Absicht}, \all{der Versuch} (+ \all{zu}).
\end{itemize}
\all{Die Firma plant, die Preise zu senken.} (L'entreprise prévoit de baisser les prix.)
\all{Es ist wichtig, die Zielgruppe zu kennen.} (Il est important de connaître le groupe cible.)
\end{exemplebox}

\courssub{La finalité : \all{um ... zu} vs \all{damit}}

Pour exprimer le \textbf{but} (la finalité), l'allemand utilise deux structures principales. Le choix dépend de l'identité du sujet.

\begin{propriete}[Règle de choix : \all{um ... zu} ou \all{damit} ?]
\begin{center}
\begin{tabularx}{\linewidth}{|X|X|X|}
\hline
\rowcolor{popblueL} Structure & Condition sur le sujet & Exemple \\
\hline
\all{um ... zu} + infinitif & \textbf{Sujet identique} dans la principale et l'infinitif & \all{Ich lerne Deutsch, um in Deutschland zu arbeiten.} \\
\hline
\all{damit} + verbe conjugué (fin de phrase) & \textbf{Sujets différents} OU même sujet mais accent sur le but/résultat & \all{Ich lerne Deutsch, damit ich in Deutschland arbeiten kann.} \\
\hline
\end{tabularx}
\end{center}
\end{propriete}

\begin{aretenir}[Points de vigilance pour le Bac]
\begin{itemize}
\item \all{um ... zu} : le verbe est à l'infinitif \textbf{à la fin} de la subordonnée. \all{zu} se place devant l'infinitif (règle de l'intercalation pour verbes séparables : \all{um das Radio anzumachen}).
\item \all{damit} : conjonction de subordination $\rightarrow$ verbe conjugué \textbf{rejeté à la toute fin}. Le sujet \textbf{doit être exprimé} (pas de sujet caché).
\item \textbf{Négation :} \all{um nicht zu} / \all{damit ... nicht}.
\item \textbf{Antériorité / Postériorité :} Si l'action du but est antérieure, on utilise \all{um ... zu haben} / \all{zu sein} (infinitif parfait) ou \all{damit ... hat / ist}.
\end{itemize}
\end{aretenir}

\begin{exemplebox}[Exemples du thème « Publicité »]
\begin{itemize}
\item \all{Die Firma gibt Geld aus, um das Produkt bekannt zu machen.} (Sujet identique : \all{Die Firma})
\item \all{Die Firma macht Werbung, damit die Kunden das Produkt kaufen.} (Sujets différents : \all{Die Firma} vs \all{die Kunden})
\item \all{Wir schalten einen Spot, damit unsere Marke bekannter wird.} (Sujets différents : \all{Wir} vs \all{unsere Marke})
\item \all{Er spart Geld, um sich ein Handy zu kaufen.} (Sujet identique : \all{Er})
\end{itemize}
\end{exemplebox}

\begin{attention}[Pièges classiques des francophones]
\begin{itemize}
\item \textbf{Confusion \all{um ... zu} / \all{damit} :} En français, « pour » + infinitif couvre les deux cas. En allemand, \textbf{testez le sujet} : si je peux dire « pour que \textbf{je/il/ils} fasse(nt) » $\rightarrow$ \all{damit}. Si « pour faire » (même acteur) $\rightarrow$ \all{um ... zu}.
\item \textbf{Place de \all{zu} :} \all{anzurufen} (pas \all{zu anrufen}), \all{auszugeben} (pas \all{zu ausgeben}).
\item \textbf{Oubli du verbe à la fin avec \all{damit} :} \all{damit er kauft} (pas \all{damit er kauft das Produkt} — l'objet \all{das Produit} vient \textbf{avant} le verbe final).
\item \textbf{Sujet manquant avec \all{damit} :} \all{Ich rufe an, damit wir uns treffen.} (Pas \all{damit uns treffen}).
\end{itemize}
\end{attention}

\begin{exercice}[Entraînement : Finalité]
Reliez les deux phrases par \all{um ... zu} ou \all{damit} (choisissez la structure correcte).
\begin{enumerate}
\item Die Firma will Gewinn machen. Die Firma senkt die Kosten.
\item Der Kunde soll das Produkt kaufen. Die Werbung zeigt Vorteile.
\item Ich gehe in den Supermarkt. Ich will Brot kaufen.
\item Der Lehrer erklärt die Grammatik. Die Schüler sollen sie verstehen.
\item Wir lernen Vokabeln. Wir wollen die Klausur bestehen.
\end{enumerate}
\end{exercice}

\begin{corrige}[Entraînement : Finalité]
\begin{enumerate}
\item \all{Die Firma senkt die Kosten, um Gewinn zu machen.} (Sujet identique : \all{Die Firma})
\item \all{Die Werbung zeigt Vorteile, damit der Kunde das Produkt kauft.} (Sujets différents : \all{Die Werbung} vs \all{der Kunde})
\item \all{Ich gehe in den Supermarkt, um Brot zu kaufen.} (Sujet identique : \all{Ich})
\item \all{Der Lehrer erklärt die Grammatik, damit die Schüler sie verstehen.} (Sujets différents : \all{Der Lehrer} vs \all{die Schüler})
\item \all{Wir lernen Vokabeln, um die Klausur zu bestehen.} (Sujet identique : \all{Wir}) — ou \all{damit wir die Klausur bestehen} (possible si insistance sur le résultat).
\end{enumerate}
\end{corrige}

\coursec{Méthode du commentaire et production guidée}

\courssub{Méthode : Rédiger un résumé (Zusammenfassung) au Bac}

\begin{methode}[Étapes pour un résumé réussi]
\begin{enumerate}
\item \textbf{Lisez le texte deux fois} (global + détaillé). Soulignez les \cle{arguments principaux} (un par paragraphe).
\item \textbf{Notez les mots-clés} en allemand (noms, verbes, connecteurs logiques : \all{zuerst}, \all{außerdem}, \all{trotzdem}, \all{am Ende}).
\item \textbf{Rédigez en allemand}, au \textbf{présent} (temps du résumé), à la \textbf{3e personne du singulier} (\all{Der Text berichtet ...}, \all{Der Autor erklärt ...}).
\item \textbf{Structure type} :
      \begin{itemize}
      \item Introduction : \all{Der Text „Werbung überall“ von ... handelt von ... / thematisiert ...}
      \item Corps : \all{Zunächst ... / Außerdem ... / Ein weiterer Punkt ist ... / Kritisch sieht der Autor ...}
      \item Conclusion : \all{Zusammenfassend lässt sich sagen, dass ...}
      \end{itemize}
\item \textbf{Relisez-vous} : accord sujet/verbe, cas (Akk/Dat), place du verbe (principal / subordonnée), majuscules aux noms.
\end{enumerate}
\end{methode}

\courssub{Production écrite guidée : Résumé du texte}

\begin{exoresolu}[Modèle de résumé (environ 60 mots)]
\textbf{Énoncé :} Faites un résumé du texte \all{„Werbung überall“} en allemand (max. 70 mots).

\tcblower

\textbf{Proposition de correction :}

\all{Der Text „Werbung überall“ thematisiert die Omnipräsenz der Werbung und ihre Strategien. Firmen investieren Milliarden, um Produkte bekannt zu machen und Kunden zum Kauf zu bewegen. Slogans im Imperativ wie „Kauf jetzt!“ sprechen die Zielgruppe direkt an. Medien wie Radio, Handy oder Fernsehen werden gezielt genutzt. Kritiker warnen vor Manipulation und unnötigem Konsum. Dennoch informiert Werbung über Preise und Angebote. Die Frage bleibt offen: Schafft Werbung Wünsche oder spiegelt sie nur wider?}

\textbf{Analyse :} Présent de narration. Connecteurs (\all{zunächst} implicite, \all{dennoch}, \all{denn}). Structure claire. Vocabulaire du texte réemployé (\all{thematisiert}, \all{Investition}, \all{Zielgruppe}, \all{Manipulation}).
\end{exoresolu}

\courssub{Production écrite libre : Argumentation courte}

\begin{exercice}[Sujet de rédaction type Bac]
\textbf{Consigne :} \all{„Werbung manipuliert uns – oder informiert sie nur?“} Nehmen Sie Stellung. Schreiben Sie einen kurzen Text (ca. 80–100 Wörter). Nutzen Sie Vokabeln der Lektion und die Strukturen \all{um ... zu} / \all{damit}.

\textbf{Aide au plan :}
\begin{enumerate}
\item These : \all{Meiner Meinung nach ... / Ich denke, dass ...}
\item Argument 1 (Pour) : \all{Erstens ... / Zum einen ...} + \all{um ... zu} / \all{damit}
\item Argument 2 (Contre) : \all{Zweitens ... / Zum anderen ...} + \all{um ... zu} / \all{damit}
\item Conclusion : \all{Zusammenfassend ... / Abschließend ...}
\end{enumerate}
\end{exercice}

\begin{corrige}[Proposition de correction (exemple)]
\all{Meiner Meinung nach hat die Werbung zwei Seiten. Zum einen informiert sie uns über neue Produkte und Preise, damit wir gute Entscheidungen treffen können. Firmen müssen werben, um ihre Produkte zu verkaufen. Zum anderen manipuliert Werbung oft unsere Gefühle. Slogans wie „Kauf jetzt!“ drängen uns zum spontanen Kauf. Wir kaufen Dinge, die wir nicht brauchen. Zusammenfassend: Werbung ist nützlich, aber man muss kritisch bleiben und nicht jedem Slogan glauben.}
\end{corrige}

\coursec{Expression orale : Jeu de rôle et débat}

\begin{experience}[Activité orale : « Der perfekte Werbespot »]
\textbf{Objectif :} Réinvestir l'impératif, le vocabulaire thématique et les structures de finalité à l'oral.

\textbf{Déroulement (15 min) :}
\begin{enumerate}
\item \textbf{Préparation (5 min) :} Par groupes de 3. Inventez un produit imaginaire (ex. \all{„Schoko-Banane für den Morgen“}, \all{„Solar-Rucksack für Schüler“}). Définissez : \all{Produktname}, \all{Zielgruppe}, \all{Preis}, \all{besonderes Merkmal}.
\item \textbf{Production (5 min) :} Rédigez un mini-spot de 30 secondes.
      \begin{itemize}
      \item Accroche (impératif) : \all{„Probier den neuen Solar-Rucksack!“}
      \item Argument + finalité : \all{Er lädt dein Handy, damit du immer erreichbar bist.}
      \item Slogan final : \all{„Solar-Rucksack – Energie für deinen Tag!“}
      \end{itemize}
\item \textbf{Présentation (5 min) :} Chaque groupe joue son spot devant la classe. La classe vote pour le plus convaincant (\all{„Am überzeugendsten“}).
\end{enumerate}

\textbf{Grille d'évaluation (compétences Bac) :}
\begin{itemize}
\item Prononciation / Intonation (impératif énergique)
\item Correction grammaticale (impératif, \all{um ... zu} / \all{damit}, place du verbe)
\item Richesse lexicale (mots de la leçon)
\item Cohérence et persuasion
\end{itemize}
\end{experience}

\begin{experience}[Débat guidé : « Werbung in der Schule – ja oder nein? »]
\textbf{Consigne :} « Faut-il autoriser la publicité dans l'enceinte du lycée (distributeurs, affiches, sponsoring d'événements) ? »
\textbf{Méthode :} Deux camps (\all{Pro} / \all{Contra}). 5 min préparation d'arguments avec \all{damit} / \all{um ... zu}. 10 min débat modéré par le prof. Exemples d'arguments :
\begin{itemize}
\item \all{Pro:} \all{Die Schule braucht Geld, um neue Computer zu kaufen. Sponsoren geben Geld, damit ihr Logo gezeigt wird.}
\item \all{Contra:} \all{Schüler sollen lernen, nicht konsumieren. Werbung lenkt ab, damit man nicht mehr auf den Unterricht achtet.}
\end{itemize}
\end{experience}

\begin{savaistu}[Culture : La « Werbewirtschaft » en chiffres]
En 2023, les dépenses publicitaires brutes en Allemagne ont atteint environ \textbf{30 milliards d'euros}. La télévision reste le médium le plus fort (\all{ca. 4,5 Mrd. €}), suivie par l'online (\all{ca. 12 Mrd. €} incluant search, display, video, social media). Le marché camerounais, bien que plus petit, suit la même tendance : le digital (Facebook, WhatsApp, Instagram) dépasse désormais la radio et l'affichage traditionnel pour toucher la jeunesse urbaine.
\end{savaistu}

\begin{attention}[Check-list finale avant le Bac (cette leçon)]
\begin{itemize}
\item[$\square$] Je sais former l'impératif \all{du}, \all{ihr}, \all{Sie} (réguliers, irréguliers, séparables, \all{sein/haben}).
\item[$\square$] Je place la particule séparable \textbf{à la fin} à l'impératif (\all{Mach das Fenster auf!}).
\item[$\square$] Je distingue \all{um ... zu} (même sujet) et \all{damit} (sujets différents).
\item[$\square$] Je construis correctement \all{zu} + infinitif (intercalation \all{aufzumachen}).
\item[$\square$] Je place le verbe \textbf{à la fin} après \all{damit}.
\item[$\square$] Je connais le vocabulaire de base : \all{die Werbung, das Produkt, die Marke, der Kunde, die Zielgruppe, der Slogan, manipulieren, überzeugen, der Verbraucher, der Umsatz}.
\item[$\square$] Je sais écrire un résumé structuré au présent, 3e personne, avec connecteurs.
\item[$\square$] Je sais argumenter brièvement sur l'influence de la publicité.
\end{itemize}
\end{attention}

\coursec{Compréhension du texte}

Dans la section précédente, nous avons découvert le texte support \all{„Werbung überall“}, lu à voix haute, observé dans sa forme et son vocabulaire de première approche. Il s'agit maintenant de passer à l'étape décisive de toute leçon de type « texte » : la \cle{compréhension}. Au baccalauréat comme dans la vie quotidienne, lire un texte allemand ne suffit pas : il faut savoir en dégager le sens global, repérer les informations précises, justifier ses réponses par des citations et restituer l'essentiel dans un résumé structuré. Cette section vous donne la méthode complète, appliquée pas à pas à notre texte.

\begin{textebox}[Texte support : „Werbung überall“ (rappel)]
Werbung begegnet uns überall: im Fernsehen, im Radio, im Internet, auf Plakaten in der Stadt und sogar auf den Trikots der Fußballspieler. Niemand kann ihr entkommen.

Die Firmen geben jedes Jahr Milliarden aus, um ihre Produkte bekannt zu machen. Ein gutes Produkt allein reicht heute nicht mehr: Die Kunden müssen zuerst erfahren, dass es existiert. Deshalb entwickeln die Marketingabteilungen Strategien, damit die Verbraucher das Angebot bemerken und kaufen.

Ein guter Slogan ist kurz, einfach und bleibt im Kopf. Er soll positive Gefühle wecken und die Marke von der Konkurrenz unterscheiden. Denken Sie an die bekannten Werbesprüche der großen Automarken: Oft genügen drei oder vier Wörter.

Die Werbung richtet sich immer an bestimmte Zielgruppen. Für Kinder gibt es bunte Spots mit lustigen Figuren, für Jugendliche Werbung mit Musik und Sportstars, für Erwachsene seriöse Spots über Versicherungen oder Autos. Auch in Kamerun sieht man überall Reklame: an den Wänden der Geschäfte in Douala, auf den Bendskins in Yaoundé und im Radio.

Manche Leute kritisieren jedoch die Werbung. Sie sagen, sie manipulieren die Verbraucher und machen sie unglücklich, weil sie immer mehr kaufen wollen. Außerdem verschmutzen zu viele Plakate die Städte. Andere meinen dagegen, Werbung sei nützlich, denn sie informiert über neue Produkte und finanziert das Fernsehen und viele Internetseiten.

Eines ist sicher: Wer die Werbung versteht, kann besser entscheiden, was er wirklich kaufen will.
\end{textebox}

\textbf{Glossaire :} die Werbung (-) la publicité ; das Plakat (-e) l'affiche ; das Trikot (-s) le maillot ; ausgeben (gibt aus, gab aus, hat ausgegeben) dépenser ; bekannt machen faire connaître ; erfahren (erfährt) apprendre ; die Marketingabteilung (-en) le service marketing ; der Verbraucher (-) le consommateur ; der Slogan (-s) le slogan ; wecken éveiller ; die Zielgruppe (-n) le groupe cible ; der Spot (-s) le spot publicitaire ; die Reklame (-n) la réclame, la publicité ; manipulieren manipuler ; verschmutzen polluer ; nützlich utile.

\courssub{La méthode générale de compréhension}

Avant toute question, il faut organiser son travail. La compréhension d'un texte se fait toujours en cinq étapes, du plus général au plus précis.

\begin{popfigure}
\begin{tikzpicture}[node distance=0.55cm, every node/.style={draw=popblue, fill=popblueL, rounded corners, align=center, minimum width=5.2cm, minimum height=0.9cm, font=\small}, arrow/.style={-{Stealth}, thick, poporange}]
\node[align=center] (a){\textbf{1. Titre et source}\\Que promet le texte ?};
\node (b) [below=of a, align=center]{\textbf{2. Lecture globale}\\Thème, type de texte, intention};
\node (c) [below=of b, align=center]{\textbf{3. Lecture détaillée}\\Souligner les informations clés};
\node (d) [below=of c, align=center]{\textbf{4. Questions}\\Répondre par des phrases complètes};
\node (e) [below=of d, align=center]{\textbf{5. Résumé}\\Restituer l'essentiel avec connecteurs};
\draw[arrow] (a) -- (b);
\draw[arrow] (b) -- (c);
\draw[arrow] (c) -- (d);
\draw[arrow] (d) -- (e);
\end{tikzpicture}
\legende{Organigramme de la méthode de compréhension en cinq étapes}
\end{popfigure}

\begin{propriete}[L'ordre des questions au Bac]
Au baccalauréat, les questions de compréhension suivent presque toujours \cle{l'ordre du texte}. La réponse à la question 2 se trouve donc \emph{après} la réponse à la question 1, et la réponse à la question 3 après celle de la question 2. Ce repérage vous fait gagner un temps précieux : inutile de relire tout le texte pour chaque question, avancez ligne après ligne.
\end{propriete}

\begin{methode}[Comment répondre à une question de compréhension]
\begin{enumerate}
\item \textbf{Repérer} dans la question les mots-clés (noms, verbes) et les retrouver dans le texte, souvent sous forme de synonymes.
\item \textbf{Reprendre les mots de la question} dans la réponse : à \all{Wo sieht man Werbung?} on répond \all{Man sieht Werbung im Fernsehen, im Radio...} et non un mot isolé.
\item \textbf{Répondre par une phrase complète} : sujet + verbe conjugué en deuxième position + compléments.
\item \textbf{Citer le texte} entre guillemets si la consigne le demande (\all{Zitieren Sie}, \all{mit einer Textstelle belegen}).
\item \textbf{Ne jamais recopier tout un paragraphe} : une citation courte et ciblée suffit ; recopier cinq lignes montre qu'on n'a pas compris.
\end{enumerate}
\end{methode}

\courssub{Compréhension globale : thème, type de texte, intention}

La lecture globale répond à trois questions fondamentales, à se poser avant toute analyse détaillée.

\begin{exemplebox}[Questions de compréhension globale et réponses modèles]
\textbf{1. Worum geht es in diesem Text?} (De quoi parle ce texte ?)\par
\all{Der Text handelt von der Werbung: wo man sie sieht, wie sie funktioniert und warum manche Leute sie kritisieren.} « Le texte parle de la publicité : où on la voit, comment elle fonctionne et pourquoi certaines personnes la critiquent. »\par\medskip
\textbf{2. Um welche Textsorte handelt es sich?} (Quel est le type de texte ?)\par
\all{Es ist ein informativer Sachtext, ein Zeitungsartikel über Marketing und Werbung.} « C'est un texte informatif, un article de presse sur le marketing et la publicité. »\par\medskip
\textbf{3. Was will der Autor erreichen?} (Quelle est l'intention de l'auteur ?)\par
\all{Der Autor will die Leser informieren und sie dazu bringen, kritisch über Werbung nachzudenken.} « L'auteur veut informer les lecteurs et les amener à réfléchir de manière critique sur la publicité. »
\end{exemplebox}

\begin{definition}[Les trois questions de la lecture globale]
\all{Worum geht es?} = le thème ; \all{Welche Textsorte?} = le type de texte (article, interview, récit, publicité, lettre) ; \all{Was will der Autor?} = l'intention (informer, convaincre, critiquer, amuser). Ces trois réponses tiennent chacune en une phrase et constituent souvent les premières questions de l'épreuve.
\end{definition}

\courssub{Compréhension détaillée : les cinq questions clés}

Passons maintenant à la lecture détaillée, paragraphe par paragraphe, en suivant l'ordre du texte.

\begin{exoresolu}[Questions détaillées sur „Werbung überall“]
\textbf{Énoncé.} Répondez en allemand, par des phrases complètes.\par
1. Wo sieht man Werbung?\par
2. Wie viel Geld geben die Firmen für Werbung aus?\par
3. Was ist ein guter Slogan?\par
4. Welche Zielgruppen nennt der Text?\par
5. Was kritisieren manche Leute an der Werbung?
\tcblower
\textbf{Solution détaillée.}\par
1. \all{Man sieht Werbung im Fernsehen, im Radio, im Internet, auf Plakaten und auf den Trikots der Fußballspieler.} « On voit la publicité à la télévision, à la radio, sur Internet, sur les affiches et sur les maillots des footballeurs. » (paragraphe 1)\par
2. \all{Die Firmen geben jedes Jahr Milliarden aus, um ihre Produkte bekannt zu machen.} « Les entreprises dépensent chaque année des milliards pour faire connaître leurs produits. » (paragraphe 2)\par
3. \all{Ein guter Slogan ist kurz, einfach und bleibt im Kopf.} « Un bon slogan est court, simple et reste en tête. » (paragraphe 3)\par
4. \all{Der Text nennt drei Zielgruppen: Kinder, Jugendliche und Erwachsene.} « Le texte cite trois groupes cibles : les enfants, les adolescents et les adultes. » (paragraphe 4)\par
5. \all{Manche Leute kritisieren, dass die Werbung die Verbraucher manipuliert und dass zu viele Plakate die Städte verschmutzen.} « Certains reprochent à la publicité de manipuler les consommateurs et aux affiches trop nombreuses de polluer les villes. » (paragraphe 5)
\end{exoresolu}

\begin{exemplebox}[Exemple traduit : la construction um ... zu]
\all{Firmen geben Milliarden aus, um ihre Produkte bekannt zu machen.} = « Les entreprises dépensent des milliards pour faire connaître leurs produits. »\par
Observez : \all{um ... zu} exprime la finalité (« pour »), l'infinitif \all{zu machen} se place à la fin, et le verbe à particule \all{bekannt machen} se réunit à l'infinitif : \all{bekannt zu machen}.
\end{exemplebox}

\begin{attention}[Ne traduisez pas mot à mot dans vos réponses]
Le piège classique du francophone : \all{bekannt machen} signifie « faire connaître, rendre célèbre », et non « faire connu ». De même, \all{im Kopf bleiben} = « rester en tête » (et non « rester dans la tête »), \all{sich richten an} = « s'adresser à ». Une réponse au Bac doit être en allemand idiomatique, pas en français déguisé. Autre piège : \all{die Reklame} est un quasi-synonyme de \all{die Werbung}, mais le français « réclame » a un autre sens — faux ami à surveiller.
\end{attention}

\courssub{Vrai ou faux : justifier par une citation}

L'exercice \all{Richtig oder Falsch?} est un classique du baccalauréat. Règle d'or : une réponse sans justification ne rapporte que la moitié des points. Il faut citer la phrase du texte entre guillemets et indiquer le paragraphe ou la ligne.

\begin{exoresolu}[Richtig oder Falsch? Begründen Sie mit einer Textstelle]
\textbf{Énoncé.} Dites si les affirmations sont vraies ou fausses et justifiez par une citation.\par
1. Man kann der Werbung leicht entkommen.\par
2. Ein gutes Produkt allein reicht heute nicht mehr.\par
3. Ein Slogan soll lang und kompliziert sein.\par
4. In Kamerun gibt es keine Werbung.\par
5. Werbung finanziert das Fernsehen.
\tcblower
\textbf{Solution détaillée.}\par
1. \all{Falsch. Der Text sagt: „Niemand kann ihr entkommen.“} (paragraphe 1)\par
2. \all{Richtig. „Ein gutes Produkt allein reicht heute nicht mehr.“} (paragraphe 2)\par
3. \all{Falsch. „Ein guter Slogan ist kurz, einfach und bleibt im Kopf.“} (paragraphe 3)\par
4. \all{Falsch. „Auch in Kamerun sieht man überall Reklame: an den Wänden der Geschäfte in Douala, auf den Bendskins in Yaoundé.“} (paragraphe 4)\par
5. \all{Richtig. „...sie informiert über neue Produkte und finanziert das Fernsehen und viele Internetseiten.“} (paragraphe 5)
\end{exoresolu}

\begin{aretenir}[La formule de justification]
Pour justifier, utilisez le schéma : \all{Richtig/Falsch} + \all{Der Text sagt: „...“} ou \all{Im Text steht: „...“} + indication du paragraphe. La citation doit être \emph{mot à mot}, entre guillemets allemands „...“, et suffisamment courte.
\end{aretenir}

\courssub{Recherche lexicale dans le texte}

L'examinateur demande souvent de retrouver dans le texte l'équivalent d'un mot français, un synonyme ou un antonyme. C'est un exercice de repérage rapide : le mot cherché se trouve dans la zone du texte indiquée par l'ordre des questions.

\begin{center}
\begin{tabularx}{\linewidth}{|X|X|X|}
\hline
\rowcolor{popblueL} \textbf{Français} & \textbf{Équivalent dans le texte} & \textbf{Remarque} \\
\hline
dépenser (de l'argent) & \all{Geld ausgeben} & verbe à particule : \all{gibt ... aus} \\
\hline
faire connaître & \all{bekannt machen} & pas « faire connu » \\
\hline
le consommateur & \all{der Verbraucher} & famille : \all{verbrauchen} \\
\hline
le groupe cible & \all{die Zielgruppe} & \all{das Ziel} = le but \\
\hline
éveiller (un sentiment) & \all{wecken} & \all{Gefühle wecken} \\
\hline
polluer, salir & \all{verschmutzen} & préfixe \all{ver-} \\
\hline
utile & \all{nützlich} & contraire : \all{unnütz} \\
\hline
\end{tabularx}
\end{center}

\begin{exemplebox}[Synonymes et antonymes dans le texte]
\textbf{Synonymes :} \all{die Werbung} $\approx$ \all{die Reklame} ; \all{der Kunde} $\approx$ \all{der Verbraucher} ; \all{der Slogan} $\approx$ \all{der Werbespruch}.\par
\textbf{Antonymes :} \all{kurz} $\leftrightarrow$ \all{lang} ; \all{einfach} $\leftrightarrow$ \all{kompliziert} ; \all{nützlich} $\leftrightarrow$ \all{schädlich} ; \all{positiv} $\leftrightarrow$ \all{negativ}.
\end{exemplebox}

\courssub{Reformulation et résumé guidé}

La reformulation est l'exercice le plus noble : elle prouve que vous avez réellement compris. Consigne type : \all{Erklären Sie den letzten Satz mit eigenen Worten.}

\begin{exemplebox}[Reformulation de la dernière phrase]
Phrase du texte : \all{„Wer die Werbung versteht, kann besser entscheiden, was er wirklich kaufen will.“}\par
Reformulation possible : \all{Wenn man die Tricks der Werbung kennt, wird man nicht manipuliert und kauft nur die Produkte, die man wirklich braucht.} « Quand on connaît les techniques de la publicité, on ne se laisse pas manipuler et on n'achète que les produits dont on a vraiment besoin. »\par
Technique : remplacer les mots par des synonymes (\all{verstehen} $\rightarrow$ \all{kennen}), transformer la structure (proposition relative $\rightarrow$ phrase avec \all{wenn}), garder le sens.
\end{exemplebox}

\begin{methode}[Rédiger un résumé guidé en 3-4 phrases]
\begin{enumerate}
\item Identifier les quatre idées principales, une par paragraphe ou groupe de paragraphes.
\item Rédiger une phrase par idée, avec ses propres mots, sans citation.
\item Enchaîner avec les connecteurs imposés : \all{zuerst} (d'abord), \all{dann} (ensuite), \all{außerdem} (de plus), \all{schließlich} (finalement).
\item Vérifier : verbe en deuxième position, connecteur en position 1 suivi du verbe, pas de détail superflu.
\end{enumerate}
\end{methode}

\begin{exoresolu}[Zusammenfassung : résumé avec connecteurs imposés]
\textbf{Énoncé.} Résumez le texte en quatre phrases en utilisant : \all{zuerst}, \all{dann}, \all{außerdem}, \all{schließlich}.
\tcblower
\textbf{Solution détaillée.}\par
\all{Zuerst erklärt der Text, dass Werbung überall ist und dass die Firmen viel Geld dafür ausgeben.} « D'abord, le texte explique que la publicité est partout et que les entreprises y dépensent beaucoup d'argent. »\par
\all{Dann beschreibt er, wie ein guter Slogan funktioniert und an welche Zielgruppen sich die Werbung richtet.} « Ensuite, il décrit comment fonctionne un bon slogan et à quels groupes cibles la publicité s'adresse. »\par
\all{Außerdem nennt der Autor die Kritik an der Werbung: Sie manipuliert die Verbraucher und verschmutzt die Städte.} « De plus, l'auteur mentionne les critiques : elle manipule les consommateurs et pollue les villes. »\par
\all{Schließlich sagt er, dass man besser entscheiden kann, wenn man die Werbung versteht.} « Finalement, il affirme qu'on peut mieux décider quand on comprend la publicité. »
\end{exoresolu}

\begin{attention}[Connecteur en tête de phrase : inversion obligatoire]
Après \all{zuerst}, \all{dann}, \all{außerdem}, \all{schließlich} placés en position 1, le verbe conjugué vient \emph{immédiatement}, puis le sujet : \all{Dann beschreibt \textbf{er}...} et jamais \all{Dann er beschreibt...}. C'est la règle d'or de l'ordre des mots allemands : le verbe conjugué occupe toujours la deuxième place de la phrase énonciative.
\end{attention}

\begin{exercice}[À vous de jouer]
1. Répondez par des phrases complètes : \all{Warum entwickeln die Firmen Marketingstrategien? Welche positiven Seiten der Werbung nennt der Text?}\par
2. \all{Richtig oder Falsch?} Justifiez : \all{„Werbung für Kinder benutzt lustige Figuren.“}\par
3. Trouvez dans le texte l'équivalent de : « rester en tête », « s'adresser à », « cependant ».\par
4. Reformulez avec vos propres mots : \all{„Ein gutes Produkt allein reicht heute nicht mehr.“}
\end{exercice}

\begin{corrige}[Exercice « À vous de jouer »]
1. \all{Die Firmen entwickeln Strategien, damit die Verbraucher das Angebot bemerken und kaufen.} — \all{Der Text nennt zwei positive Seiten: Die Werbung informiert über neue Produkte und finanziert das Fernsehen und viele Internetseiten.}\par
2. \all{Richtig. Der Text sagt: „Für Kinder gibt es bunte Spots mit lustigen Figuren.“} (paragraphe 4)\par
3. « rester en tête » = \all{im Kopf bleiben} ; « s'adresser à » = \all{sich richten an} ; « cependant » = \all{jedoch}.\par
4. \all{Heute muss man auch gute Produkte bekannt machen, sonst kaufen die Kunden sie nicht.} « Aujourd'hui, il faut aussi faire connaître les bons produits, sinon les clients ne les achètent pas. »
\end{corrige}

\begin{aretenir}[L'essentiel de la compréhension de texte]
Lire deux fois (globale puis détaillée), suivre l'ordre du texte, répondre par des phrases complètes en reprenant les mots de la question, justifier par des citations courtes entre guillemets, reformuler avec des synonymes et résumer avec les connecteurs \all{zuerst}, \all{dann}, \all{außerdem}, \all{schließlich} suivis de l'inversion. Ces automatismes, acquis sur \all{„Werbung überall“}, vous serviront pour tout texte du baccalauréat.
\end{aretenir}

\coursec{Lexique thématique : Marketing und Werbung}

Dans la section précédente, nous avons lu et compris le texte \all{„Werbung überall“}. Nous y avons rencontré de nombreux mots nouveaux : \all{der Kunde}, \all{das Produkt}, \all{die Werbung}, \all{das Angebot}... Le moment est venu de les organiser, de les mémoriser avec leurs \cle{articles} et leurs \cle{pluriels}, et d'apprendre à les employer dans des phrases complètes. Un bon lexique, c'est la moitié de la réussite à l'épreuve d'allemand : sans vocabulaire précis, impossible de résumer un texte économique ni de rédiger un paragraphe sur la publicité.

\courssub{La technique des couleurs : apprendre l'article avec le nom}

\begin{propriete}[Règle d'or du vocabulaire allemand]
En allemand, un nom ne s'apprend \textbf{jamais seul}. On mémorise toujours ensemble : l'\cle{article} (der / die / das), le \cle{nom} avec sa majuscule, et le \cle{pluriel}. Exemple : \all{der Kunde, -n} se lit « der Kunde, pluriel die Kunden ».
\end{propriete}

\begin{methode}[La méthode des trois couleurs]
Pour mémoriser le genre des noms, utilisez un code couleur dans vos cahiers et vos fiches :
\begin{enumerate}
\item \textbf{Bleu} pour \all{der} (masculin) : écrivez \all{der Kunde} en bleu.
\item \textbf{Rouge} pour \all{die} (féminin) : écrivez \all{die Marke} en rouge.
\item \textbf{Vert} pour \all{das} (neutre) : écrivez \all{das Produkt} en vert.
\end{enumerate}
Au bout de quelques semaines, votre mémoire visuelle associera automatiquement la couleur au genre, et vous ne direz plus jamais « le Marke » !
\end{methode}

\courssub{Champ lexical 1 : Les acteurs (die Akteure)}

Observons d'abord ces phrases tirées de situations réelles :

\all{Der Kunde kauft ein Produkt.} « Le client achète un produit. »

\all{Die Firma wirbt um neue Kunden.} « L'entreprise fait de la publicité pour attirer de nouveaux clients. »

\begin{center}
\begin{tabularx}{\linewidth}{|X|X|X|}\hline
\rowcolor{popblueL} Deutsch (Singular / Plural) & Français & Beispiel \\ \hline
\all{der Kunde, -n} & le client & \all{Der Kunde ist König.} (Le client est roi.) \\ \hline
\all{die Kundin, -nen} & la cliente & \all{Die Kundin zahlt bar.} \\ \hline
\all{der Verbraucher, -} & le consommateur & \all{Die Verbraucher wollen Qualität.} \\ \hline
\all{die Firma, Firmen} & l'entreprise, la société & \all{Die Firma hat 200 Mitarbeiter.} \\ \hline
\all{der Hersteller, -} & le fabricant, le producteur & \all{Der Hersteller gibt zwei Jahre Garantie.} \\ \hline
\all{die Zielgruppe, -n} & le groupe cible & \all{Die Zielgruppe sind Jugendliche.} \\ \hline
\end{tabularx}
\end{center}

\begin{definition}[die Zielgruppe, -n]
\all{Die Zielgruppe} désigne le groupe de personnes auquel une publicité ou un produit s'adresse : les jeunes, les familles, les seniors, les sportifs... Une publicité pour des bonbons visera les enfants ; une publicité pour une voiture familiale visera les parents.
\all{Diese Werbung spricht eine junge Zielgruppe an.} « Cette publicité s'adresse à un public jeune. »
\end{definition}

\begin{attention}[der Kunde : un nom faible (n-Deklination)]
\all{Der Kunde} appartient à la \cle{déclinaison faible} : il prend un \textbf{-n} à tous les cas sauf au nominatif singulier : \all{der Kunde}, mais \all{den Kunden} (accusatif), \all{dem Kunden} (datif), \all{des Kunden} (génitif).
\all{Die Verkäuferin fragt den Kunden.} « La vendeuse interroge le client. » (et non \all{den Kunde} !)
Autre piège : ne confondez pas \all{der Kunde} (le client) avec le français « la connaissance ». Le féminin est \all{die Kundin, -nen}. De même : \all{der Verkäufer} (le vendeur) / \all{die Verkäuferin} (la vendeuse). Le suffixe \all{-in} forme le féminin des métiers et des rôles.
\end{attention}

\courssub{Champ lexical 2 : Le produit et la marque (Produkt und Marke)}

\begin{center}
\begin{tabularx}{\linewidth}{|X|X|X|}\hline
\rowcolor{popblueL} Deutsch & Français & Beispiel \\ \hline
\all{das Produkt, -e} & le produit & \all{Das Produkt ist neu auf dem Markt.} \\ \hline
\all{die Marke, -n} & la marque & \all{Diese Marke ist weltbekannt.} \\ \hline
\all{das Angebot, -e} & l'offre & \all{Das Angebot gilt nur bis Sonntag.} \\ \hline
\all{der Preis, -e} & le prix & \all{Der Preis ist zu hoch.} \\ \hline
\all{die Qualität, -en} & la qualité & \all{Gute Qualität hat ihren Preis.} \\ \hline
\all{der Markt, die Märkte} & le marché & \all{Der Markt in Douala ist riesig.} \\ \hline
\all{die Ware, -n} & la marchandise & \all{Die Ware ist frisch.} \\ \hline
\end{tabularx}
\end{center}

\begin{exemplebox}[Exemples traduits]
\all{Die Firma bringt ein neues Produkt auf den Markt.} « L'entreprise lance un nouveau produit sur le marché. »

\all{Diese Woche ist der Fernseher im Angebot.} « Cette semaine, la télévision est en promotion. »

\all{Die Qualität stimmt, aber der Preis ist zu hoch.} « La qualité est bonne, mais le prix est trop élevé. »
\end{exemplebox}

\begin{attention}[das Angebot : deux sens à distinguer]
\all{Das Angebot} signifie d'abord « l'offre » en général : \all{ein Angebot machen} (faire une offre). Mais dans le langage commercial, \all{im Angebot sein} signifie « être en promotion, en solde » : \all{Die Schuhe sind im Angebot.} « Les chaussures sont en promotion. » Ne traduisez donc jamais \all{im Angebot} par « dans l'offre » !
\end{attention}

\courssub{Champ lexical 3 : La publicité (die Werbung)}

\begin{center}
\begin{tabularx}{\linewidth}{|X|X|X|}\hline
\rowcolor{popblueL} Deutsch & Français & Beispiel \\ \hline
\all{die Werbung, -en} & la publicité & \all{Die Werbung beeinflusst die Kunden.} \\ \hline
\all{der Werbespot, -s} & le spot publicitaire & \all{Der Werbespot dauert 30 Sekunden.} \\ \hline
\all{das Plakat, -e} & l'affiche & \all{Das Plakat hängt an der Wand.} \\ \hline
\all{der Slogan, -s} & le slogan & \all{Der Slogan bleibt im Kopf.} \\ \hline
\all{die Anzeige, -n} & l'annonce (presse) & \all{Ich lese die Anzeigen in der Zeitung.} \\ \hline
\all{die Reklame, -n} & la publicité, la réclame & \all{Reklame machen} = faire de la publicité \\ \hline
\all{die Werbekampagne, -n} & la campagne publicitaire & \all{Die Kampagne startet im Mai.} \\ \hline
\end{tabularx}
\end{center}

\begin{attention}[die Reklame ou die Werbung ?]
\all{Die Reklame} et \all{die Werbung} sont synonymes au sens de « publicité » : \all{Reklame machen} = \all{Werbung machen} (faire de la publicité). Mais attention : le verbe \all{werben} a aussi un sens plus ancien : \all{um eine Frau werben} signifie « faire la cour à une femme » ! Le contexte est donc obligatoire pour comprendre. Dans un texte économique, \all{die Werbung} = la publicité ; dans un roman d'amour, méfiez-vous !
\end{attention}

\courssub{Champ lexical 4 : Vendre et acheter (Verkaufen und Kaufen)}

\begin{center}
\begin{tabularx}{\linewidth}{|X|X|X|}\hline
\rowcolor{popblueL} Deutsch & Français & Beispiel \\ \hline
\all{verkaufen} & vendre & \all{Der Händler verkauft Obst.} \\ \hline
\all{kaufen} & acheter & \all{Ich kaufe ein neues Handy.} \\ \hline
\all{anbieten (séparable)} & proposer, offrir & \all{Das Geschäft bietet Rabatte an.} \\ \hline
\all{bestellen} & commander & \all{Wir bestellen online.} \\ \hline
\all{Geld ausgeben (séparable)} & dépenser de l'argent & \all{Jugendliche geben viel Geld aus.} \\ \hline
\all{sparen} & économiser & \all{Ich spare für ein Fahrrad.} \\ \hline
\all{der Verkauf, die Verkäufe} & la vente & \all{Der Verkauf steigt.} \\ \hline
\all{der Rabatt, -e} & la remise, le rabais & \all{Es gibt 20 \% Rabatt.} \\ \hline
\end{tabularx}
\end{center}

\begin{attention}[Verbes à particule séparable]
\all{anbieten} et \all{ausgeben} sont des verbes à \cle{particule séparable} : au Präsens, la particule va à la fin de la phrase. \all{Das Geschäft \textbf{bietet} heute günstige Preise \textbf{an}.} « Le magasin propose aujourd'hui des prix avantageux. » Erreur typique des francophones : \all{Das Geschäft anbietet...} — interdit !
\end{attention}

\courssub{Les familles de mots (Wortfamilien)}

L'allemand construit des familles entières de mots autour d'une même racine. Apprendre la famille, c'est apprendre trois ou quatre mots pour le prix d'un.

\begin{center}
\begin{tabularx}{\linewidth}{|X|X|X|}\hline
\rowcolor{popblueL} Verbe & Nom d'action & Personne / composé \\ \hline
\all{werben} (faire de la pub) & \all{die Werbung} (la publicité) & \all{der Werbespot} (le spot) \\ \hline
\all{verkaufen} (vendre) & \all{der Verkauf} (la vente) & \all{der Verkäufer} (le vendeur) \\ \hline
\all{kaufen} (acheter) & \all{der Kauf} (l'achat) & \all{der Käufer} (l'acheteur) \\ \hline
\all{der Kunde} (le client) & \all{der Kundendienst} (le SAV) & \all{die Kundin} (la cliente) \\ \hline
\all{anbieten} (proposer) & \all{das Angebot} (l'offre) & \all{der Anbieter} (le fournisseur) \\ \hline
\end{tabularx}
\end{center}

\begin{savaistu}[Le Kundendienst, un mot composé typique]
L'allemand économique adore les mots composés : \all{der Kundendienst} (le service après-vente) = \all{Kunde} (client) + \all{Dienst} (service). On trouve aussi \all{die Kundenzufriedenheit} (la satisfaction du client) ou \all{das Werbeplakat} (l'affiche publicitaire). Astuce : le genre d'un mot composé est toujours celui du \textbf{dernier} élément : \all{der Dienst} donne donc \all{der Kundendienst}.
\end{savaistu}

\courssub{Expressions utiles à réutiliser}

\begin{itemize}
\item \all{ein Produkt auf den Markt bringen} : lancer un produit sur le marché.
\item \all{eine Werbekampagne starten} : lancer une campagne publicitaire.
\item \all{Kunden gewinnen} : gagner / attirer des clients (\all{gewinnen} est un verbe fort : gewinnt, gewann, hat gewonnen).
\item \all{im Angebot sein} : être en promotion.
\item \all{die Preise senken / erhöhen} : baisser / augmenter les prix.
\item \all{Werbung machen für ein Produkt} : faire de la publicité pour un produit.
\item \all{mit der Qualität zufrieden sein} : être satisfait de la qualité.
\end{itemize}

\begin{popfigure}
\begin{center}
\begin{tikzpicture}[mindmap, grow cyclic,
  every node/.style={concept, align=center, font=\small},
  root concept/.append style={concept color=popblue, font=\small\bfseries},
  level 1/.append style={level distance=3.6cm, sibling angle=120},
  level 2/.append style={level distance=2.4cm, sibling angle=45, font=\scriptsize}]
  \node[root concept]{Marketing und Werbung}
    child[concept color=poporange] { node {Akteure}
      child { node {der Kunde, -n} }
      child { node {der Verbraucher, -} }
      child { node {die Firma, Firmen} }
      child { node {der Hersteller, -} }
      child { node {die Zielgruppe, -n} } }
    child[concept color=popgreen] { node {Produkt \& Marke}
      child { node {das Produkt, -e} }
      child { node {die Marke, -n} }
      child { node {der Preis, -e} }
      child { node {die Qualität, -en} }
      child { node {das Angebot, -e} } }
    child[concept color=poppurple] { node {Werbung \& Verkauf}
      child { node {die Werbung, -en} }
      child { node {der Werbespot, -s} }
      child { node {das Plakat, -e} }
      child { node {verkaufen / kaufen} }
      child { node {der Rabatt, -e} } };
\end{tikzpicture}
\end{center}
\legende{Carte mentale du lexique : trois branches pour organiser sa mémorisation.}
\end{popfigure}

\courssub{Texte d'application : mini-lecture lexicale}

\begin{textebox}[„Eine Firma aus Yaoundé startet eine Kampagne“]
Die Firma \emph{KamerunKakao} ist ein kleiner Hersteller von Schokolade in Yaoundé. Sie bringt ein neues Produkt auf den Markt: eine Schokolade mit 70 \% Kakao. Um Kunden zu gewinnen, startet die Firma eine große Werbekampagne. Es gibt bunte Plakate in der Stadt, Anzeigen in der Zeitung und einen Werbespot im Fernsehen. Der Slogan lautet: „Unser Kakao — unsere Stärke!“ Die Zielgruppe sind junge Familien und Schokoladenliebhaber.

In der ersten Woche ist die neue Schokolade im Angebot: Die Kunden bekommen 15 \% Rabatt. Viele Verbraucher kaufen das Produkt und sind mit der Qualität zufrieden. „Der Preis ist gut, und die Schokolade schmeckt fantastisch“, sagt eine Kundin. Der Verkauf steigt schnell. Auch der Kundendienst der Firma arbeitet gut: Er beantwortet alle Fragen der Kunden freundlich und schnell.

Der Chef der Firma erklärt: „Werbung ist wichtig, aber die Qualität ist wichtiger. Nur ein gutes Produkt macht die Kunden glücklich — und glückliche Kunden kommen wieder.“
\end{textebox}

\textbf{Glossaire} : \all{der Hersteller, -} (le fabricant) ; \all{auf den Markt bringen} (lancer sur le marché) ; \all{der Schokoladenliebhaber, -} (l'amateur de chocolat) ; \all{der Rabatt, -e} (la remise) ; \all{zufrieden} (satisfait) ; \all{steigen} (augmenter) ; \all{freundlich} (aimable) ; \all{lauten} (être formulé, se lire).

\textbf{Questions de compréhension} :
\begin{enumerate}
\item Welches Produkt bringt die Firma auf den Markt ?
\item Wer ist die Zielgruppe der Kampagne ?
\item Was bedeutet hier „im Angebot“ ?
\end{enumerate}

\begin{corrige}[Questions de compréhension]
1. \all{Die Firma bringt eine Schokolade mit 70 \% Kakao auf den Markt.} (L'entreprise lance un chocolat à 70 \% de cacao.)

2. \all{Die Zielgruppe sind junge Familien und Schokoladenliebhaber.} (Le groupe cible : les jeunes familles et les amateurs de chocolat.)

3. \all{„Im Angebot“ bedeutet hier: Die Schokolade ist billiger, die Kunden bekommen 15 \% Rabatt.} (Ici, « en promotion » : le produit est moins cher, avec 15 \% de remise.)
\end{corrige}

\courssub{Exercice résolu et entraînement}

\begin{exoresolu}[Complétez avec le bon article et le bon mot]
Énoncé — Complétez avec : \all{der Kunde}, \all{die Marke}, \all{das Angebot}, \all{die Werbung}, \all{der Rabatt}.

1. \all{... ist sehr aggressiv: Man sieht sie überall.} \quad 2. \all{... kauft drei Tafeln Schokolade.} \quad 3. \all{... dieser Woche ist interessant: Alles ist billiger.} \quad 4. \all{... ist weltbekannt.} \quad 5. \all{Wir bekommen 10 \% ... .}
\tcblower
Solution détaillée :
\begin{enumerate}
\item \all{Die Werbung} — nominatif féminin (sujet ; l'indice est le pronom \all{sie} : « on la voit partout »).
\item \all{Der Kunde} — nominatif masculin (sujet du verbe \all{kauft}).
\item \all{Das Angebot} — nominatif neutre (sujet, « l'offre de cette semaine »).
\item \all{Die Marke} — nominatif féminin.
\item \all{Rabatt} — après un pourcentage, on emploie le nom sans article : \all{10 \% Rabatt bekommen} (obtenir 10 \% de remise).
\end{enumerate}
Méthode : identifiez d'abord la fonction (sujet = nominatif), puis le genre du nom grâce à vos fiches couleurs.
\end{exoresolu}

\begin{exercice}[À vous de jouer]
1. Traduisez en allemand : a) « L'entreprise lance un nouveau produit sur le marché. » b) « Cette semaine, le téléphone portable est en promotion. » c) « La publicité s'adresse aux jeunes. »

2. Formez la famille de mots : \all{verkaufen} → nom d'action → personne.

3. Recopiez la carte mentale et ajoutez deux mots nouveaux à chaque branche.

4. Mettez le nom entre parenthèses au cas correct : \all{Die Verkäuferin fragt (der Kunde).} / \all{Wir danken (der Kunde) für sein Vertrauen.}
\end{exercice}

\begin{corrige}[Exercice « À vous de jouer »]
1. a) \all{Die Firma bringt ein neues Produkt auf den Markt.} b) \all{Diese Woche ist das Handy im Angebot.} c) \all{Die Werbung spricht Jugendliche an.} (attention : \all{ansprechen} est séparable et régit l'accusatif).

2. \all{verkaufen} → \all{der Verkauf} → \all{der Verkäufer / die Verkäuferin}.

3. Exemples possibles : Akteure : \all{der Händler}, \all{der Chef} ; Produkt \& Marke : \all{der Markt}, \all{die Ware} ; Werbung \& Verkauf : \all{die Anzeige}, \all{der Slogan}.

4. \all{Die Verkäuferin fragt den Kunden.} (accusatif, déclinaison faible) ; \all{Wir danken dem Kunden für sein Vertrauen.} (datif après \all{danken}, déclinaison faible).
\end{corrige}

\begin{aretenir}[L'essentiel du lexique]
Retenez les six mots-clés du programme avec article et pluriel : \all{die Werbung (-en)}, \all{das Produkt (-e)}, \all{der Kunde (-n)}, \all{das Angebot (-e)}, \all{der Verkauf (Verkäufe)}, \all{die Marke (-n)}. Apprenez-les par familles, avec la méthode des couleurs, et réutilisez les expressions \all{auf den Markt bringen}, \all{Kunden gewinnen} et \all{im Angebot sein} dans vos productions écrites. N'oubliez pas : \all{der Kunde} est un nom faible (\all{den Kunden}, \all{dem Kunden}).
\end{aretenir}

\coursec{Grammaire 1 : L'impératif (der Imperativ)}

Après avoir enrichi notre vocabulaire du marketing et de la publicité, nous allons maintenant observer de plus près la grammaire des slogans que nous venons de lire. En effet, la publicité possède une forme verbale favorite : l'\cle{impératif}. C'est le mode de l'ordre, du conseil, de l'invitation — exactement ce que fait une publicité quand elle nous dit : « Achète ! », « Essaie ! », « Découvre ! ».

\courssub{Observation : les slogans parlent à l'impératif}

\begin{textebox}[Slogans et consignes publicitaires]
„Probier es einfach!“ — „Kauft jetzt und spart Geld!“ — „Besuchen Sie uns in unserem neuen Geschäft!“ — „Lies das Angebot genau!“ — „Vergleicht die Preise!“ — „Sei clever und wähle unsere Marke!“ — „Schau dir unsere Produkte an!“ — „Seien Sie willkommen bei uns!“

\emph{Glossaire :} \all{einfach} = simplement ; \all{sparen} = économiser ; \all{das Geschäft, -e} = le magasin ; \all{genau} = attentivement, exactement ; \all{clever} = malin ; \all{wählen} = choisir ; \all{das Angebot, -e} = l'offre ; \all{willkommen} = bienvenu(e).
\end{textebox}

Observons ces phrases. Trois faits frappent immédiatement :

\begin{enumerate}
\item Le \cle{verbe} se trouve en \cle{première position} de la phrase, suivi d'un point d'exclamation.
\item Il n'y a \cle{pas de sujet exprimé} devant le verbe — sauf dans \all{Besuchen Sie uns!}, où le pronom \all{Sie} apparaît \emph{après} le verbe.
\item Les formes verbales changent selon la personne à qui l'on s'adresse : \all{Probier!} (une personne, tutoiement), \all{Kauft!} (plusieurs personnes), \all{Besuchen Sie!} (vouvoiement).
\end{enumerate}

C'est exactement la logique de l'impératif allemand : il possède \cle{trois formes}, correspondant aux trois façons de s'adresser à quelqu'un : \all{du}, \all{ihr} et \all{Sie}.

\begin{definition}[L'impératif]
L'\cle{impératif} (\all{der Imperativ}) est le mode verbal qui sert à donner un \cle{ordre}, un \cle{conseil}, une \cle{invitation}, une \cle{consigne} ou une \cle{recommandation}. En allemand, il existe à trois personnes : \all{du} (2\up{e} personne du singulier), \all{ihr} (2\up{e} personne du pluriel) et \all{Sie} (forme de politesse, singulier et pluriel). Le verbe conjugué occupe la \cle{première place} de la phrase.
\end{definition}

\begin{exemplebox}[Emplois de l'impératif]
\begin{itemize}
\item \textbf{Ordre :} \all{Komm sofort her!} « Viens tout de suite ! »
\item \textbf{Conseil :} \all{Vergleich die Preise vor dem Kauf!} « Compare les prix avant l'achat ! »
\item \textbf{Invitation :} \all{Besuchen Sie unsere Ausstellung!} « Visitez notre exposition ! »
\item \textbf{Consigne (mode d'emploi) :} \all{Öffnen Sie die Packung vorsichtig!} « Ouvrez l'emballage avec précaution ! »
\item \textbf{Recommandation publicitaire :} \all{Probier unser neues Produkt!} « Essaie notre nouveau produit ! »
\end{itemize}
\end{exemplebox}

\courssub{Formation de l'impératif à la 2\up{e} personne du singulier (du)}

C'est la forme la plus fréquente dans la publicité, car la publicité s'adresse au consommateur individuellement, comme à un ami.

\begin{propriete}[Impératif « du » : la règle de base]
À la 2\up{e} personne du singulier, l'impératif se forme avec le \cle{radical} du présent de l'indicatif, \cle{sans pronom}, auquel on peut ajouter un \cle{-e facultatif} (surtout à l'écrit soutenu) :
\begin{center}
\all{komm(en) $\rightarrow$ Komm!} \quad \all{mach(en) $\rightarrow$ Mach! / Mache!} \quad \all{probier(en) $\rightarrow$ Probier! / Probiere!}
\end{center}
Le pronom \all{du} n'apparaît \cle{jamais}.
\end{propriete}

\begin{exemplebox}[Exemples traduits]
\all{Kauf das Produkt!} « Achète le produit ! »

\all{Mach die Tür zu!} « Ferme la porte ! »

\all{Probiere unseren Kaffee!} « Goûte notre café ! »

\all{Hilf mir bitte!} « Aide-moi, s'il te plaît ! »
\end{exemplebox}

\begin{attention}[Verbes à changement de voyelle e $\rightarrow$ i/ie : la voyelle change !]
Les verbes forts qui changent leur voyelle radicale \all{e} en \all{i} ou \all{ie} à la 2\up{e} et 3\up{e} personnes du singulier du présent (\all{du liest, er spricht}) conservent ce changement à l'impératif, et dans ce cas on n'ajoute \cle{jamais} de -e :
\begin{center}
\all{lesen $\rightarrow$ Lies!} \quad \all{sprechen $\rightarrow$ Sprich!} \quad \all{nehmen $\rightarrow$ Nimm!} \quad \all{geben $\rightarrow$ Gib!} \quad \all{essen $\rightarrow$ Iss!}
\end{center}
\all{Lies das Angebot genau!} « Lis l'offre attentivement ! »

\all{Nimm das billigere Produkt!} « Prends le produit le moins cher ! »

\textbf{Mais attention :} les verbes qui changent \all{a} en \all{ä} au présent (\all{du fährst, er schläft}) ne changent \cle{pas} de voyelle à l'impératif :
\begin{center}
\all{fahren $\rightarrow$ Fahr(e)!} \quad \all{schlafen $\rightarrow$ Schlaf(e)!} \quad \all{tragen $\rightarrow$ Trag(e)!}
\end{center}
On dit \all{Fahr vorsichtig!} « Conduis prudemment ! », jamais \all{Fähr!}.
\end{attention}

\begin{attention}[Verbes en -eln, -ern et radicaux en -d, -t, -ig]
Deux cas où le \cle{-e} se comporte différemment :
\begin{itemize}
\item Les verbes en \all{-eln} et \all{-ern} gardent toujours le -e : \all{Lächle!} « Souris ! », \all{Wandere mit uns!} « Randonne avec nous ! » (à l'oral, on entend souvent \all{Lächl!}, mais à l'écrit on garde le -e).
\item Les verbes dont le radical se termine en \all{-d}, \all{-t}, \all{-ig} ou en groupes de consonnes difficiles (\all{-tm}, \all{-fn}) prennent \cle{obligatoirement} le -e : \all{Warte!} « Attends ! », \all{Öffne!} « Ouvre ! », \all{Entschuldige!} « Excuse-toi ! », \all{Atme tief!} « Respire profondément ! »
\end{itemize}
\end{attention}

\courssub{Formation de l'impératif à la 2\up{e} personne du pluriel (ihr)}

\begin{propriete}[Impératif « ihr »]
À la 2\up{e} personne du pluriel, l'impératif est \cle{identique à la forme du présent} de l'indicatif, simplement \cle{sans le pronom} \all{ihr} :
\begin{center}
\all{ihr kommt $\rightarrow$ Kommt!} \quad \all{ihr lest $\rightarrow$ Lest!} \quad \all{ihr macht $\rightarrow$ Macht!}
\end{center}
Il n'y a ici \cle{aucun changement de voyelle} à gérer : la forme du présent suffit.
\end{propriete}

\begin{exemplebox}[Exemples traduits]
\all{Vergleicht die Preise!} « Comparez les prix ! » (à plusieurs personnes)

\all{Kauft jetzt und spart Geld!} « Achetez maintenant et économisez de l'argent ! »

\all{Kommt zu unserem Markt in Yaoundé!} « Venez à notre marché à Yaoundé ! »
\end{exemplebox}

\courssub{Forme de politesse (Sie) : l'inversion}

\begin{propriete}[Impératif « Sie » : infinitif + Sie]
À la forme de politesse, l'impératif se forme avec la forme verbale du présent à la 3\up{e} personne du pluriel (identique à l'\cle{infinitif}), suivie du pronom \all{Sie} placé \cle{après} le verbe (inversion) :
\begin{center}
\all{Kommen Sie!} \quad \all{Lesen Sie den Text!} \quad \all{Besuchen Sie uns!}
\end{center}
Contrairement au français (« Venez ! », « Lisez ! »), l'allemand \cle{conserve le pronom} \all{Sie} à l'impératif. C'est la seule forme d'impératif qui garde un pronom.
\end{propriete}

\begin{exemplebox}[Exemples traduits]
\all{Besuchen Sie unser Geschäft!} « Visitez notre magasin ! »

\all{Probieren Sie unsere neue Marke!} « Essayez notre nouvelle marque ! »

\all{Rufen Sie uns an!} « Appelez-nous ! »

\all{Entschuldigen Sie, wo ist die Kasse?} « Excusez-moi, où est la caisse ? »
\end{exemplebox}

\courssub{Cas particuliers : sein et les verbes à particule}

\begin{propriete}[L'impératif du verbe sein]
Le verbe \all{sein} (être) a des formes d'impératif irrégulières, à mémoriser :
\begin{center}
\begin{tabular}{|l|l|l|}
\hline
\rowcolor{popblueL} \textbf{Personne} & \textbf{Forme} & \textbf{Exemple} \\
\hline
du & \all{Sei!} & \all{Sei ruhig!} « Sois tranquille ! » \\
\hline
ihr & \all{Seid!} & \all{Seid vorsichtig!} « Soyez prudents ! » \\
\hline
Sie & \all{Seien Sie!} & \all{Seien Sie willkommen!} « Soyez le bienvenu ! » \\
\hline
\end{tabular}
\end{center}
\end{propriete}

\begin{propriete}[Verbes à particule séparable à l'impératif]
Avec les verbes à particule séparable, la règle habituelle s'applique : le verbe conjugué est en \cle{première position} et la \cle{particule va en fin de phrase} :
\begin{center}
\all{an/schauen $\rightarrow$ Schau dir das Angebot an!} « Regarde l'offre ! »

\all{mit/kommen $\rightarrow$ Kommt mit!} « Venez avec nous ! »

\all{an/rufen $\rightarrow$ Rufen Sie uns an!} « Appelez-nous ! »
\end{center}
\end{propriete}

\courssub{Tableau récapitulatif et comparaison avec le français}

\begin{popfigure}
\begin{tikzpicture}
\node[draw=popblue, fill=popblueL, rounded corners, font=\bfseries, align=center, minimum width=13.5cm] (titre) at (0,0) {Der Imperativ : drei Formen für drei Adressaten};
\node[draw=popgreen, fill=popgreenL, rounded corners, align=center, minimum width=4.2cm, minimum height=4.6cm, anchor=north] (du) at (-4.6,-0.6) {\textbf{du} (une personne)\\[2pt] radical (+ e)\\[4pt] \all{Komm!}\\ \all{Lies!}\\ \all{Sei!}\\ \all{Fang an!}};
\node[draw=poporange, fill=poporangeL, rounded corners, align=center, minimum width=4.2cm, minimum height=4.6cm, anchor=north] (ihr) at (0,-0.6) {\textbf{ihr} (plusieurs personnes)\\[2pt] présent sans \all{ihr}\\[4pt] \all{Kommt!}\\ \all{Lest!}\\ \all{Seid!}\\ \all{Fangt an!}};
\node[draw=poppurple, fill=poppurpleL, rounded corners, align=center, minimum width=4.2cm, minimum height=4.6cm, anchor=north] (sie) at (4.6,-0.6) {\textbf{Sie} (politesse)\\[2pt] infinitif + \all{Sie}\\[4pt] \all{Kommen Sie!}\\ \all{Lesen Sie!}\\ \all{Seien Sie!}\\ \all{Fangen Sie an!}};
\draw[-{Stealth}, thick, popblue] (titre.south) -- (du.north);
\draw[-{Stealth}, thick, popblue] (titre.south) -- (ihr.north);
\draw[-{Stealth}, thick, popblue] (titre.south) -- (sie.north);
\end{tikzpicture}
\legende{Les trois formes de l'impératif allemand avec les verbes kommen, lesen, sein et anfangen.}
\end{popfigure}

\begin{center}
\begin{tabularx}{\linewidth}{|X|X|X|}
\hline
\rowcolor{popblueL} \textbf{Français} & \textbf{Deutsch} & \textbf{Remarque} \\
\hline
Viens ! / Venez ! (à un ami) & \all{Komm!} & radical seul, sans pronom \\
\hline
Venez ! (à plusieurs) & \all{Kommt!} & comme au présent \\
\hline
Venez ! (politesse) & \all{Kommen Sie!} & le pronom \all{Sie} est obligatoire \\
\hline
Sois sage ! & \all{Sei brav!} & forme irrégulière de \all{sein} \\
\hline
Lis le texte ! & \all{Lies den Text!} & changement de voyelle e $\rightarrow$ ie \\
\hline
Conduis lentement ! & \all{Fahr langsam!} & pas de changement a $\rightarrow$ ä \\
\hline
\end{tabularx}
\end{center}

\begin{attention}[Erreur fréquente des francophones]
Ne dites \cle{jamais} \all{Kommen du!} ou \all{Du komm!}. À la forme \all{du}, l'impératif allemand n'a \cle{jamais} de terminaison \all{-en} et \cle{jamais} de pronom. Le français dit « Viens, toi ! » pour insister ; l'allemand dira \all{Komm du mal her!} seulement dans un oral très familier et insistant — à éviter à l'écrit. Autre piège : ne pas oublier le pronom \all{Sie} à la forme de politesse : \all{Besuchen Sie uns!} et non \all{Besuchen uns!}.
\end{attention}

\begin{savaistu}[L'impératif dans la publicité allemande]
Les slogans publicitaires allemands utilisent massivement l'impératif à la forme \all{du}, car la publicité veut créer une relation proche et amicale avec le client : \all{Mach dir keine Sorgen!} « Ne te fais pas de souci ! ». Mais il existe une exception célèbre : le slogan du fabricant de bonbons Haribo, \all{„Haribo macht Kinder froh, und Erwachsene ebenso“} (« Haribo rend les enfants joyeux, et les adultes aussi »), utilise la 3\up{e} personne du présent — une phrase affirmative, pas un impératif ! Ce slogan, créé dans les années 1930, est l'un des plus connus de l'histoire de la publicité allemande. Au Cameroun aussi, les slogans publicitaires adorent l'impératif : « Goûtez ! », « Essayez ! », « Rejoignez-nous ! ».
\end{savaistu}

\begin{exoresolu}[Mets les verbes à l'impératif]
Énoncé : mets les verbes entre parenthèses à l'impératif, à la personne indiquée.

1. (lesen, du) \all{\dots\dots{} das Angebot genau!}

2. (besuchen, Sie) \all{\dots\dots{} unser Geschäft in Douala!}

3. (vergleichen, ihr) \all{\dots\dots{} die Preise!}

4. (sein, du) \all{\dots\dots{} vorsichtig mit dem Geld!}

5. (an/rufen, Sie) \all{\dots\dots{} uns heute \dots\dots!}

6. (nehmen, du) \all{\dots\dots{} das billigere Produkt!}

\tcblower
Solution détaillée :

1. \all{Lies das Angebot genau!} — verbe fort e $\rightarrow$ ie : la voyelle change à l'impératif \all{du}, sans -e.

2. \all{Besuchen Sie unser Geschäft in Douala!} — forme de politesse : infinitif + \all{Sie} après le verbe.

3. \all{Vergleicht die Preise!} — forme \all{ihr} identique au présent, sans pronom.

4. \all{Sei vorsichtig mit dem Geld!} — impératif irrégulier de \all{sein} à la forme \all{du}.

5. \all{Rufen Sie uns heute an!} — verbe à particule : la particule \all{an} va en fin de phrase.

6. \all{Nimm das billigere Produkt!} — verbe fort e $\rightarrow$ i : \all{nimm}, sans -e.
\end{exoresolu}

\begin{exercice}[À toi de jouer]
1. Mets les verbes suivants aux trois formes de l'impératif : \all{kommen}, \all{helfen}, \all{sein}, \all{mit/bringen}.

2. Traduis en allemand : « Achetez (Sie) nos produits ! », « Essaie (du) ce yaourt ! », « Soyez (ihr) les bienvenus ! », « Regarde (du) la publicité ! » (verbe \all{an/schauen}).
\end{exercice}

\begin{corrige}[Exercice « À toi de jouer »]
1. \all{Komm! / Kommt! / Kommen Sie!} ; \all{Hilf! / Helft! / Helfen Sie!} (changement e $\rightarrow$ i à la forme \all{du}) ; \all{Sei! / Seid! / Seien Sie!} ; \all{Bring mit! / Bringt mit! / Bringen Sie mit!}

2. \all{Kaufen Sie unsere Produkte!} ; \all{Probier(e) diesen Joghurt!} ; \all{Seid willkommen!} ; \all{Schau die Werbung an!}
\end{corrige}

\begin{aretenir}[L'essentiel de l'impératif]
\begin{itemize}
\item Trois formes : \all{du} = radical (+ e facultatif) ; \all{ihr} = forme du présent sans pronom ; \all{Sie} = infinitif + \all{Sie} (le pronom est obligatoire).
\item Les verbes à changement de voyelle \all{e $\rightarrow$ i/ie} changent aussi à l'impératif \all{du} (\all{Lies!}, \all{Nimm!}), mais pas ceux en \all{a $\rightarrow$ ä} (\all{Fahr!}).
\item \all{sein} est irrégulier : \all{Sei! / Seid! / Seien Sie!}
\item Verbes à particule : la particule va en fin de phrase (\all{Schau dir das an!}).
\item Le verbe est toujours en première position ; c'est le mode typique de la publicité et des consignes.
\end{itemize}
\end{aretenir}

\coursec{Grammaire 2 : L'infinitif avec zu et la finalité (um ... zu / damit)}

Dans la section précédente, nous avons étudié l'impératif, ce mode qui permet à la publicité de s'adresser directement au consommateur : \all{Kaufen Sie jetzt!}, \all{Probieren Sie unser neues Produkt!} Mais la publicité et le marketing ne donnent pas seulement des ordres : ils expliquent aussi des buts, des intentions. Pourquoi les entreprises dépensent-elles des milliards ? Pourquoi baissent-elles leurs prix ? Pour répondre à ces questions, l'allemand dispose de deux constructions de la finalité : \cle{um ... zu} et \cle{damit}. C'est l'objet de cette deuxième leçon de grammaire, précédée d'un point complet sur l'infinitif avec \all{zu}.

\courssub{Observation : la finalité dans notre texte}

Reprenons deux phrases du texte support \all{„Werbung überall“} et observons-les attentivement.

\begin{exemplebox}[Deux phrases à observer]
\all{Firmen geben Milliarden aus, um ihre Produkte bekannt zu machen.} « Les entreprises dépensent des milliards pour faire connaître leurs produits. »

\all{Sie wollen den Kunden überzeugen, damit er ihre Marke kauft.} « Elles veulent convaincre le client pour qu'il achète leur marque. »
\end{exemplebox}

Posons-nous trois questions :

\begin{enumerate}
\item Dans la première phrase, qui dépense ? \emph{Die Firmen}. Qui fait connaître les produits ? \emph{Die Firmen} également. Le sujet est \textbf{identique} dans les deux actions. La construction utilisée est \all{um ... zu} + infinitif, placé en fin de phrase.
\item Dans la deuxième phrase, qui veut convaincre ? \emph{Sie} (les entreprises). Qui achète ? \emph{Er} (le client). Les sujets sont \textbf{différents}. La construction utilisée est \all{damit} + subordonnée, avec le verbe conjugué (\all{kauft}) rejeté en fin de phrase.
\item En français, les deux phrases emploient la même préposition « pour ». L'allemand, lui, \textbf{oblige à choisir} entre deux structures selon l'identité ou la différence des sujets.
\end{enumerate}

C'est toute la logique de la finalité allemande. Avant d'énoncer la règle, nous devons maîtriser l'outil de base : l'infinitif avec \all{zu}.

\courssub{L'infinitif avec zu (der Infinitiv mit zu)}

\begin{propriete}[Formation de l'infinitif avec zu]
Le \all{zu} se place \textbf{devant l'infinitif}, en un seul bloc, généralement en fin de proposition :

\all{Ich versuche, die Kunden zu überzeugen.} « J'essaie de convaincre les clients. »

\textbf{Cas particulier des verbes à particule séparable} : le \all{zu} s'insère \textbf{entre la particule et le radical}, et le tout s'écrit en \textbf{un seul mot} :

\all{anfangen} $\rightarrow$ \all{anzufangen} ; \all{bekannt machen} $\rightarrow$ \all{bekannt zu machen} ; \all{einkaufen} $\rightarrow$ \all{einzukaufen}.
\end{propriete}

\begin{exemplebox}[Exemples traduits]
\all{Es ist wichtig, früh aufzustehen.} « Il est important de se lever tôt. »

\all{Die Firma beginnt, ihre Produkte im Internet anzubieten.} « L'entreprise commence à offrir ses produits sur Internet. »

\all{Vergiss nicht, die Rechnung mitzunehmen!} « N'oublie pas d'emporter la facture ! »
\end{exemplebox}

\begin{attention}[La place du zu dans les verbes à particule]
Erreur classique du francophone : écrire \all{zu aufstehen} ou \all{zu einkaufen}. C'est faux. Le \all{zu} se glisse \textbf{à l'intérieur} du verbe : \all{aufzustehen}, \all{einzukaufen}, \all{mitzunehmen}. Retenez le schéma : particule + \all{zu} + radical, en un seul mot. Attention aussi aux verbes \textbf{inséparables} (\all{verstehen}, \all{bekommen}, \all{übersetzen}) : le \all{zu} se place devant, normalement : \all{zu verstehen}, \all{zu übersetzen}.
\end{attention}

\begin{propriete}[Verbes suivis de zu + infinitif]
De nombreux verbes construisent leur complément avec \all{zu} + infinitif (la virgule est recommandée, souvent obligatoire). À connaître pour le Bac :

\begin{center}
\begin{tabularx}{\linewidth}{|l|X|X|}
\hline
\rowcolor{popblueL} \textbf{Verbe} & \textbf{Exemple} & \textbf{Traduction} \\
\hline
\all{versuchen} & \all{Er versucht, den Kunden zu überzeugen.} & Il essaie de convaincre le client. \\
\hline
\all{beginnen} & \all{Sie beginnt, Deutsch zu lernen.} & Elle commence à apprendre l'allemand. \\
\hline
\all{hoffen} & \all{Wir hoffen, mehr zu verkaufen.} & Nous espérons vendre davantage. \\
\hline
\all{vergessen} & \all{Ich habe vergessen, zu zahlen.} & J'ai oublié de payer. \\
\hline
\all{beschließen} & \all{Die Firma beschließt, die Preise zu senken.} & L'entreprise décide de baisser les prix. \\
\hline
\end{tabularx}
\end{center}
\end{propriete}

\begin{aretenir}[Verbes SANS zu : le rappel de Première]
Certains verbes sont directement suivis de l'infinitif \textbf{sans} \all{zu} :
\begin{itemize}
\item les \textbf{verbes modaux} : \all{Ich kann Deutsch sprechen.} « Je peux parler allemand. » ; \all{Wir müssen arbeiten.} « Nous devons travailler. »
\item \all{sehen}, \all{hören}, \all{lassen} + infinitif : \all{Ich sehe ihn kommen.} « Je le vois venir. » ; \all{Sie hört das Kind weinen.} « Elle entend l'enfant pleurer. » ; \all{Er lässt das Auto reparieren.} « Il fait réparer la voiture. »
\end{itemize}
Dire \all{Ich kann zu sprechen} est une faute grave.
\end{aretenir}

\courssub{La finalité avec um ... zu}

\begin{propriete}[Règle de um ... zu]
Pour exprimer un but (« pour + infinitif » en français) lorsque le \textbf{sujet est le même} dans les deux actions, on utilise \all{um} en début de proposition infinitive et \all{zu} + infinitif \textbf{en toute fin} :

\all{Er lernt Deutsch, um in Deutschland zu studieren.} « Il apprend l'allemand pour étudier en Allemagne. »

Structure : [proposition principale] + \all{um} + [compléments] + \all{zu} + infinitif.
\end{propriete}

\begin{exemplebox}[Exemples traduits]
\all{Ich spare Geld, um ein Handy zu kaufen.} « J'économise de l'argent pour acheter un portable. »

\all{Die Firma wirbt im Fernsehen, um neue Kunden zu gewinnen.} « L'entreprise fait de la publicité à la télévision pour gagner de nouveaux clients. »

\all{Amina geht auf den Markt, um frisches Obst einzukaufen.} « Amina va au marché pour acheter des fruits frais. » (verbe à particule : \all{einzukaufen})

\all{Wir lernen fleißig, um das Bac zu bestehen.} « Nous travaillons avec application pour réussir le Bac. »
\end{exemplebox}

\courssub{La finalité avec damit}

\begin{propriete}[Règle de damit]
Lorsque les \textbf{sujets sont différents} dans les deux actions, on ne peut pas utiliser \all{um ... zu}. On construit alors une \textbf{subordonnée} introduite par la conjonction \all{damit}, avec le verbe \textbf{conjugué rejeté en fin de phrase} :

\all{Die Firma senkt die Preise, damit die Kunden mehr kaufen.} « L'entreprise baisse les prix pour que les clients achètent davantage. »

Structure : [proposition principale] + \all{damit} + sujet + [compléments] + verbe conjugué en fin.
\end{propriete}

\begin{exemplebox}[Exemples traduits]
\all{Der Vater gibt mir Geld, damit ich ein Handy kaufe.} « Mon père me donne de l'argent pour que j'achète un portable. »

\all{Die Lehrerin erklärt die Regel, damit die Schüler sie verstehen.} « La professeure explique la règle pour que les élèves la comprennent. »

\all{Ich spreche langsam, damit du mich verstehst.} « Je parle lentement pour que tu me comprennes. »
\end{exemplebox}

\begin{propriete}[Règle de la finalité : synthèse]
\begin{itemize}
\item \textbf{Même sujet} dans les deux propositions $\rightarrow$ \cle{um ... zu} + infinitif (pas de verbe conjugué).
\item \textbf{Sujets différents} $\rightarrow$ \cle{damit} + subordonnée (verbe conjugué en fin).
\end{itemize}
Le français traduit les deux par « pour » ou « pour que » ; c'est l'analyse des sujets qui impose le choix en allemand.
\end{propriete}

\begin{center}
\begin{tabularx}{\linewidth}{|X|X|X|}
\hline
\rowcolor{popblueL} \textbf{Critère} & \textbf{um ... zu} & \textbf{damit} \\
\hline
Sujets & identiques & différents \\
\hline
Verbe & infinitif avec \all{zu}, en fin & verbe conjugué, en fin \\
\hline
Français & pour + infinitif & pour que + subjonctif \\
\hline
Exemple & \all{Ich spare, um ein Handy zu kaufen.} & \all{Der Vater gibt mir Geld, damit ich ein Handy kaufe.} \\
\hline
\end{tabularx}
\end{center}

\begin{popfigure}
\begin{tikzpicture}[
  node distance=1.1cm,
  question/.style={diamond, draw=popblue, fill=popblueL, text width=3.2cm, align=center, inner sep=2pt, font=\small\bfseries},
  box/.style={rectangle, rounded corners, draw=popdark, fill=popcream, text width=5.2cm, align=center, inner sep=5pt, font=\small},
  fleche/.style={-{Stealth}, thick, popdark}]
\node[question] (q) {Même sujet dans les deux actions ?};
\node[box, below left=1.2cm and 0.2cm of q, draw=popgreen, fill=popgreenL, align=center] (oui){\textbf{OUI} : \all{um ... zu} + infinitif\\ \all{Er lernt, um das Bac zu bestehen.}};
\node[box, below right=1.2cm and 0.2cm of q, draw=poporange, fill=poporangeL, align=center] (non){\textbf{NON} : \all{damit} + verbe conjugué en fin\\ \all{Er erklärt, damit wir verstehen.}};
\draw[fleche] (q) -- (oui) node[midway, left, font=\small\bfseries, popgreen]{oui};
\draw[fleche] (q) -- (non) node[midway, right, font=\small\bfseries, poporange]{non};
\end{tikzpicture}
\legende{Schéma de décision : le choix entre \all{um ... zu} et \all{damit} dépend de l'identité des sujets.}
\end{popfigure}

\begin{methode}[Comment choisir entre um ... zu et damit]
\begin{enumerate}
\item \textbf{Identifier le sujet} de chaque verbe : qui fait la première action ? qui fait la deuxième ?
\item Si les sujets sont \textbf{identiques} $\rightarrow$ \all{um ... zu} + infinitif en fin de phrase.
\item Si les sujets sont \textbf{différents} $\rightarrow$ \all{damit} + sujet + verbe conjugué en fin de phrase.
\item Vérifier la place du \all{zu} : devant l'infinitif, ou entre la particule et le radical (\all{anzufangen}).
\end{enumerate}
\end{methode}

\begin{attention}[Pièges majeurs pour les francophones]
\begin{itemize}
\item \all{um dass} \textbf{n'existe pas}. Le français « pour que » se traduit \textbf{toujours} par \all{damit}, jamais par \all{um ... zu}. Faux : \all{Ich gebe dir Geld, um du ein Handy zu kaufen.} Correct : \all{Ich gebe dir Geld, damit du ein Handy kaufst.}
\item Jamais de \textbf{verbe conjugué} après \all{um ... zu} : seul l'infinitif est possible.
\item \all{damit} ne s'élide jamais : on écrit \all{damit ich}, \all{damit er}, sans apostrophe.
\item Ne pas confondre \all{damit} conjonction (« afin que ») et \all{damit} pronom adverbial (« avec cela ») : \all{Ich schreibe damit.} « J'écris avec cela. » Le contexte et la place du verbe permettent de trancher.
\end{itemize}
\end{attention}

\courssub{Texte d'application : la stratégie d'une jeune entreprise}

\begin{textebox}[Ein Start-up in Yaoundé]
Die junge Firma „MboaTech“ in Yaoundé hat ein neues Handy entwickelt. Um ihre Produkte bekannt zu machen, investiert sie viel Geld in Werbung: Plakate in der Stadt, Spots im Radio und Videos im Internet. Die Chefin, Frau Ngo Bell, erklärt: „Wir arbeiten Tag und Nacht, um die Kunden von unserer Qualität zu überzeugen.“

Zuerst versucht die Firma, die Jugendlichen zu erreichen. Sie organisiert Konzerte, damit die jungen Leute ihre Marke entdecken. Außerdem senkt sie die Preise, damit auch Schüler und Studenten das Handy kaufen können. „Es ist wichtig, billig anzufangen“, sagt Frau Ngo Bell, „denn niemand vergisst gern, ein gutes Angebot zu nutzen.“

Die Verkäufer lernen jeden Tag neue Methoden, um besser zu beraten. Sie hören den Kunden zu, um ihre Wünsche zu verstehen. Die Firma beschließt auch, ein Callcenter zu öffnen, damit die Kunden schnell Hilfe bekommen. „Wir hoffen, bald in Douala und Bafoussam zu verkaufen“, sagt die Chefin. „Und wir sparen Geld, um nächstes Jahr nach Deutschland zu exportieren.“

Die Strategie funktioniert: nach sechs Monaten kennt fast jeder in der Hauptstadt die Marke „MboaTech“. Die Firma hat es geschafft, ohne große Erfahrung anzufangen und trotzdem zu gewinnen.
\end{textebox}

\textbf{Glossaire} : \all{entwickeln} (développer) ; \all{das Plakat, -e} (l'affiche) ; \all{der Spot, -s} (le spot publicitaire) ; \all{erreichen} (atteindre) ; \all{das Angebot, -e} (l'offre) ; \all{nutzen} (profiter de) ; \all{beraten} (conseiller) ; \all{der Wunsch, „-e} (le souhait) ; \all{das Callcenter, -} (le centre d'appels) ; \all{exportieren} (exporter) ; \all{die Erfahrung, -en} (l'expérience) ; \all{es schaffen} (y arriver).

\textbf{Questions de compréhension} :
\begin{enumerate}
\item \all{Was macht die Firma, um ihre Produkte bekannt zu machen?}
\item \all{Warum organisiert sie Konzerte?}
\item \all{Warum senkt sie die Preise?}
\item Relevez dans le texte trois constructions de finalité et précisez pour chacune s'il s'agit de \all{um ... zu} ou de \all{damit}, en justifiant par les sujets.
\end{enumerate}

\begin{exoresolu}[Exercice résolu : transformation um ... zu / damit]
\textbf{Énoncé.} Transformez les phrases suivantes : remplacez \all{um ... zu} par \all{damit} (ou l'inverse) en adaptant le sujet entre parenthèses.

1. \all{Die Firma wirbt, um mehr Kunden zu gewinnen.} (les produits / être connus — sujet : \all{die Produkte})

2. \all{Ich lerne Deutsch, damit ich in Berlin studieren kann.} (sujet unique : \all{ich})

\tcblower
\textbf{Solution détaillée.}

1. Le nouveau sujet de la deuxième action est \all{die Produkte}, différent de \all{die Firma} : il faut donc \all{damit} + subordonnée, verbe conjugué en fin. On choisit le verbe \all{sein} avec le participe : \all{Die Firma wirbt, damit ihre Produkte bekannter werden.} « L'entreprise fait de la publicité pour que ses produits deviennent plus connus. »

2. Le sujet est identique (\all{ich}) dans les deux actions : on peut passer à \all{um ... zu}. Le verbe modal \all{kann} disparaît (les modaux n'ont pas d'infinitif avec \all{zu} dans cette construction) : \all{Ich lerne Deutsch, um in Berlin zu studieren.} « J'apprends l'allemand pour étudier à Berlin. »
\end{exoresolu}

\begin{exercice}[À vous]
Traduisez en allemand :
\begin{enumerate}
\item Les élèves travaillent beaucoup pour réussir l'examen.
\item Le professeur répète la règle pour que les élèves la comprennent.
\item L'entreprise fait de la publicité pour vendre plus de produits.
\item Ma mère me donne de l'argent pour que j'achète des fruits au marché.
\end{enumerate}
\end{exercice}

\begin{corrige}[Exercice : la finalité]
1. Sujet identique (\all{die Schüler}) : \all{Die Schüler arbeiten viel, um die Prüfung zu bestehen.}

2. Sujets différents (\all{der Lehrer} / \all{die Schüler}) : \all{Der Lehrer wiederholt die Regel, damit die Schüler sie verstehen.}

3. Sujet identique : \all{Die Firma macht Werbung, um mehr Produkte zu verkaufen.}

4. Sujets différents ; attention au verbe à particule \all{einkaufen} : \all{Meine Mutter gibt mir Geld, damit ich Obst auf dem Markt kaufe.} (ou \all{... damit ich Obst auf dem Markt einkaufe.})
\end{corrige}

\begin{savaistu}[Um ... zu et damit au Bac]
Au baccalauréat, la transformation \all{um ... zu} $\leftrightarrow$ \all{damit} est un exercice de grammaire quasi systématique dans les épreuves de traduction (thème). Le correcteur vérifie trois choses : le bon choix de la structure selon les sujets, la place du verbe (infinitif ou conjugué, toujours en fin), et la place exacte du \all{zu} dans les verbes à particule. Un entraînement régulier sur ce point rapporte des points faciles. Petit bonus culturel : dans la publicité allemande réelle, on trouve souvent des slogans construits sur la finalité, comme \all{Wir tun alles, um Sie zufrieden zu stellen} — « Nous faisons tout pour vous satisfaire ».
\end{savaistu}

\begin{aretenir}[L'essentiel de la section]
\begin{itemize}
\item L'infinitif avec \all{zu} : \all{zu} devant l'infinitif ; entre la particule et le radical pour les verbes séparables (\all{anzufangen}, \all{bekannt zu machen}).
\item Pas de \all{zu} après les modaux, \all{sehen}, \all{hören}, \all{lassen}.
\item Finalité, même sujet : \all{um ... zu} + infinitif. Sujets différents : \all{damit} + verbe conjugué en fin.
\item Le français « pour / pour que » correspond aux deux structures ; c'est l'analyse des sujets qui décide.
\end{itemize}
\end{aretenir}

\coursec{Méthode du commentaire et production guidée}

Nous savons désormais employer l'impératif pour formuler des slogans et construire des propositions de but avec \all{um ... zu} et \all{damit}. Ces outils grammaticaux ne sont pas des fins en soi : ils vont nous servir à \cle{commenter} le texte « Werbung überall » et à \cle{produire} nos propres écrits. Cette dernière étape est décisive : au baccalauréat, l'épreuve d'allemand demande toujours de passer de la compréhension à l'expression personnelle. Apprenons donc une méthode simple, rigoureuse et transferable à tous les textes.

\courssub{La méthode du commentaire de texte en quatre étapes}

\begin{methode}[Comment rédiger un commentaire de texte au Bac]
Un commentaire de texte (\all{der Kommentar}) se construit toujours en quatre étapes, dans cet ordre :
\begin{enumerate}
\item \textbf{Présenter le texte} : indiquer le titre, le thème et la source. Formule obligatoire : \all{Der Text mit dem Titel „...“ handelt von ...} (Le texte intitulé « ... » traite de ...).
\item \textbf{Dégager l'idée principale et la structure} : résumer le contenu en une ou deux phrases et annoncer le plan avec \all{Zuerst ..., dann ..., schließlich ...} (d'abord..., ensuite..., enfin...).
\item \textbf{Analyser les arguments et les procédés} : relever les slogans, les impératifs, les chiffres, les exemples, et expliquer leur effet sur le lecteur.
\item \textbf{Donner son avis argumenté} : exprimer son opinion personnelle avec \all{meiner Meinung nach}, la justifier avec \all{denn} ou \all{weil}, et conclure avec \all{zum Schluss}.
\end{enumerate}
\end{methode}

\begin{popfigure}
\begin{center}
\begin{tikzpicture}[
node distance=0.9cm,
etape/.style={rectangle, rounded corners, draw=popblue, fill=popblueL, text width=4.6cm, align=center, font=\small\bfseries, minimum height=1cm},
fleche/.style={-{Stealth[length=3mm]}, thick, poporange}
]
\node[etape, align=center] (e1){1. Présentation\\ \footnotesize Titel, Thema, Quelle};
\node[etape, below=of e1, fill=popgreenL, draw=popgreen, align=center] (e2){2. Idée principale\\ \footnotesize Zuerst, dann, schließlich};
\node[etape, below=of e2, fill=poporangeL, draw=poporange, align=center] (e3){3. Analyse\\ \footnotesize Slogans, Imperativ, Zahlen};
\node[etape, below=of e3, fill=poppinkL, draw=poppink, align=center] (e4){4. Avis personnel\\ \footnotesize Meiner Meinung nach ...};
\draw[fleche] (e1) -- (e2);
\draw[fleche] (e2) -- (e3);
\draw[fleche] (e3) -- (e4);
\end{tikzpicture}
\end{center}
\legende{Organigramme du commentaire de texte : quatre étapes enchaînées par des flèches, de la présentation à l'avis personnel.}
\end{popfigure}

\courssub{La banque de connecteurs du commentaire}

Les connecteurs (\all{die Konnektoren}) sont les articulations du commentaire. Sans eux, le texte est une suite de phrases isolées ; avec eux, il devient un raisonnement. Voici le tableau de référence à apprendre par cœur.

\begin{center}
\begin{tabularx}{\linewidth}{|l|X|X|}
\hline
\rowcolor{popblueL} \textbf{Connecteur} & \textbf{Français} & \textbf{Exemple} \\
\hline
\all{zuerst} & d'abord & \all{Zuerst zeigt der Autor das Problem.} \\
\hline
\all{dann} & ensuite & \all{Dann erklärt er die Strategien.} \\
\hline
\all{einerseits ... andererseits} & d'une part ... d'autre part & \all{Einerseits informiert die Werbung, andererseits manipuliert sie.} \\
\hline
\all{außerdem} & de plus & \all{Außerdem nennt er Beispiele.} \\
\hline
\all{zum Beispiel} & par exemple & \all{Viele Firmen, zum Beispiel Getränkefirmen, benutzen kurze Slogans.} \\
\hline
\all{meiner Meinung nach} & à mon avis & \all{Meiner Meinung nach ist Werbung nützlich.} \\
\hline
\all{denn / weil} & car / parce que & \all{Sie ist nützlich, denn sie informiert.} \\
\hline
\all{deshalb} & c'est pourquoi & \all{Deshalb muss man vorsichtig sein.} \\
\hline
\all{zum Schluss} & pour conclure & \all{Zum Schluss: Werbung ist überall.} \\
\hline
\end{tabularx}
\end{center}

\begin{attention}[L'inversion après « meiner Meinung nach »]
Piège classique du francophone : en français, « à mon avis, la publicité est importante » garde l'ordre sujet-verbe. En allemand, \all{meiner Meinung nach} occupe la première position de la phrase : le verbe conjugué doit donc venir en \cle{deuxième position}, avant le sujet. On dit : \all{Meiner Meinung nach \textbf{ist} die Werbung wichtig.} et jamais \all{Meiner Meinung nach die Werbung ist wichtig.} Même règle après \all{deshalb}, \all{zuerst}, \all{dann} et \all{außerdem} placés en tête : \all{Deshalb \textbf{kaufe} ich das Produkt.} « C'est pourquoi j'achète le produit. »
\end{attention}

\begin{attention}[« denn » ou « weil » ?]
Les deux se traduisent par « parce que / car », mais ils n'ont pas la même syntaxe. Après \all{denn} (coordination), l'ordre est normal : \all{Sie ist nützlich, denn sie \textbf{informiert} über Angebote.} Après \all{weil} (subordination), le verbe conjugué va \cle{à la fin} : \all{Sie ist nützlich, weil sie über Angebote \textbf{informiert}.} Au baccalauréat, \all{denn} est plus sûr pour un élève de Tle A.
\end{attention}

\begin{exemplebox}[Exemples traduits]
\all{Meiner Meinung nach beeinflusst die Werbung unsere Wünsche, denn sie zeigt uns immer neue Produkte.} « À mon avis, la publicité influence nos désirs, car elle nous montre toujours de nouveaux produits. »

\all{Einerseits informiert die Werbung über Angebote, andererseits drängt sie uns zum Kaufen.} « D'une part, la publicité informe sur les offres ; d'autre part, elle nous pousse à acheter. »

\all{Außerdem benutzen die Firmen kurze Slogans, damit die Kunden sie nicht vergessen.} « De plus, les entreprises utilisent des slogans courts afin que les clients ne les oublient pas. »

\all{Deshalb sollte man die Preise vergleichen, um nicht zu viel Geld auszugeben.} « C'est pourquoi il faudrait comparer les prix afin de ne pas dépenser trop d'argent. »
\end{exemplebox}

\courssub{Commentaire modèle sur « Werbung überall »}

\begin{textebox}[Commentaire modèle — niveau Tle A]
Der Text „Werbung überall“ handelt von der Werbung und ihren Strategien. Zuerst zeigt der Autor, dass Werbung überall ist: in der Stadt, im Radio und im Internet. Dann erklärt er, wie Firmen die Kunden überzeugen, zum Beispiel mit kurzen Slogans im Imperativ. Außerdem nennt er die Kritik: Die Werbung manipuliert uns. Meiner Meinung nach ist Werbung nützlich, denn sie informiert über Angebote. Aber man muss vorsichtig sein, um nicht zu viel Geld auszugeben. Zum Schluss kann man sagen: Werbung ist ein Teil unseres Lebens.
\end{textebox}

Analysons les connecteurs de ce modèle. \all{Zuerst} et \all{dann} organisent le résumé de la structure du texte ; \all{außerdem} ajoute un élément d'analyse (la critique) ; \all{meiner Meinung nach} introduit l'avis personnel, avec inversion du verbe (\all{ist} avant \all{Werbung}) ; \all{denn} justifie cet avis ; la proposition \all{um nicht zu viel Geld auszugeben} exprime la finalité avec un infinitif à \all{zu} (noter la particule séparable : \all{auszugeben} en un seul mot à la fin) ; enfin \all{zum Schluss} annonce la conclusion. Remarque : ce commentaire tient en huit lignes — c'est exactement la longueur attendue au baccalauréat.

\courssub{Production guidée 1 : le résumé en quatre phrases}

\begin{exoresolu}[Résumé du texte avec connecteurs imposés]
Énoncé : résume le texte « Werbung überall » en quatre phrases, en utilisant obligatoirement les connecteurs \all{zuerst}, \all{dann}, \all{außerdem} et \all{zum Schluss}.
\tcblower
Solution détaillée :

\all{Zuerst zeigt der Text, dass die Werbung überall in unserem Leben ist.} « D'abord, le texte montre que la publicité est partout dans notre vie. »

\all{Dann beschreibt er die Strategien der Firmen, zum Beispiel kurze Slogans und schöne Bilder.} « Ensuite, il décrit les stratégies des entreprises, par exemple des slogans courts et de belles images. »

\all{Außerdem erklärt er, dass die Werbung die Kunden manchmal manipuliert.} « De plus, il explique que la publicité manipule parfois les clients. »

\all{Zum Schluss sagt der Autor, dass man kritisch denken muss.} « Pour conclure, l'auteur dit qu'il faut penser de manière critique. »

Méthode : une phrase par grande partie du texte, un connecteur par phrase, et le verbe conjugué toujours en deuxième position dans la proposition principale (à la fin après \all{dass}).
\end{exoresolu}

\courssub{Production guidée 2 : créer des slogans camerounais}

Un bon slogan publicitaire combine souvent un \cle{impératif} (pour interpeller le client) et une proposition de \cle{finalité} avec \all{um ... zu} (pour promettre un bénéfice). Structure modèle : \all{Impératif + Produit, um + Infinitif !} Rappel : on choisit \all{um ... zu} quand le sujet des deux verbes est le même (toi, le client) ; on choisit \all{damit} quand les sujets sont différents (\all{Kauf unser Produkt, damit wir wachsen können!} « Achète notre produit afin que nous puissions grandir ! »).

\begin{exemplebox}[Trois slogans pour des produits camerounais]
\all{Probier unseren Kaffee aus dem Westen, um den Tag gut zu beginnen!} « Goûte notre café de l'Ouest pour bien commencer la journée ! »

\all{Kauf unseren Kakao aus dem Zentrum, um die beste Schokolade zu genießen!} « Achète notre cacao du Centre pour savourer le meilleur chocolat ! »

\all{Trink unseren Fruchtsaft, um gesund und stark zu bleiben!} « Bois notre jus de fruits pour rester en bonne santé et fort ! »
\end{exemplebox}

\begin{exercice}[À toi de créer]
Rédige deux slogans supplémentaires : un pour le miel du Nord (\all{der Honig}) et un pour l'ananas de la péninsule (\all{die Ananas}). Utilise chaque fois un impératif à la deuxième personne du singulier (\all{probier}, \all{kauf}, \all{trink}, \all{genieß}) et une proposition avec \all{um ... zu}.
\end{exercice}

\begin{corrige}[Production guidée 2]
Réponses possibles : \all{Probier unseren Honig aus dem Norden, um deine Familie gesund zu ernähren!} « Goûte notre miel du Nord pour nourrir sainement ta famille ! » — \all{Genieß unsere Ananas, um die Frische Kameruns zu schmecken!} « Savoure nos ananas pour goûter la fraîcheur du Cameroun ! » Toute réponse respectant l'impératif en début de phrase et l'infinitif avec \all{zu} en fin de proposition est correcte.
\end{corrige}

\courssub{Production guidée 3 : le paragraphe d'opinion}

\begin{methode}[Rédiger un court paragraphe d'opinion en cinq ou six lignes]
Sujet : \all{Beeinflusst die Werbung unsere Wünsche?} (La publicité influence-t-elle nos désirs ?)
\begin{enumerate}
\item Commence par ta position : \all{Meiner Meinung nach beeinflusst die Werbung unsere Wünsche.} (attention à l'inversion).
\item Justifie avec \all{denn} ou \all{weil} : \all{denn sie zeigt uns jeden Tag neue Produkte.}
\item Ajoute un exemple avec \all{zum Beispiel} et un impératif de slogan : \all{Zum Beispiel sagen die Slogans: „Kauf jetzt!“}
\item Nuance ou conseille avec \all{um ... zu} ou \all{damit} : \all{Wir müssen die Preise vergleichen, um nicht alles zu kaufen.}
\item Conclus avec \all{deshalb} ou \all{zum Schluss} : \all{Deshalb ist es wichtig, kritisch zu bleiben.}
\end{enumerate}
\end{methode}

\begin{exemplebox}[Paragraphe d'opinion modèle]
\all{Meiner Meinung nach beeinflusst die Werbung unsere Wünsche sehr stark, denn wir sehen sie überall: im Fernsehen, im Internet und in der Stadt. Zum Beispiel sagen viele Slogans: „Kauf jetzt, um glücklich zu sein!“ Die Firmen benutzen schöne Bilder, damit wir ihre Produkte kaufen. Deshalb müssen wir vorsichtig sein und die Preise vergleichen, um nicht zu viel Geld auszugeben. Zum Schluss: Die Werbung ist nützlich, aber wir dürfen ihr nicht alles glauben.} « À mon avis, la publicité influence fortement nos désirs, car nous la voyons partout : à la télévision, sur Internet et en ville. Par exemple, beaucoup de slogans disent : "Achète maintenant pour être heureux !" Les entreprises utilisent de belles images afin que nous achetions leurs produits. C'est pourquoi nous devons être prudents et comparer les prix afin de ne pas dépenser trop d'argent. Pour conclure : la publicité est utile, mais nous ne devons pas tout croire. »
\end{exemplebox}

\begin{exercice}[Mini-commentaire personnel]
Rédige un mini-commentaire de cinq à six phrases sur le sujet : \all{Ist Werbung in Kamerun nützlich oder gefährlich?} (La publicité au Cameroun est-elle utile ou dangereuse ?) Contraintes : commence par \all{Meiner Meinung nach}, utilise au moins un impératif de slogan entre guillemets allemands, une proposition avec \all{um ... zu} ou \all{damit}, et termine par \all{Zum Schluss}.
\end{exercice}

\begin{corrige}[Mini-commentaire personnel]
Réponse possible : \all{Meiner Meinung nach ist die Werbung in Kamerun nützlich, denn sie informiert die Kunden über neue Produkte und Preise. In Yaoundé und Douala sieht man überall Slogans wie „Kauf jetzt, um Geld zu sparen!“ Einerseits hilft die Werbung den Firmen, andererseits drängt sie die Leute zum Kaufen. Deshalb müssen die Kunden die Preise vergleichen, um nicht zu viel auszugeben. Zum Schluss: Die Werbung ist nützlich, aber man muss kritisch bleiben.} « À mon avis, la publicité au Cameroun est utile, car elle informe les clients sur les nouveaux produits et les prix. À Yaoundé et à Douala, on voit partout des slogans comme "Achète maintenant pour économiser de l'argent !" D'une part, la publicité aide les entreprises ; d'autre part, elle pousse les gens à acheter. C'est pourquoi les clients doivent comparer les prix afin de ne pas trop dépenser. Pour conclure : la publicité est utile, mais il faut rester critique. »
\end{corrige}

\courssub{Grille d'auto-évaluation}

Avant de rendre ta production, vérifie chaque critère de cette grille. Elle correspond exactement aux attentes du correcteur au baccalauréat.

\begin{center}
\begin{tabularx}{\linewidth}{|X|l|l|}
\hline
\rowcolor{popblueL} \textbf{Critère} & \textbf{Exigence} & \textbf{Oui / Non} \\
\hline
Vocabulaire du thème & au moins 3 mots : \all{die Werbung}, \all{das Produkt}, \all{der Kunde}, \all{die Marke}... & \\
\hline
Impératif & au moins 1 impératif correct (\all{Kauf!}, \all{Probier!}) & \\
\hline
Finalité & au moins 1 \all{um ... zu} ou 1 \all{damit} & \\
\hline
Connecteurs & au moins 3 connecteurs différents & \\
\hline
Place du verbe & verbe conjugué en 2e position, inversion après \all{meiner Meinung nach} & \\
\hline
Orthographe & majuscules aux noms, Umlaute, infinitif avec \all{zu} en fin de phrase & \\
\hline
\end{tabularx}
\end{center}

\begin{aretenir}[L'essentiel de la section]
Un commentaire de texte suit quatre étapes : présentation (\all{Der Text mit dem Titel ... handelt von ...}), idée principale et plan (\all{Zuerst ..., dann ..., schließlich ...}), analyse des procédés, avis personnel (\all{Meiner Meinung nach ...}). Les connecteurs structurent le raisonnement ; après \all{meiner Meinung nach} et \all{deshalb}, le verbe conjugué passe en deuxième position. Les slogans publicitaires combinent impératif et finalité (\all{um ... zu} si le sujet est identique, \all{damit} s'il est différent). Toute production se relit avec la grille d'auto-évaluation.
\end{aretenir}

\begin{savaistu}[Des slogans allemands pour le Cameroun]
Au Cameroun, les slogans publicitaires sont le plus souvent rédigés en français, en anglais ou en pidgin, sur les murs des villes, les bendskins et les panneaux des marchés de Yaoundé et de Douala. Créer des slogans en allemand pour des produits locaux — le café de l'Ouest, le cacao du Centre, le jus de fruits du Littoral — est un excellent exercice de transfert culturel : il oblige à penser en allemand une réalité camerounaise, exactement ce que demande l'épreuve d'expression écrite du baccalauréat.
\end{savaistu}

\newpage

\coursec{Activité d'intégration}

\courssub{Objectifs de l'activité}

À la fin de cette activité d'intégration, tu seras capable de :
\begin{itemize}
\item mobiliser le \cle{vocabulaire thématique} du marketing et de la publicité (\all{die Werbung}, \all{das Produkt}, \all{der Kunde}, \all{die Marke}, \all{die Zielgruppe}…) dans une situation de communication authentique ;
\item employer correctement l'\cle{impératif} (formes \all{du}, \all{ihr}, \all{Sie}) dans un slogan publicitaire adapté à la cible ;
\item exprimer la \cle{finalité} avec \all{um … zu} et \all{damit} dans un texte de présentation ;
\item présenter oralement un projet en deux minutes avec une prononciation soignée ;
\item justifier un choix en allemand avec \all{weil} (\all{Ich wähle Gruppe 2, weil …}).
\end{itemize}

\courssub{Situation}

Tu travailles dans la \all{Werbeagentur} (agence de publicité) « Kamerun Kreativ », une jeune agence installée à Douala. Une entreprise allemande d'importation, la \all{Afrika Import GmbH} de Hambourg, cherche des produits camerounais de qualité pour le marché allemand. Elle lance un appel d'offres : quelle agence saura le mieux promouvoir un produit camerounais auprès des consommateurs allemands ? Ton agence veut gagner ce contrat !

\begin{textebox}[E-Mail der Afrika Import GmbH]
\textbf{Betreff: Werbekampagne für kamerunische Produkte}\par
\medskip
Sehr geehrte Damen und Herren,\par
\medskip
unsere Firma möchte kamerunische Produkte auf dem deutschen Markt verkaufen: Kaffee, Kakao, Palmöl oder das traditionelle Ndop-Gewebe. Wir suchen eine Werbeagentur, die für eines dieser Produkte eine Kampagne erstellt.\par
\medskip
Wir erwarten: ein Plakat mit einem Slogan, einen kurzen Präsentationstext (fünf bis sechs Sätze) und eine mündliche Präsentation von zwei Minuten. Definieren Sie bitte Ihre Zielgruppe genau: Wer soll das Produkt kaufen? Junge Leute? Familien? Sportler?\par
\medskip
Die beste Kampagne bekommt den Auftrag. Wir freuen uns auf Ihre Ideen!\par
\medskip
Mit freundlichen Grüßen\par
Sabine Berger, Marketingchefin
\end{textebox}

\begin{definition}[Glossaire de l'e-mail]
\begin{itemize}
\item \all{der Auftrag, die Aufträge} : la commande, le contrat
\item \all{das Plakat, die Plakate} : l'affiche
\item \all{der Slogan, die Slogans} : le slogan
\item \all{die Zielgruppe, die Zielgruppen} : le groupe cible
\item \all{die mündliche Präsentation, -en} : la présentation orale
\item \all{das Gewebe, die Gewebe} : le tissu
\item \all{erwarten} : attendre, exiger (ici : « nous attendons »)
\item \all{sich freuen auf} + accusatif : se réjouir de (quelque chose à venir)
\end{itemize}
\end{definition}

\courssub{Consignes}

Par groupes de quatre, vous formez une équipe de l'agence « Kamerun Kreativ ». Suivez les étapes dans l'ordre.

\begin{enumerate}
\item \textbf{Lecture de l'e-mail (compréhension).} Lisez le courriel de Madame Berger et répondez par écrit en allemand : \all{Was möchte die Firma verkaufen? Was erwartet sie von der Agentur? Wer ist Sabine Berger?}
\item \textbf{Choix du produit et de la cible.} Choisissez un produit : \all{der Kaffee}, \all{der Kakao}, \all{das Palmöl} ou \all{das Ndop-Gewebe}. Définissez votre \all{Zielgruppe} en une phrase : \all{Unsere Zielgruppe sind junge Leute zwischen 18 und 30 Jahren.}
\item \textbf{Création de l'affiche.} Rédigez un slogan à l'\cle{impératif}. Attention au choix de la forme : \all{du} ou \all{ihr} pour un public jeune, \all{Sie} pour un public adulte ou formel. Ajoutez le nom de la marque et une image dessinée.
\item \textbf{Rédaction du texte de présentation.} Écrivez cinq à six phrases qui présentent le produit et la campagne. Votre texte doit contenir \textbf{au moins deux structures \all{um … zu}} et \textbf{au moins une phrase avec \all{damit}}.
\item \textbf{Présentation orale.} Présentez votre campagne à la classe en deux minutes : chaque membre du groupe parle environ trente secondes. Soignez la prononciation (le \all{w} se prononce « v », le \all{v} se prononce « f »).
\item \textbf{Vote et justification.} La classe vote pour la meilleure campagne. Chaque élève justifie son vote à voix haute : \all{Ich wähle Gruppe 2, weil ihr Slogan originell ist.} Attention : après \all{weil}, le verbe conjugué va à la fin !
\end{enumerate}

\begin{attention}[Critères d'évaluation]
L'enseignant évalue selon quatre critères : (1) le \cle{lexique thématique} du marketing est-il riche et correct ? (2) l'\cle{impératif} est-il bien formé et adapté à la cible ? (3) les structures de \cle{finalité} (\all{um … zu}, \all{damit}) sont-elles correctes ? (4) la \cle{prononciation} et l'aisance orale sont-elles satisfaisantes ?
\end{attention}

\courssub{Ressources à mobiliser}

\begin{aretenir}[Boîte à outils linguistique]
\textbf{Lexique :} \all{das Produkt, die Produkte} (le produit) ; \all{die Marke, die Marken} (la marque) ; \all{der Kunde, die Kunden} (le client) ; \all{das Angebot, die Angebote} (l'offre) ; \all{der Verkauf, die Verkäufe} (la vente) ; \all{die Qualität, die Qualitäten} (la qualité) ; \all{der Preis, die Preise} (le prix) ; \all{natürlich} (naturel) ; \all{gesund} (sain) ; \all{fair} (équitable).\par
\medskip
\textbf{Impératif :} \all{du} : radical sans terminaison (\all{Probier} !) ; \all{ihr} : radical + \all{-t} (\all{Probiert} !) ; \all{Sie} : infinitif + \all{Sie} (\all{Probieren Sie} !). Irréguliers : \all{Sei} ! / \all{Seid} ! / \all{Seien Sie} ! ; \all{Nimm} ! / \all{Nehmt} ! / \all{Nehmen Sie} !\par
\medskip
\textbf{Finalité :} \all{um … zu} + infinitif (même sujet) : \all{Wir werben, um das Produkt bekannt zu machen.} ; \all{damit} + proposition subordonnée (sujets différents) : \all{Wir machen Werbung, damit die Kunden das Produkt kaufen.}
\end{aretenir}

\begin{center}
\begin{tabular}{|l|l|l|l|}
\hline
\rowcolor{popblueL} Infinitif & du & ihr & Sie \\
\hline
\all{probieren} & \all{Probier(e)!} & \all{Probiert!} & \all{Probieren Sie!} \\
\hline
\all{kaufen} & \all{Kauf(e)!} & \all{Kauft!} & \all{Kaufen Sie!} \\
\hline
\all{nehmen} & \all{Nimm!} & \all{Nehmt!} & \all{Nehmen Sie!} \\
\hline
\all{lesen} & \all{Lies!} & \all{Lest!} & \all{Lesen Sie!} \\
\hline
\all{sein} & \all{Sei!} & \all{Seid!} & \all{Seien Sie!} \\
\hline
\end{tabular}
\end{center}

\begin{attention}[Pièges à éviter]
\begin{itemize}
\item Ne dis pas \all{*Probieren du !} : à l'impératif \all{du}, on utilise le radical seul : \all{Probiere !} ou \all{Probier !}
\item Les verbes qui changent \all{e} en \all{i} ou \all{ie} au radical prennent cette voyelle à l'impératif \all{du} : \all{nehmen} $\rightarrow$ \all{Nimm!}, \all{lesen} $\rightarrow$ \all{Lies!} (mais jamais d'Umlaut : \all{fahren} $\rightarrow$ \all{Fahr!}, et non \all{*Fähr!}).
\item Après \all{um … zu}, l'infinitif se place \textbf{à la fin} : \all{um mehr Kunden zu gewinnen}.
\item Après \all{damit} et \all{weil}, le verbe conjugué va \textbf{à la fin} de la subordonnée : \all{…, damit die Deutschen unseren Kaffee entdecken.}
\item N'oublie pas la majuscule à tous les noms : \all{der Kaffee}, \all{die Marke}, \all{das Plakat}.
\end{itemize}
\end{attention}

\courssub{Proposition de production et corrigé commenté}

\begin{exoresolu}[Modèle de campagne : le café de l'Ouest]
Énoncé : rédigez pour le groupe « Kamerun Kreativ » le slogan, la définition de la cible et le texte de présentation (5-6 phrases) d'une campagne pour le café camerounais destiné à des jeunes Allemands.
\tcblower
\textbf{Solution détaillée.}\par
\medskip
\textbf{Slogan (impératif \all{du}, cible jeune) :}\par
\all{Probier Kamerun-Kaffee — starte deinen Tag mit Energie!}\par
« Goûte le café du Cameroun — commence ta journée avec de l'énergie ! »\par
\medskip
\textbf{Zielgruppe :}\par
\all{Unsere Zielgruppe sind junge Leute zwischen 18 und 30 Jahren, die gern Kaffee trinken und faire Produkte kaufen.}\par
« Notre groupe cible, ce sont les jeunes entre 18 et 30 ans, qui aiment boire du café et acheter des produits équitables. »\par
\medskip
\textbf{Texte de présentation :}\par
\all{Unser Produkt ist Kaffee aus den Bergen Westkameruns. Die Marke heißt „Bamenda Gold“. Wir machen eine Kampagne in den sozialen Medien, um junge Kunden zu erreichen. Auf dem Plakat sieht man einen jungen Mann mit einer Tasse Kaffee, um Lust auf das Produkt zu wecken. Wir erzählen die Geschichte der Bauern, damit die Kunden wissen, woher der Kaffee kommt. Der Kaffee ist natürlich und fair, und der Preis ist attraktiv. Probieren Sie „Bamenda Gold“ — oder besser: Probier es einfach!}\par
\medskip
« Notre produit est un café des montagnes de l'Ouest du Cameroun. La marque s'appelle „Bamenda Gold“. Nous faisons une campagne sur les réseaux sociaux pour atteindre de jeunes clients. Sur l'affiche, on voit un jeune homme avec une tasse de café, pour susciter l'envie du produit. Nous racontons l'histoire des paysans, afin que les clients sachent d'où vient le café. Le café est naturel et équitable, et le prix est attractif. Goûtez „Bamenda Gold“ — ou plutôt : goûte-le tout simplement ! »\par
\medskip
\textbf{Commentaire des choix de langue :}
\begin{itemize}
\item \textbf{Impératif \all{du}} dans le slogan (\all{Probier}, \all{starte}) : la cible est jeune, le tutoiement crée la proximité. Pour un public plus âgé, on aurait choisi \all{Probieren Sie} !
\item \textbf{Deux structures \all{um … zu}} : \all{um junge Kunden zu erreichen} et \all{um Lust auf das Produkt zu wecken} — l'infinitif est bien placé à la fin, et le sujet de la principale et de l'infinitive est le même (\all{wir} / \all{man} renvoyant à l'agence).
\item \textbf{Une structure \all{damit}} : \all{damit die Kunden wissen, woher der Kaffee kommt} — ici les sujets sont différents (\all{wir} / \all{die Kunden}), donc \all{damit} est obligatoire ; le verbe \all{wissen} est à la fin.
\item \textbf{Lexique thématique} : \all{das Produkt}, \all{die Marke}, \all{die Kampagne}, \all{die Kunden}, \all{das Plakat}, \all{der Preis}.
\item \textbf{Justification de vote possible} : \all{Ich wähle diese Gruppe, weil der Slogan kurz und stark ist und weil der Text die Geschichte der Bauern erzählt.}
\end{itemize}
\end{exoresolu}

\courssub{Bilan de l'activité}

Cette activité t'a permis de réinvestir l'ensemble des acquis de la leçon dans une situation professionnelle réaliste. Retiens l'essentiel :
\begin{itemize}
\item le vocabulaire du marketing (\all{die Werbung}, \all{das Produkt}, \all{die Marke}, \all{die Zielgruppe}, \all{der Kunde}) s'emploie dans des contextes concrets de la vie économique ;
\item le choix de la forme de l'impératif (\all{du}, \all{ihr} ou \all{Sie}) dépend toujours du \textbf{public cible} : c'est une décision de communication, pas seulement de grammaire ;
\item \all{um … zu} s'emploie quand le sujet est identique, \all{damit} quand les sujets sont différents — et dans les deux cas, l'ordre des mots allemand doit être respecté ;
\item une bonne présentation orale repose sur un texte clair, une prononciation soignée et une répartition équilibrée de la parole dans le groupe.
\end{itemize}

\begin{experience}[Pour aller plus loin]
À la maison, transforme ton affiche en véritable maquette : dessine le produit, écris le slogan en gros caractères et ajoute deux phrases publicitaires avec \all{um … zu}. Tu peux aussi enregistrer ta présentation orale et vérifier ta prononciation du \all{w} (« v ») et du \all{v} (« f »).
\end{experience}

\newpage

\coursec{Méthodes et exercices résolus}

\courssub{Lexique thématique : Marketing und Werbung}

\begin{definition}[Les mots-clés du thème : die Werbung und das Marketing]
\begin{center}
\begin{tabularx}{\linewidth}{|l|l|X|}
\hline
\rowcolor{popblueL} \textbf{Deutsch} & \textbf{Français} & \textbf{Exemple} \\
\hline
die Werbung, -en & la publicité & \all{Die Werbung beeinflusst die Kunden.} \\
\hline
das Produkt, -e & le produit & \all{Das Produkt ist neu auf dem Markt.} \\
\hline
der Kunde, -n / die Kundin, -nen & le client / la cliente & \all{Der Kunde vergleicht die Preise.} \\
\hline
das Angebot, -e & l'offre & \all{Das Angebot gilt nur bis Freitag.} \\
\hline
der Verkauf, -e & la vente & \all{Der Verkauf beginnt um acht Uhr.} \\
\hline
die Marke, -n & la marque & \all{Diese Marke ist bei Jugendlichen beliebt.} \\
\hline
der Preis, -e & le prix & \all{Die Firma senkt die Preise.} \\
\hline
der Markt, -e & le marché & \all{Der Markt in Douala ist sehr groß.} \\
\hline
der Verkäufer, - / die Verkäuferin, -nen & le vendeur / la vendeuse & \all{Die Verkäuferin erklärt das Angebot.} \\
\hline
der Konsument, -en / die Konsumentin, -nen & le consommateur & \all{Die Konsumenten sollen kritisch bleiben.} \\
\hline
\end{tabularx}
\end{center}
\end{definition}

\begin{definition}[Verbes et expressions utiles]
\begin{itemize}
\item \all{werben (wirbt, warb, hat geworben)} : faire de la publicité ; \all{für ein Produkt werben} = « faire de la publicité pour un produit ».
\item \all{Werbung machen} : faire de la publicité. \all{Die Firma macht Werbung im Fernsehen.} « L'entreprise fait de la publicité à la télévision. »
\item \all{anbieten (bietet an, bot an, hat angeboten)} : proposer, offrir. \all{Wir bieten ein neues Produkt an.}
\item \all{überzeugen (+ Akk.)} : convaincre. \all{Die Werbung überzeugt die Kunden.}
\item \all{vergleichen (vergleicht, verglich, hat verglichen)} : comparer. \all{Vergleichen Sie die Angebote!}
\item \all{senken / erhöhen} : baisser / augmenter (les prix). \all{Die Firma senkt die Preise.}
\item \all{gehören (+ Dat.)} : appartenir à. \all{Er möchte zu einer Gruppe gehören.}
\item \all{kennenlernen} : découvrir, faire connaissance avec. \all{Die Jugendlichen lernen die Marke kennen.}
\end{itemize}
\end{definition}

\begin{attention}[Familles de mots et faux amis]
\begin{itemize}
\item Famille de \all{der Verkauf} : \all{verkaufen} (vendre), \all{der Verkäufer} (le vendeur), \all{der Ausverkauf} (les soldes). Attention : \all{verkaufen} = « vendre », mais « acheter » = \all{kaufen} (le préfixe \all{ver-} inverse le sens).
\item \all{das Angebot} = « l'offre » commerciale, jamais « l'offrande ».
\item \all{die Marke} = « la marque » commerciale ; une « marque » au sens de trace = \all{das Zeichen}.
\item \all{die Reklame} existe aussi, mais \all{die Werbung} est le terme général et moderne ; au Bac, utilisez \all{Werbung machen}.
\end{itemize}
\end{attention}

\courssub{Fiche méthode 1 : comprendre un texte sur le marketing et la publicité}

\begin{methode}[Lire et exploiter un texte publicitaire ou journalistique]
\begin{enumerate}
\item \textbf{Première lecture globale} : repérer le thème, la source probable (article de presse, brochure, site web) et la thèse de l'auteur. Souligner les mots du champ lexical : \all{die Werbung}, \all{das Produkt}, \all{der Kunde}, \all{die Marke}, \all{der Verkauf}.
\item \textbf{Deuxième lecture détaillée} : entourer les connecteurs (\all{weil}, \all{damit}, \all{um ... zu}, \all{aber}, \all{deshalb}) : ils portent l'argumentation.
\item \textbf{Repérer les impératifs} : dans un texte publicitaire, l'impératif (\all{Kaufen Sie jetzt!}, \all{Probieren Sie es!}) signale l'appel à l'action. C'est un indice de la fonction du texte : convaincre et faire acheter.
\item \textbf{Répondre aux questions} : reformuler avec ses propres mots, citer le texte entre guillemets allemands „...“, répondre par des phrases complètes avec le verbe conjugué en deuxième position.
\item \textbf{Résumé} : 3 à 5 phrases au Präsens, sans citation longue, en gardant l'ordre logique : problème, arguments, conclusion.
\end{enumerate}
\end{methode}

\begin{textebox}[Werbung überall (Auszug)]
„Werbung begleitet uns überall: im Fernsehen, im Internet, auf Plakaten in der Stadt. Die Firmen wollen die Kunden überzeugen, ihre Produkte zu kaufen. Um die Aufmerksamkeit der Konsumenten zu gewinnen, benutzen sie kurze Sätze, starke Bilder und Imperative wie „Kauf jetzt!“ oder „Probier es aus!“. Viele Jugendliche kaufen eine Marke, damit sie zu einer Gruppe gehören. Deshalb ist es wichtig, dass die Konsumenten kritisch bleiben und die Angebote vergleichen.“
\end{textebox}

« La publicité nous accompagne partout : à la télévision, sur Internet, sur des affiches en ville. Les entreprises veulent convaincre les clients d'acheter leurs produits. Pour attirer l'attention des consommateurs, elles utilisent des phrases courtes, des images fortes et des impératifs comme „Kauf jetzt!“ ou „Probier es aus!“. Beaucoup de jeunes achètent une marque pour appartenir à un groupe. C'est pourquoi il est important que les consommateurs restent critiques et comparent les offres. »

\begin{exoresolu}[Compréhension de texte : « Werbung überall »]
\textbf{Énoncé.} Lisez l'extrait ci-dessus puis répondez en allemand par des phrases complètes.

\begin{enumerate}
\item Wo findet man Werbung? Nennen Sie drei Beispiele.
\item Warum benutzen die Firmen Imperative in der Werbung?
\item Was sollen die Konsumenten nach dem Text tun?
\end{enumerate}

\tcblower
\textbf{Solution détaillée.}

\textbf{Question 1.} Le texte donne la liste dès la première phrase : \all{im Fernsehen, im Internet, auf Plakaten}. On reformule en phrase complète, verbe conjugué en position 2 :

\all{Man findet Werbung im Fernsehen, im Internet und auf Plakaten in der Stadt.}

« On trouve de la publicité à la télévision, sur Internet et sur des affiches en ville. »

\textbf{Question 2.} Le signal est \all{um ... zu gewinnen} (finalité). On reprend cette structure dans la réponse :

\all{Die Firmen benutzen Imperative, um die Aufmerksamkeit der Konsumenten zu gewinnen und sie zu überzeugen, ihre Produkte zu kaufen.}

« Les entreprises utilisent des impératifs pour attirer l'attention des consommateurs et les convaincre d'acheter leurs produits. »

Justification : \all{um ... zu} + infinitif exprime le but ; l'infinitif précédé de \all{zu} se place en fin de proposition.

\textbf{Question 3.} Le dernier membre de phrase donne la réponse : \all{kritisch bleiben} et \all{die Angebote vergleichen}.

\all{Nach dem Text sollen die Konsumenten kritisch bleiben und die Angebote vergleichen.}

« D'après le texte, les consommateurs doivent rester critiques et comparer les offres. »

\textbf{Conclusion méthodologique} : chaque réponse est une phrase complète, commence par le sujet ou un complément suivi du verbe conjugué, et reprend le vocabulaire du texte sans le recopier mot à mot.
\end{exoresolu}

\courssub{Fiche méthode 2 : employer l'impératif (du, ihr, Sie)}

\begin{propriete}[Formation de l'impératif allemand]
\begin{center}
\begin{tabular}{|l|l|l|l|}
\hline
\rowcolor{popblueL} \textbf{Personne} & \textbf{Formation} & \textbf{Exemple (kommen)} & \textbf{Exemple (kaufen)} \\
\hline
\all{du} & radical (2\textsuperscript{e} pers. sing. sans \all{-st}) & \all{Komm!} & \all{Kauf(e)!} \\
\hline
\all{ihr} & forme du Präsens sans pronom & \all{Kommt!} & \all{Kauft!} \\
\hline
\all{Sie} & infinitif + \all{Sie} & \all{Kommen Sie!} & \all{Kaufen Sie!} \\
\hline
\end{tabular}
\end{center}
\begin{itemize}
\item \textbf{Changement de voyelle e $\rightarrow$ i/ie} : il se conserve à l'impératif \all{du} : \all{gib!} (geben), \all{sprich!} (sprechen), \all{lies!} (lesen), \all{nimm!} (nehmen).
\item \textbf{Changement a $\rightarrow$ ä} : il ne se conserve \emph{pas} : \all{fahr!} (fahren), \all{schlaf!} (schlafen), \all{lauf!} (laufen).
\item \textbf{Verbes séparables} : la particule va en fin de phrase : \all{Ruf mich an!}, \all{Hören Sie gut zu!}
\item \textbf{Cas particulier} : \all{sein} : \all{sei!}, \all{seid!}, \all{seien Sie!}
\end{itemize}
\end{propriete}

\begin{methode}[Former et utiliser l'impératif allemand pas à pas]
\begin{enumerate}
\item \textbf{Identifier le destinataire} : une personne qu'on tutoie (\all{du}), plusieurs personnes qu'on tutoie (\all{ihr}), ou une personne de politesse (\all{Sie}).
\item \textbf{Former l'impératif} selon le tableau ci-dessus.
\item \textbf{Vérifier le changement de voyelle} : e $\rightarrow$ i/ie se conserve (\all{gib!}, \all{lies!}) ; a $\rightarrow$ ä disparaît (\all{fahr!}, et non \all{fähr!}).
\item \textbf{Placer la particule séparable} en fin de phrase : \all{Rufen Sie uns an!}
\item \textbf{Emploi au Bac} : dans la publicité, l'impératif \all{Sie} domine (politesse commerciale) ; le \all{du} vise les jeunes (\all{Kauf jetzt!}).
\end{enumerate}
\end{methode}

\begin{exoresolu}[Transformation : mettre à l'impératif]
\textbf{Énoncé.} Une agence de publicité à Douala prépare des slogans. Mettez les phrases suivantes à l'impératif demandé.

\begin{enumerate}
\item (kaufen, Sie) Sie kaufen unser neues Produkt.
\item (probieren, du) Du probierst den neuen Schokoriegel.
\item (anrufen, ihr) Ihr ruft uns heute an.
\item (sein, Sie) Sie sind kritisch gegenüber der Werbung.
\item (lesen, du) Du liest das Angebot genau.
\end{enumerate}

\tcblower
\textbf{Solution détaillée.}

\textbf{1.} Verbe faible, forme \all{Sie} = infinitif + \all{Sie} :

\all{Kaufen Sie unser neues Produkt!} « Achetez notre nouveau produit ! »

\textbf{2.} Forme \all{du} = radical sans \all{-st} ; \all{probieren} est un verbe faible sans changement de voyelle (le \all{-e} final est facultatif) :

\all{Probier den neuen Schokoriegel!} (ou \all{Probiere den neuen Schokoriegel!}) « Goûte la nouvelle barre chocolatée ! »

\textbf{3.} Verbe séparable \all{anrufen} : la particule \all{an} se détache et va en fin de phrase ; forme \all{ihr} = \all{ruft} :

\all{Ruft uns heute an!} « Appelez-nous aujourd'hui ! »

\textbf{4.} Verbe irrégulier \all{sein}, forme spéciale :

\all{Seien Sie kritisch gegenüber der Werbung!} « Soyez critique envers la publicité ! »

\textbf{5.} Verbe fort \all{lesen} (e $\rightarrow$ ie) : le changement de voyelle se conserve à l'impératif \all{du} :

\all{Lies das Angebot genau!} « Lis attentivement l'offre ! »

\textbf{Conclusion} : trois réflexes à vérifier à chaque impératif : la personne visée, le changement de voyelle éventuel, et la particule séparable renvoyée en fin de phrase.
\end{exoresolu}

\courssub{Fiche méthode 3 : exprimer la finalité avec um ... zu et damit}

\begin{propriete}[L'infinitif avec zu et les deux structures de finalité]
\begin{itemize}
\item \textbf{L'infinitif avec \all{zu}} : après certains verbes et expressions (\all{hoffen}, \all{versuchen}, \all{es ist wichtig}, \all{überzeugen ... zu}), l'infinitif est précédé de \all{zu} et placé en fin : \all{Die Firmen wollen die Kunden überzeugen, ihre Produkte zu kaufen.} Avec un verbe séparable, \all{zu} s'insère entre la particule et le radical : \all{kennenzulernen}, \all{anzurufen}.
\item \textbf{\all{um ... zu} + infinitif} : exprime le but quand le sujet des deux actions est \emph{identique} : \all{Er lernt Deutsch, um in Deutschland zu studieren.} (« Il apprend l'allemand pour étudier en Allemagne » — c'est lui qui étudie.)
\item \textbf{\all{damit} + subordonnée} : exprime le but quand les sujets sont \emph{différents} ; le verbe conjugué va en toute fin : \all{Die Firma senkt die Preise, damit die Kunden mehr kaufen.} (« L'entreprise baisse les prix pour que les clients achètent davantage ».)
\end{itemize}
\end{propriete}

\begin{methode}[Choisir entre um ... zu et damit]
\begin{enumerate}
\item \textbf{Poser la question du sujet} : les deux sujets (celui de la principale et celui de l'action-but) sont-ils identiques ?
\item \textbf{Sujet identique} $\Rightarrow$ \all{um ... zu} + infinitif en fin de phrase.
\item \textbf{Sujets différents} $\Rightarrow$ \all{damit} + subordonnée, verbe conjugué en fin.
\item \textbf{Interdit} : jamais \all{um ... zu} avec un verbe conjugué ; jamais \all{für} + infinitif (calque du français « pour »).
\item \textbf{Verbe séparable} : \all{zu} s'insère entre particule et radical : \all{um das Produkt kennenzulernen}.
\end{enumerate}
\end{methode}

\begin{exoresolu}[Thème : traduction de phrases à finalité]
\textbf{Énoncé.} Traduisez en allemand.

\begin{enumerate}
\item L'entreprise fait de la publicité pour vendre plus de produits.
\item Elle baisse les prix pour que les clients achètent davantage.
\item Les jeunes achètent cette marque pour appartenir à un groupe.
\item Le vendeur explique le produit pour que le client comprenne l'offre.
\end{enumerate}

\tcblower
\textbf{Solution détaillée.}

\textbf{1.} Même sujet (l'entreprise fait de la publicité et vend) $\Rightarrow$ \all{um ... zu} :

\all{Die Firma macht Werbung, um mehr Produkte zu verkaufen.}

\textbf{2.} Sujets différents (elle / les clients) $\Rightarrow$ \all{damit}, verbe conjugué \all{kaufen} en fin de subordonnée :

\all{Sie senkt die Preise, damit die Kunden mehr kaufen.}

\textbf{3.} Même sujet (les jeunes) $\Rightarrow$ \all{um ... zu} ; « appartenir à » = \all{gehören} + datif :

\all{Die Jugendlichen kaufen diese Marke, um zu einer Gruppe zu gehören.}

\textbf{4.} Sujets différents (le vendeur / le client) $\Rightarrow$ \all{damit} ; « comprendre » = \all{verstehen}, conjugué \all{versteht} en fin :

\all{Der Verkäufer erklärt das Produkt, damit der Kunde das Angebot versteht.}

\textbf{Conclusion} : le choix \all{um ... zu} / \all{damit} repose uniquement sur l'identité ou la différence des sujets. C'est une question classique du Bac : vérifier le sujet avant d'écrire la première lettre.
\end{exoresolu}

\courssub{Fiche méthode 4 : version — traduire un passage publicitaire en français}

\begin{methode}[Traduire de l'allemand vers le français]
\begin{enumerate}
\item \textbf{Lire tout le passage} avant de traduire la première phrase : repérer le temps dominant (souvent le Präsens dans la publicité).
\item \textbf{Découper chaque phrase} : sujet, verbe conjugué (position 2), compléments, infinitif ou participe en fin.
\item \textbf{Traduire les structures de finalité} : \all{um ... zu} = « pour » + infinitif ; \all{damit} = « pour que » + subjonctif (ou « afin que »).
\item \textbf{Rendre les impératifs} par des impératifs français : \all{Kaufen Sie jetzt!} = « Achetez maintenant ! »
\item \textbf{Éviter les calques} : \all{Werbung machen} = « faire de la publicité » ; \all{das Angebot} = « l'offre » (et non « l'offrande »).
\item \textbf{Relire} le français : il doit être idiomatique, pas mot à mot.
\end{enumerate}
\end{methode}

\begin{textebox}[Ein neues Angebot]
„Unsere Firma bietet ab heute ein neues Produkt an. Um die Kunden zu überzeugen, haben wir die Preise gesenkt. Besuchen Sie unseren Laden in der Stadtmitte und vergleichen Sie unsere Angebote! Wir machen Werbung im Radio, damit auch die Jugendlichen unsere Marke kennenlernen.“
\end{textebox}

\begin{exoresolu}[Version : un texte publicitaire]
\textbf{Énoncé.} Traduisez en français le texte « Ein neues Angebot » ci-dessus.

\tcblower
\textbf{Solution détaillée.}

\textbf{Phrase 1.} \all{Unsere Firma bietet ab heute ein neues Produkt an.} — Verbe séparable \all{anbieten} (« proposer ») : la particule \all{an} est en fin. \all{ab heute} = « à partir d'aujourd'hui ».

« Notre entreprise propose à partir d'aujourd'hui un nouveau produit. »

\textbf{Phrase 2.} \all{Um die Kunden zu überzeugen, haben wir die Preise gesenkt.} — \all{um ... zu} = « pour » ; Perfekt \all{haben ... gesenkt} = passé composé « avons baissé » ; quand le tour \all{um ... zu} ouvre la phrase, il occupe la position 1 et le verbe conjugué de la principale suit immédiatement (\all{haben wir}).

« Pour convaincre les clients, nous avons baissé les prix. »

\textbf{Phrase 3.} Deux impératifs \all{Sie} : \all{Besuchen Sie} et \all{vergleichen Sie} ; \all{in der Stadtmitte} = « au centre-ville ».

« Visitez notre magasin au centre-ville et comparez nos offres ! »

\textbf{Phrase 4.} \all{damit} = « pour que » + subjonctif ; sujets différents (nous / les jeunes) ; \all{kennenlernen} = « découvrir, faire connaissance avec ».

« Nous faisons de la publicité à la radio pour que les jeunes aussi découvrent notre marque. »

\textbf{Conclusion} : la version réussie repère les verbes séparables, rend correctement les deux structures de finalité et produit un français naturel.
\end{exoresolu}

\courssub{Fiche méthode 5 : production écrite guidée — un court texte publicitaire}

\begin{methode}[Rédiger un mini-texte publicitaire (5 à 8 phrases)]
\begin{enumerate}
\item \textbf{Plan en 3 parties} : présenter le produit (\all{Wir bieten ... an}), vanter ses avantages avec une finalité (\all{um ... zu} / \all{damit}), terminer par un appel à l'action à l'impératif (\all{Besuchen Sie ...!}, \all{Kaufen Sie jetzt!}).
\item \textbf{Employer le vocabulaire du thème} : \all{das Produkt}, \all{die Marke}, \all{das Angebot}, \all{der Preis}, \all{der Kunde}, \all{der Verkauf}, \all{überzeugen}, \all{vergleichen}.
\item \textbf{Grammaire obligatoire} : au moins un impératif, une structure \all{um ... zu} ou \all{damit}, un verbe séparable.
\item \textbf{Vérifications finales} : majuscule à tous les noms, verbe conjugué en position 2, infinitif en fin avec \all{zu}, Umlaute corrects (\all{ä, ö, ü}).
\end{enumerate}
\end{methode}

\begin{exoresolu}[Rédaction : un slogan développé pour un produit camerounais]
\textbf{Énoncé.} Une entreprise de Yaoundé lance un nouveau jus de fruits naturel, « Tropik ». Rédigez en allemand un court texte publicitaire (5 à 7 phrases) en utilisant : au moins deux impératifs, une structure \all{um ... zu}, une structure \all{damit} et le vocabulaire du thème.

\tcblower
\textbf{Solution détaillée — production modèle.}

\all{„Probieren Sie Tropik, unseren neuen Fruchtsaft aus Kamerun! Wir bieten ein natürliches Produkt ohne Zucker an. Trinken Sie Tropik, um gesund zu bleiben! Wir machen Werbung im Fernsehen, damit alle Kunden unsere Marke kennenlernen. Vergleichen Sie unser Angebot mit anderen Marken! Besuchen Sie unseren Laden in Yaoundé und kaufen Sie Tropik heute!“}

« Goûtez Tropik, notre nouveau jus de fruits du Cameroun ! Nous proposons un produit naturel sans sucre. Buvez Tropik pour rester en bonne santé ! Nous faisons de la publicité à la télévision pour que tous les clients découvrent notre marque. Comparez notre offre avec d'autres marques ! Visitez notre magasin à Yaoundé et achetez Tropik aujourd'hui ! »

\textbf{Justifications grammaticales.}
\begin{itemize}
\item Impératifs \all{Sie} : \all{Probieren Sie}, \all{Trinken Sie}, \all{Vergleichen Sie}, \all{Besuchen Sie}, \all{kaufen Sie} — infinitif + \all{Sie}.
\item Verbe séparable : \all{wir bieten ... an} (particule \all{an} en fin de phrase) ; \all{kennenlernen} reste soudé dans la subordonnée \all{damit}, verbe conjugué en fin.
\item Finalité même sujet : \all{um gesund zu bleiben} ; sujets différents : \all{damit alle Kunden unsere Marke kennenlernen}.
\item Vocabulaire du thème : \all{das Produkt}, \all{die Marke}, \all{das Angebot}, \all{die Kunden}, \all{Werbung machen}.
\end{itemize}

\textbf{Conclusion} : le texte respecte le plan en trois parties, intègre les structures grammaticales exigées et reste naturel et correct.
\end{exoresolu}

\begin{attention}[Les 5 erreurs qui coûtent des points au Bac]
\begin{enumerate}
\item \textbf{Oublier la majuscule aux noms} : \all{das produkt}, \all{die werbung} sont fautifs. Tout nom allemand prend une majuscule : \all{das Produkt}, \all{die Werbung}, \all{der Kunde}.
\item \textbf{Calquer « pour » par \all{für} + infinitif} : on ne dit jamais \all{für zu kaufen}. La finalité se rend par \all{um ... zu} (même sujet) ou \all{damit} (sujets différents).
\item \textbf{Mal placer le verbe après \all{damit}} : dans la subordonnée, le verbe conjugué va en toute fin : \all{damit die Kunden mehr kaufen} (et non \all{damit die Kunden kaufen mehr}).
\item \textbf{Conserver l'Umlaut à l'impératif \all{du}} : on écrit \all{Fahr!}, \all{Schlaf!} (sans Umlaut), mais \all{Gib!}, \all{Lies!} (le changement e $\rightarrow$ i/ie se conserve).
\item \textbf{Faux amis et calques lexicaux} : \all{das Angebot} = « l'offre » (pas « l'offrande ») ; \all{die Marke} = « la marque » commerciale ; « faire de la publicité » = \all{Werbung machen}. Et n'oubliez pas la particule séparable en fin de phrase : \all{Rufen Sie uns an!}
\end{enumerate}
\end{attention}

\begin{exercice}[Pour s'entraîner]
\begin{enumerate}
\item Mettez à l'impératif : a) (mitkommen, du) \all{Du kommst mit uns in den Laden mit.} b) (vergleichen, Sie) \all{Sie vergleichen die Preise.} c) (nehmen, du) \all{Du nimmst das günstige Angebot.}
\item Traduisez en allemand : a) « Les consommateurs comparent les offres pour payer moins cher. » b) « La vendeuse parle lentement pour que le client comprenne. »
\item Rédigez en allemand un slogan de 3 phrases pour un café camerounais, avec un impératif et une structure \all{um ... zu}.
\end{enumerate}
\end{exercice}

\begin{corrige}[Exercice « Pour s'entraîner »]
\textbf{1.} a) Verbe séparable \all{mitkommen}, forme \all{du} : \all{Komm mit uns in den Laden mit!} b) \all{Vergleichen Sie die Preise!} c) Verbe fort \all{nehmen} (e $\rightarrow$ i) : \all{Nimm das günstige Angebot!}

\textbf{2.} a) Même sujet $\Rightarrow$ \all{um ... zu} : \all{Die Konsumenten vergleichen die Angebote, um weniger zu zahlen.} b) Sujets différents $\Rightarrow$ \all{damit} : \all{Die Verkäuferin spricht langsam, damit der Kunde versteht.}

\textbf{3.} Modèle : \all{„Probieren Sie unseren Kaffee aus Kamerun! Trinken Sie ihn jeden Morgen, um gut in den Tag zu starten! Besuchen Sie unser Geschäft in Douala!“} (« Goûtez notre café du Cameroun ! Buvez-le chaque matin pour bien commencer la journée ! Visitez notre boutique à Douala ! »)
\end{corrige}

\newpage

\coursec{Exercices}

\courssub{Je vérifie mes connaissances}

\begin{exercice}[QCM : vocabulaire du marketing]
Entoure la bonne réponse.
\begin{enumerate}
\item \all{Die Werbung} signifie : a) la vente \quad b) la publicité \quad c) le client
\item \all{Das Angebot} signifie : a) l'offre \quad b) la marque \quad c) le prix
\item Quel est le pluriel de \all{der Kunde} ? a) die Kunden \quad b) die Kundes \quad c) die Künde
\item \all{Die Zielgruppe} désigne : a) le slogan \quad b) le groupe cible \quad c) le magasin
\end{enumerate}
\end{exercice}

\begin{exercice}[Vrai ou faux ? Justifie ta réponse]
Réponds par \all{richtig} ou \all{falsch} et justifie en français.
\begin{enumerate}
\item \all{Kauf dieses Produkt!} est un impératif à la deuxième personne du singulier (\all{du}).
\item On utilise \all{um ... zu} quand les deux propositions ont des sujets différents.
\item \all{Besuchen Sie unser Geschäft!} est un impératif de politesse (\all{Sie}).
\item Dans \all{Ich spare Geld, um einen Computer zu kaufen}, l'infinitif se place au début de la phrase.
\end{enumerate}
\end{exercice}

\begin{exercice}[Vocabulaire : articles et pluriels]
Complète le tableau avec l'article et le pluriel de chaque nom.
\begin{center}
\begin{tabular}{|l|l|l|}
\hline
\rowcolor{popblueL} Nom & Article & Pluriel \\
\hline
Werbung & \ldots & \ldots \\
\hline
Produkt & \ldots & \ldots \\
\hline
Marke & \ldots & \ldots \\
\hline
Verkauf & \ldots & \ldots \\
\hline
Angebot & \ldots & \ldots \\
\hline
Kunde & \ldots & \ldots \\
\hline
\end{tabular}
\end{center}
\end{exercice}

\begin{exercice}[Conjugaison : l'impératif]
Mets les verbes entre parenthèses à l'impératif à la forme demandée.
\begin{enumerate}
\item (kaufen, Sie) \all{\ldots\ Sie unser Produkt!}
\item (probieren, du) \all{\ldots\ es mal!}
\item (vergleichen, ihr) \all{\ldots\ die Preise!}
\item (besuchen, Sie) \all{\ldots\ Sie unser Geschäft in Douala!}
\item (sparen, du) \all{\ldots\ Geld für ein neues Handy!}
\item (sein, ihr) \all{\ldots\ vorsichtig mit der Werbung!}
\end{enumerate}
\end{exercice}

\courssub{Je m'entraîne}

\begin{exercice}[Thème : français $\rightarrow$ allemand]
Traduis les phrases suivantes en allemand.
\begin{enumerate}
\item Achète ce produit, il est en promotion !
\item Les entreprises dépensent beaucoup d'argent pour convaincre les clients.
\item J'économise pour acheter un ordinateur.
\item Le professeur explique la règle pour que les élèves comprennent.
\item Visitez notre magasin à Douala !
\end{enumerate}
\end{exercice}

\begin{exercice}[Transformation : la finalité]
Relie les deux phrases avec \all{um ... zu} ou \all{damit}. Attention aux sujets !
\begin{enumerate}
\item Die Firma senkt die Preise. Sie will mehr Kunden gewinnen.
\item Der Vater gibt dem Sohn Geld. Der Sohn soll ein Handy kaufen.
\item Wir machen Werbung im Internet. Wir wollen junge Kunden erreichen.
\item Die Lehrerin erklärt die Regel langsam. Die Schüler sollen sie verstehen.
\item Ich lerne Deutsch. Ich möchte in Deutschland studieren.
\end{enumerate}
\end{exercice}

\begin{exercice}[Texte à trous : une campagne publicitaire]
Complète le texte avec les mots suivants : \all{die Marke}, \all{das Angebot}, \all{der Kunde}, \all{der Verkauf}, \all{die Zielgruppe}, \all{der Slogan}.

\begin{textebox}[Eine Werbekampagne]
Eine große Firma in Yaoundé startet eine neue Werbekampagne. Zuerst definiert sie die \ldots : junge Leute zwischen 18 und 25 Jahren. Dann sucht sie einen guten \ldots : „Trink frisch, leb frisch!“ Die \ldots der Firma ist jetzt in ganz Kamerun bekannt. Jedes \ldots im Supermarkt ist sehr attraktiv : zwei Flaschen für den Preis von einer! Der \ldots findet die Produkte billig und gut. Am Ende steigt der \ldots um zwanzig Prozent.
\end{textebox}
\end{exercice}

\begin{exercice}[Appariements : mots et définitions]
Relie chaque mot allemand à sa définition française.
\begin{center}
\begin{tabular}{|l|l|}
\hline
\rowcolor{popblueL} Mot allemand & Définition française \\
\hline
1. die Werbung & a) personne qui achète un produit \\
\hline
2. der Kunde & b) action de promouvoir un produit dans les médias \\
\hline
3. das Produkt & c) phrase courte et frappante d'une publicité \\
\hline
4. der Slogan & d) ce qu'une entreprise fabrique et vend \\
\hline
5. die Marke & e) nom commercial d'un produit (ex. Nike) \\
\hline
\end{tabular}
\end{center}
\end{exercice}

\courssub{Je me prépare au Bac}

\begin{exercice}[Type Bac : compréhension écrite]
Lis le texte suivant, puis réponds aux questions.

\begin{textebox}[Werbung im Internet : Fluch oder Segen ?]
In Deutschland sehen die Menschen heute mehr Werbung im Internet als im Fernsehen. Fast jede Webseite zeigt bunte Anzeigen : Schuhe, Handys, Reisen, Essen. Die Firmen wollen die Kunden überall erreichen, besonders die jungen Leute. Sie analysieren die Daten der Nutzer, um personalisierte Werbung zu schicken. Wer zum Beispiel nach Schuhen sucht, sieht danach überall Schuhwerbung.

Viele Kunden finden das praktisch : „Ich bekomme Angebote, die mich wirklich interessieren“, sagt Anna, 22 Jahre alt, aus Hamburg. Aber andere sind skeptisch. „Die Firmen wissen zu viel über uns“, meint Herr Müller, ein Informatiklehrer aus München. Er benutzt ein spezielles Programm, damit die Werbung verschwindet.

Die Werbeindustrie verdient mit Online-Werbung jedes Jahr Milliarden. Auch kleine Firmen können billig werben : ein Geschäft in einem Dorf kann so Kunden in der ganzen Welt finden. Experten raten den Konsumenten : „Vergleichen Sie die Preise! Lesen Sie die Bewertungen! Kaufen Sie nicht zu schnell!“ Denn nicht jedes Angebot im Internet ist ehrlich. Die Verbraucherzentrale warnt : Manche Webseiten verkaufen gefälschte Produkte. Deshalb muss der moderne Kunde kritisch denken, damit er keine schlechten Erfahrungen macht.
\end{textebox}

\textbf{A. Questions de compréhension} (réponds en allemand, avec des phrases complètes)
\begin{enumerate}
\item Wo sehen die Deutschen heute mehr Werbung als früher ?
\item Warum analysieren die Firmen die Daten der Nutzer ?
\item Was findet Anna praktisch ?
\item Was macht Herr Müller gegen die Werbung ?
\item Welchen Rat geben die Experten den Konsumenten ? Nenne zwei Ratschläge.
\end{enumerate}

\textbf{B. Vrai ou faux ?} Réponds par \all{richtig} ou \all{falsch} et justifie avec une phrase du texte.
\begin{enumerate}
\item Nur große Firmen können im Internet werben.
\item Alle Webseiten im Internet sind ehrlich.
\item Die Verbraucherzentrale warnt vor gefälschten Produkten.
\end{enumerate}

\textbf{C. Questions de langue}
\begin{enumerate}
\item Relève dans le texte deux impératifs et indique à quelle forme (\all{du}, \all{ihr} ou \all{Sie}) ils sont conjugués.
\item Relève une construction avec \all{um ... zu} et une avec \all{damit}, puis explique la différence d'emploi.
\end{enumerate}
\end{exercice}

\begin{exercice}[Type Bac : thème et version]
\textbf{A. Thème} — Traduis en allemand (attention à l'impératif et à la finalité) :
\begin{enumerate}
\item Achetez nos produits, ils sont de bonne qualité !
\item La marque fait de la publicité pour attirer de nouveaux clients.
\item Mon frère travaille beaucoup afin que sa famille vive bien.
\end{enumerate}

\textbf{B. Version} — Traduis en français le texte suivant :

\begin{textebox}[Un texte publicitaire]
Liebe Kunden! Besuchen Sie unser neues Geschäft in der Stadtmitte! Wir bieten Ihnen die besten Preise der Region. Kommen Sie heute, um von unseren Sonderangeboten zu profitieren! Wir haben unsere Preise gesenkt, damit jeder Kunde etwas Schönes kaufen kann. Probieren Sie unsere Produkte — Sie werden zufrieden sein!
\end{textebox}
\end{exercice}

\begin{exercice}[Type Bac : production écrite]
\begin{textebox}[Sujet de rédaction]
\all{Schreiben Sie einen kurzen Werbetext (60–80 Wörter) für ein Produkt Ihrer Wahl (z. B. ein Getränk, ein Handy, eine Schule, ein Restaurant in Kamerun). Benutzen Sie den Imperativ (mindestens zwei Formen) und mindestens eine Finalkonstruktion (um ... zu oder damit). Vergessen Sie nicht einen Slogan!}
\end{textebox}

Consignes méthodologiques :
\begin{enumerate}
\item Choisis un produit et un groupe cible (\all{Zielgruppe}).
\item Rédige un slogan court et frappant.
\item Emploie au moins deux impératifs (\all{du}, \all{ihr} ou \all{Sie}).
\item Intègre au moins une construction finale : \all{um ... zu} (même sujet) ou \all{damit} (sujets différents).
\item Vérifie : majuscules des noms, place du verbe, infinitif avec \all{zu} en fin de proposition.
\end{enumerate}
\end{exercice}



\newpage

\coursec{Corrigés des exercices}

\courssub{Je vérifie mes connaissances}

\begin{corrige}[Exercice 1 — QCM : vocabulaire du marketing]
\begin{enumerate}
\item b) la publicité. \all{Die Werbung} (pluriel : \all{die Werbungen}) = la publicité ; la vente se dit \all{der Verkauf}, le client \all{der Kunde}.
\item a) l'offre. \all{Das Angebot} (pluriel : \all{die Angebote}) = l'offre ; la marque se dit \all{die Marke}, le prix \all{der Preis}.
\item a) \all{die Kunden}. \all{Der Kunde} est un nom faible (déclinaison en \all{-n}) : \all{der Kunde, die Kunden}.
\item b) le groupe cible. \all{Die Zielgruppe} (pluriel : \all{die Zielgruppen}) : \all{das Ziel} = le but, \all{die Gruppe} = le groupe.
\end{enumerate}
\textbf{Point de vigilance :} ne pas confondre \all{das Angebot} (l'offre) et \all{der Verkauf} (la vente) ; attention au faux ami : \all{die Promotion} n'existe pas au sens de « promotion commerciale ».
\end{corrige}

\begin{corrige}[Exercice 2 — Vrai ou faux ? Justifie ta réponse]
\begin{enumerate}
\item \all{Richtig}. À l'impératif \all{du}, on prend le radical de la 2\up{e} personne du singulier sans \all{-st} : \all{du kaufst} $\rightarrow$ \all{Kauf(e)!}.
\item \all{Falsch}. C'est le contraire : on utilise \all{um ... zu} quand les deux propositions ont le \emph{même} sujet ; quand les sujets sont différents, on utilise \all{damit} + subordonnée (verbe conjugué en fin).
\item \all{Richtig}. À la forme \all{Sie}, l'impératif reprend la forme du présent avec le pronom \all{Sie} placé après le verbe : \all{Besuchen Sie!}.
\item \all{Falsch}. Dans une construction \all{um ... zu}, l'infinitif précédé de \all{zu} se place à la \emph{fin} de la proposition : \all{..., um einen Computer zu kaufen}.
\end{enumerate}
\textbf{Point de vigilance :} la règle d'or du chapitre : même sujet $\rightarrow$ \all{um ... zu} ; sujets différents $\rightarrow$ \all{damit}.
\end{corrige}

\begin{corrige}[Exercice 3 — Vocabulaire : articles et pluriels]
\begin{center}
\begin{tabular}{|l|l|l|}
\hline
\rowcolor{popblueL} Nom & Article & Pluriel \\
\hline
Werbung & die & die Werbungen \\
\hline
Produkt & das & die Produkte \\
\hline
Marke & die & die Marken \\
\hline
Verkauf & der & die Verkäufe \\
\hline
Angebot & das & die Angebote \\
\hline
Kunde & der & die Kunden \\
\hline
\end{tabular}
\end{center}
\textbf{Point de vigilance :} \all{der Verkauf} prend un Umlaut au pluriel (\all{die Verkäufe}) et \all{der Kunde} est un nom faible : \all{die Kunden}. Apprends toujours chaque nom avec son article et son pluriel.
\end{corrige}

\begin{corrige}[Exercice 4 — Conjugaison : l'impératif]
\begin{enumerate}
\item \all{Kaufen Sie unser Produkt!} (forme \all{Sie} = présent + \all{Sie} après le verbe)
\item \all{Probier(e) es mal!} (forme \all{du} = radical sans \all{-st})
\item \all{Vergleicht die Preise!} (forme \all{ihr} = comme au présent : \all{ihr vergleicht})
\item \all{Besuchen Sie unser Geschäft in Douala!}
\item \all{Spar(e) Geld für ein neues Handy!}
\item \all{Seid vorsichtig mit der Werbung!} (impératif irrégulier de \all{sein} : \all{sei! seid! seien Sie!})
\end{enumerate}
\textbf{Point de vigilance :} les verbes forts à changement de voyelle \all{e} $\rightarrow$ \all{i/ie} gardent ce changement à l'impératif \all{du} : \all{Gib!}, \all{Lies!} ; mais jamais d'Umlaut : \all{Fahr!} (et non \all{Fähr!}).
\end{corrige}

\courssub{Je m'entraîne}

\begin{corrige}[Exercice 5 — Thème : français $\rightarrow$ allemand]
\begin{enumerate}
\item \all{Kauf dieses Produkt, es ist im Angebot!} (ou : \all{..., es ist reduziert!}) — « en promotion » se dit \all{im Angebot} ou \all{reduziert}, jamais \all{in Promotion}.
\item \all{Die Firmen geben viel Geld aus, um die Kunden zu überzeugen.} — même sujet (\all{die Firmen}) $\rightarrow$ \all{um ... zu} ; « dépenser » = \all{Geld ausgeben} (verbe à particule séparable : \all{geben ... aus}).
\item \all{Ich spare, um einen Computer zu kaufen.} — même sujet $\rightarrow$ \all{um ... zu}, infinitif en fin.
\item \all{Der Lehrer erklärt die Regel, damit die Schüler sie verstehen.} — sujets différents (\all{der Lehrer} / \all{die Schüler}) $\rightarrow$ \all{damit} + verbe conjugué en fin.
\item \all{Besuchen Sie unser Geschäft in Douala!} — impératif de politesse \all{Sie}.
\end{enumerate}
\textbf{Point de vigilance :} dans \all{um ... zu}, l'infinitif est toujours en fin de proposition ; avec \all{damit}, le verbe conjugué est en fin de subordonnée. Vérifie la majuscule à tous les noms : \all{das Produkt}, \all{das Geld}, \all{die Regel}, \all{das Geschäft}.
\end{corrige}

\begin{corrige}[Exercice 6 — Transformation : la finalité]
\begin{enumerate}
\item \all{Die Firma senkt die Preise, um mehr Kunden zu gewinnen.} (même sujet : \all{die Firma})
\item \all{Der Vater gibt dem Sohn Geld, damit der Sohn ein Handy kauft.} (sujets différents : \all{der Vater} / \all{der Sohn})
\item \all{Wir machen Werbung im Internet, um junge Kunden zu erreichen.} (même sujet : \all{wir})
\item \all{Die Lehrerin erklärt die Regel langsam, damit die Schüler sie verstehen.} (sujets différents)
\item \all{Ich lerne Deutsch, um in Deutschland zu studieren.} (même sujet : \all{ich})
\end{enumerate}
\textbf{Point de vigilance :} avant de choisir, souligne le sujet de chaque proposition. Avec \all{damit}, le sujet de la subordonnée doit être répété : \all{damit der Sohn ein Handy kauft}.
\end{corrige}

\begin{corrige}[Exercice 7 — Texte à trous : une campagne publicitaire]
Dans l'ordre : \all{Zielgruppe} (les jeunes de 18 à 25 ans = le groupe cible) ; \all{Slogan} (« Trink frisch, leb frisch ! » = la phrase publicitaire) ; \all{Marke} (la marque est connue dans tout le Cameroun) ; \all{Angebot} (deux bouteilles pour le prix d'une = une offre) ; \all{Kunde} (celui qui trouve les produits bons et bon marché) ; \all{Verkauf} (les ventes augmentent de vingt pour cent).

Texte complet : \all{Eine große Firma in Yaoundé startet eine neue Werbekampagne. Zuerst definiert sie die Zielgruppe : junge Leute zwischen 18 und 25 Jahren. Dann sucht sie einen guten Slogan : „Trink frisch, leb frisch!“ Die Marke der Firma ist jetzt in ganz Kamerun bekannt. Jedes Angebot im Supermarkt ist sehr attraktiv : zwei Flaschen für den Preis von einer! Der Kunde findet die Produkte billig und gut. Am Ende steigt der Verkauf um zwanzig Prozent.}

\textbf{Point de vigilance :} observe les cas : \all{einen guten Slogan} (accusatif après \all{suchen}), \all{der Verkauf steigt} (nominatif, sujet de \all{steigen}).
\end{corrige}

\begin{corrige}[Exercice 8 — Appariements : mots et définitions]
1 -- b ; 2 -- a ; 3 -- d ; 4 -- c ; 5 -- e.

\all{Die Werbung} = la publicité (promouvoir dans les médias) ; \all{der Kunde} = le client (celui qui achète) ; \all{das Produkt} = le produit (ce qu'on fabrique et vend) ; \all{der Slogan} = le slogan (phrase courte et frappante) ; \all{die Marke} = la marque (nom commercial).

\textbf{Point de vigilance :} \all{die Marke} et \all{das Produkt} ne sont pas synonymes : la marque est le nom commercial, le produit est l'objet fabriqué.
\end{corrige}

\courssub{Je me prépare au Bac}

\begin{corrige}[Exercice 9 — Type Bac : compréhension écrite]
\textbf{A. Questions de compréhension}
\begin{enumerate}
\item \all{Die Deutschen sehen heute mehr Werbung im Internet (als im Fernsehen).}
\item \all{Sie analysieren die Daten der Nutzer, um personalisierte Werbung zu schicken.}
\item \all{Anna findet es praktisch, dass sie Angebote bekommt, die sie wirklich interessieren.}
\item \all{Er benutzt ein spezielles Programm, damit die Werbung verschwindet.}
\item \all{Die Experten raten den Konsumenten : „Vergleichen Sie die Preise! Lesen Sie die Bewertungen! Kaufen Sie nicht zu schnell!“} (deux de ces trois conseils suffisent).
\end{enumerate}

\textbf{B. Vrai ou faux ?}
\begin{enumerate}
\item \all{Falsch} : \all{Auch kleine Firmen können billig werben.}
\item \all{Falsch} : \all{Nicht jedes Angebot im Internet ist ehrlich.}
\item \all{Richtig} : \all{Die Verbraucherzentrale warnt : Manche Webseiten verkaufen gefälschte Produkte.}
\end{enumerate}

\textbf{C. Questions de langue}
\begin{enumerate}
\item \all{Vergleichen Sie die Preise!}, \all{Lesen Sie die Bewertungen!}, \all{Kaufen Sie nicht zu schnell!} : ce sont des impératifs à la forme de politesse \all{Sie} (verbe à l'infinitif présent + pronom \all{Sie} après le verbe).
\item \all{um ... zu} : \all{Sie analysieren die Daten der Nutzer, um personalisierte Werbung zu schicken.} — même sujet (\all{sie} = les entreprises) dans les deux propositions, donc infinitif avec \all{zu} en fin. \all{damit} : \all{Er benutzt ein spezielles Programm, damit die Werbung verschwindet.} — sujets différents (\all{er} / \all{die Werbung}), donc subordonnée avec verbe conjugué en fin.
\end{enumerate}

\textbf{Barème indicatif (sur 20) :} A : 5 $\times$ 2 = 10 pts (phrase complète et correcte : 2 pts ; réponse juste mais fautes de langue : 1 pt) ; B : 3 $\times$ 2 = 6 pts (\all{richtig}/\all{falsch} : 1 pt ; justification citée : 1 pt) ; C : 2 $\times$ 2 = 4 pts.

\textbf{Point de vigilance :} réponds toujours par une phrase complète en allemand, en reprenant les termes du texte ; ne recopie jamais un passage entier sans l'adapter à la question.
\end{corrige}

\begin{corrige}[Exercice 10 — Type Bac : thème et version]
\textbf{A. Thème}
\begin{enumerate}
\item \all{Kaufen Sie unsere Produkte, sie sind von guter Qualität!} — « de bonne qualité » = \all{von guter Qualität} (datif après \all{von}, adjectif décliné).
\item \all{Die Marke macht Werbung, um neue Kunden zu gewinnen.} (ou : \all{..., um neue Kunden anzuziehen.}) — même sujet $\rightarrow$ \all{um ... zu} ; « attirer » = \all{anziehen} (particule séparable : \all{anzuziehen}) ou \all{gewinnen}.
\item \all{Mein Bruder arbeitet viel, damit seine Familie gut lebt.} — sujets différents (\all{mein Bruder} / \all{seine Familie}) $\rightarrow$ \all{damit}, verbe conjugué en fin.
\end{enumerate}

\textbf{B. Version}

« Chers clients ! Visitez notre nouveau magasin au centre-ville ! Nous vous offrons les meilleurs prix de la région. Venez aujourd'hui pour profiter de nos offres spéciales ! Nous avons baissé nos prix afin que chaque client puisse acheter quelque chose de beau. Essayez nos produits — vous serez satisfaits ! »

Points de langue : \all{die Stadtmitte} = le centre-ville ; \all{das Sonderangebot} = l'offre spéciale ; \all{etwas Schönes} = quelque chose de beau (adjectif nominalisé avec majuscule) ; \all{Sie werden zufrieden sein} = futur (\all{werden} + infinitif).

\textbf{Barème indicatif (sur 20) :} thème : 3 $\times$ 3 = 9 pts (structure correcte : 2 pts ; vocabulaire : 1 pt) ; version : 11 pts (fidélité du sens : 7 pts ; qualité du français : 4 pts).

\textbf{Point de vigilance :} en version, traduis les impératifs allemands par des impératifs français (\all{Besuchen Sie} = « visitez ») et rends \all{um ... zu} par « pour » et \all{damit} par « afin que ».
\end{corrige}

\begin{corrige}[Exercice 11 — Type Bac : production écrite]
\textbf{Modèle de production} (environ 70 mots) :

\all{„Mambo Juice — Trink frisch, leb frisch!“}

\all{Liebe Freunde! Probiert unseren neuen Fruchtsaft aus Kamerun! Er schmeckt frisch und natürlich, mit Mangos und Ananas aus der Region Douala. Kommt heute in unser Geschäft, um unsere Sonderangebote zu entdecken! Wir haben die Preise gesenkt, damit jeder junge Kunde unser Produkt kaufen kann. Kauft zwei Flaschen und bezahlt nur eine! Mambo Juice — die frische Marke für ganz Kamerun!}

\textbf{Analyse du modèle :}
\begin{itemize}
\item Slogan : \all{Trink frisch, leb frisch!} (deux impératifs \all{du} dans le slogan).
\item Impératifs : \all{Probiert}, \all{Kommt}, \all{Kauft}, \all{bezahlt} (forme \all{ihr}, car le groupe cible est un groupe de jeunes).
\item Finalité : \all{um unsere Sonderangebote zu entdecken} (même sujet \all{ihr}) ; \all{damit jeder junge Kunde unser Produkt kaufen kann} (sujets différents : \all{wir} / \all{jeder junge Kunde}).
\item Vocabulaire du thème : \all{das Sonderangebot}, \all{der Preis}, \all{der Kunde}, \all{das Produkt}, \all{die Marke}.
\end{itemize}

\textbf{Barème indicatif (sur 20) :} respect des consignes (slogan, 2 impératifs, 1 construction finale, 60--80 mots) : 8 pts ; correction grammaticale (impératifs, place du verbe, cas) : 6 pts ; vocabulaire thématique : 4 pts ; présentation et orthographe (majuscules, Umlaute) : 2 pts.

\textbf{Point de vigilance :} choisis une seule forme d'impératif (\all{du}, \all{ihr} ou \all{Sie}) et garde-la du début à la fin ; relis-toi en vérifiant la majuscule de chaque nom et la place de l'infinitif avec \all{zu}.
\end{corrige}

\newpage

\coursec{Fiche bilan}

\begin{popfigure}
\begin{tikzpicture}[mindmap, grow cyclic, every node/.style={concept, text=white, font=\small\bfseries}, concept color=popblue, level 1/.append style={level distance=3.6cm, sibling angle=60}, level 2/.append style={level distance=2.2cm, sibling angle=40}]
\node[concept, font=\normalsize\bfseries] {Marketing und Werbung}
  child[concept color=poporange] { node {Lexique}
    child { node[concept color=poporange!80] {die Werbung} }
    child { node[concept color=poporange!80] {das Produkt} }
    child { node[concept color=poporange!80] {der Kunde} }
  }
  child[concept color=popgreen] { node {Impératif}
    child { node[concept color=popgreen!80] {du : Kauf(e)!} }
    child { node[concept color=popgreen!80] {ihr : Kauft!} }
    child { node[concept color=popgreen!80] {Sie : Kaufen Sie!} }
  }
  child[concept color=poppurple] { node {Infinitif mit zu}
    child { node[concept color=poppurple!80] {zu + Inf.} }
    child { node[concept color=poppurple!80] {anzuprobieren} }
  }
  child[concept color=popgold] { node {Finalité}
    child { node[concept color=popgold!80] {um ... zu} }
    child { node[concept color=popgold!80] {damit} }
  }
  child[concept color=poppink] { node {Compréhension}
    child { node[concept color=poppink!80] {globale} }
    child { node[concept color=poppink!80] {détaillée} }
  }
  child[concept color=popdark] { node {Production}
    child { node[concept color=popdark!80] {résumé} }
    child { node[concept color=popdark!80] {slogan} }
  };
\end{tikzpicture}
\legende{Carte mentale de la leçon AL07 : Marketing et publicité}
\end{popfigure}

\begin{aretenir}[L'essentiel de la leçon]
\textbf{1. Lexique clé (à connaître avec article et pluriel)}
\begin{itemize}\setlength{\itemsep}{1pt}
  \item \all{die Werbung, -en} : la publicité ; \all{der Werbespot, -s} : le spot publicitaire ; \all{das Plakat, -e} : l'affiche publicitaire
  \item \all{das Produkt, -e} : le produit ; \all{die Marke, -n} : la marque ; \all{das Angebot, -e} : l'offre
  \item \all{der Kunde, -n / die Kundin, -nen} : le client / la cliente ; \all{der Verkauf, die Verkäufe} : la vente ; \all{der Preis, -e} : le prix
  \item \all{der Markt, die Märkte} : le marché ; \all{die Reklame, -n} : la réclame, la publicité ; \all{der Verbraucher, - / die Verbraucherin, -nen} : le consommateur / la consommatrice
  \item Verbes : \all{werben (wirbt, warb, hat geworben)} : faire de la publicité ; \all{anbieten (bietet an, bot an, hat angeboten)} : proposer ; \all{verkaufen} : vendre ; \all{kaufen} : acheter ; \all{überzeugen} : convaincre ; \all{beeinflussen} : influencer
  \item Expressions utiles : \all{Werbung machen für ein Produkt} : faire de la publicité pour un produit ; \all{die Kunden überzeugen} : convaincre les clients ; \all{ein neues Produkt auf den Markt bringen} : lancer un nouveau produit sur le marché
\end{itemize}

\textbf{2. Grammaire à connaître par c\oe ur}
\begin{center}
\begin{tabularx}{\linewidth}{|l|X|X|}
\hline
\rowcolor{popblueL} \textbf{Point} & \textbf{Règle} & \textbf{Exemple} \\
\hline
Impératif \all{du} & radical du présent (+\all{-e} facultatif) ; la mutation \all{e>i/ie} est conservée, mais pas \all{a>ä} & \all{Kauf(e)! Lies! Gib!} mais \all{Fahr(e)!} (et non *\all{Fähr!}) \\
\hline
Impératif \all{ihr} & comme la 2e personne du pluriel, sans le pronom & \all{Kauft! Lest! Gebt!} \\
\hline
Impératif \all{Sie} & forme de politesse inversée : verbe + \all{Sie} & \all{Kaufen Sie! Lesen Sie!} \\
\hline
\all{sein} & impératif irrégulier & \all{Sei! Seid! Seien Sie!} \\
\hline
Infinitif avec \all{zu} & après \all{versuchen, hoffen, vergessen, planen, versprechen}... ; l'infinitif en \all{zu} va à la fin ; verbe à particule : \all{zu} s'insère (\all{anzuprobieren}) & \all{Ich versuche, das Produkt zu kaufen.} \\
\hline
\all{um ... zu} & finalité, même sujet dans les deux verbes & \all{Die Firma wirbt, um neue Kunden zu gewinnen.} \\
\hline
\all{damit} & finalité, sujets différents ; verbe conjugué en fin de subordonnée & \all{Die Firma senkt den Preis, damit die Kunden mehr kaufen.} \\
\hline
\end{tabularx}
\end{center}

\textbf{3. Méthodes express}
\begin{itemize}\setlength{\itemsep}{1pt}
  \item Compréhension : lire d'abord le titre et la source, repérer le thème, puis faire une lecture détaillée en soulignant les mots-clés (\all{das Produkt, der Preis, der Kunde, die Marke}).
  \item Résumé : une phrase par idée principale, au présent, sans longue citation, avec des connecteurs (\all{zuerst, dann, außerdem, schließlich}).
  \item Slogan publicitaire : impératif + adjectif + nom de marque, court et rythmé : \all{Probier es! Kauf jetzt!} (« Essaie-le ! Achète maintenant ! »)
\end{itemize}
\end{aretenir}

\begin{attention}[Pièges fréquents des francophones]
\begin{itemize}\setlength{\itemsep}{1pt}
  \item Ne pas garder le \all{-st} de la 2e personne : \all{du kaufst} devient \all{Kauf!} (et non *\all{Kaufst!}).
  \item Ne pas confondre \all{die Werbung} (la publicité en général) et \all{die Reklame} (l'annonce, la réclame) ; \all{der Markt} (le marché) n'est pas \all{die Marke} (la marque).
  \item Après \all{damit}, le verbe conjugué va obligatoirement à la fin : \all{..., damit die Kunden mehr kaufen.} (et non *\all{damit die Kunden kaufen mehr}).
  \item \all{um ... zu} exige le même sujet ; sinon, utiliser \all{damit}.
  \item Majuscule à tous les noms : \all{das Angebot, der Verkauf, die Werbung}.
\end{itemize}
\end{attention}

\begin{aretenir}[Checklist avant le Bac]
\begin{itemize}\setlength{\itemsep}{2pt}
  \item $\square$ Je connais le vocabulaire du marketing avec article et pluriel (\all{die Werbung, das Produkt, der Kunde, das Angebot, der Verkauf, die Marke}).
  \item $\square$ Je sais former l'impératif aux trois personnes : \all{du}, \all{ihr}, \all{Sie}.
  \item $\square$ Je connais les impératifs irréguliers : \all{Sei! Seid! Seien Sie!}, \all{Gib! Nimm! Lies!}
  \item $\square$ Je sais que l'impératif \all{du} prend le radical sans \all{-st} : \all{du kaufst} devient \all{Kauf!}
  \item $\square$ Je sais que la mutation \all{a>ä} disparaît à l'impératif : \all{du fährst} devient \all{Fahr!}
  \item $\square$ Je place \all{zu} devant l'infinitif en fin de proposition : \all{..., das Produkt zu kaufen.}
  \item $\square$ Je sais insérer \all{zu} dans les verbes à particule : \all{anzuprobieren, einzukaufen}.
  \item $\square$ Je choisis \all{um ... zu} quand le sujet est identique, \all{damit} quand les sujets diffèrent.
  \item $\square$ Je mets le verbe conjugué à la fin après \all{damit}.
  \item $\square$ Je n'oublie pas la majuscule à tous les noms allemands.
  \item $\square$ Je sais rédiger un résumé court au présent avec des connecteurs.
  \item $\square$ Je sais créer un slogan publicitaire simple à l'impératif.
  \item $\square$ Je vérifie l'ordre des mots : verbe en position 2 dans la principale, en fin dans la subordonnée.
\end{itemize}
\end{aretenir}

\findecours

\end{document}
